 % !TEX root = main.tex


\section{Tight ROM Security via Circular Discrete-Logarithm} \label{sec-fs}

The Fiat–Shamir (FS) transform~\cite{C:FiaSha86} has notoriously loose security proofs. One line of work attains tight ROM proofs when the underlying $\Sigma$-protocol is lossy~\cite{JC:AFLT16,C:KilMasPan16}, but Schnorr’s is not. Another line proposes alternative transforms with tight ROM proofs~\cite{C:Fischlin05,AC:BelPoeSte16,RotemT25}, yet these are not used in practice. So, we focus on FS and particularly for Schnorr and its derivatives.

%Indeed, decades of cryptanalysis  suggests there are no ``better-than-trivial'' attacks on these schemes, \emph{e.g.} better than breaking DL in the case of  Schnorr signatures.  However, absent rigorous justification of this, we  are at risk  of  surprise attacks. 
To justify their security at practical parameters, we propose a new agenda: tight ROM proofs under novel ``circular'' strengthenings of 
classic assumptions (unrelated to circular security for encryption). 
As tight ROM proofs under classic assumptions have resisted decades of work, 
our approach offers a timely and promising path forward. 
%In parallel, we will vet new group-theoretic assumptions within idealized models, 
%focusing the idealization on the assumption rather than the scheme. 
%This shift provides sharper insight and increases trustworthiness for practice.


%In particular, our goal is to show that such schemes admit tight ROM proofs under  novel ``circular'' assumptions. %, resolving decades-old looseness 

%Indeed, decades of cryptanalysis  suggests there are no ``better-than-trivial'' attacks on these schemes, \emph{e.g.} better than breaking DL in the case of  Schnorr signatures.  However, absent rigorous justification of this, we  are at risk  of  surprise attacks. Accordingly, we propose a research agenda aimed at showing  tight ROM proofs under mild strengthenings of the assumptions. %Prior work showed a tight security proof when applying Fiat-Shamir to a  \emph{lossy} identification schemes~\cite{}. However, common identification schemes used in practice like Schnorr are not lossy. On the other hand, cryptanalysis  indicates that a tight security proof for Schnorr and related schemes would  be the ``right answer.''
%  Schnorr will thus be a representative application of our techniques, but it is by no means the only one. 

%We aim to provide better foundations for group-based cryptographic protocols, particularly those that combine groups and hash functions. While there is a push towards post-quantum cryptography, it remains vital to better analyze group-based schemes as they are still pervasively used.  

%We will first focus on applications \emph{discrete-log-based} cryptography. While there is a push towards post-quantum, DL-based cryptography remains vital to study as it is pervasively used in practice. 

\subsection{Tight Security  of Schnorr Signatures}

%By far the most popular FS-based scheme in practice is Schnorr signatures, so we start there.

% We begin with the most popular case of Schnorr signatures. 
 
 \heading{The scheme.} Schnorr signatures~\cite{C:Schnorr89} are ubiquitous, often in the form of EdDSA~\cite{eddsa}. %implemented over twisted Edwards curves, \emph{e.g.}, in TLS, Signal, and in Bitcoin since the Taproot soft-fork upgrade \cite{add:BIP340}.  
 To define the scheme, fix a cyclic group $\GG$  of prime order $p$ with generator $g \in \GG$, and let $\hashH \colon \GG \times \bits^* \to \Z_p$ be hash function. Define the Schnorr signature scheme $\mathsf{Sch}[\GG,\hashH] = (\calK,\calS,\calV)$ as follows.  The key-generation $\calK$ algorithm   chooses random $x\in \Z_p$ and returns $g^x$ as the verification key and $x$ as the signing key.  Given $x$ and a message $m\in\bits^*$, the signing algorithm $\calS$ returns $(g^r,(r+\hashH(g^r \| m) \cdot x) \bmod p)$ for random $r\in \Z_p$.  Given $X = g^x$,  message $m\in\bits^*$, and  signature $(R,z)$, the verification algorithm $\calV$ returns 1 iff  $g^z = R \cdot X^{\hashH(R \| m)}$. 

%We will initially  focus on the most popular case of Schnorr signatures~\cite{JC:Schnorr91}. 

%The popular Schnorr signature scheme~\cite{JC:Schnorr91} is obtained from applying  FS to Schnorr's $\Sigma$-protocol.  
 
%The key-generation algorithm $\calK$ chooses random $x\getsr Z_p$ and returns $g^x$ as the verification key and $x$ as the signing key.  Given $x$ and message $m\in\bits^*$, the signing algorithm $\calS$ returns $(g^r,r+\Hash(g^r \| m) \cdot x \bmod n)$.  Given $X = g^x$,  message $m\in\bits^*$, and a purported signature $(R,z)$, the verification algorithm $\calV$ returns 1 iff  $g^z = R \cdot X^{\Hash(R \| m)}$.
%Let us give intuition for why kPGG1 or kPGG2 are relevant here. To start with, we assume the forger's signing queries $m_1,\ldots,m_q$ and forgery message $m^*$ are all output upfrront. Consider a KPGG1 source $\advs$ that queries random $x$ to get $X$ and  for each $i \in [q]$ queries and  random $r_i \in \Z_p$ to get $R_i$, then outputs leakage $L = (X,(R_1,r_1 + \Hash(R_1 \| m_1)\cdot x),\ldots,(R_q,r_q + \Hash(R_q \| m_q)\cdot x))$. Given $L$, 
%We ultimately want to eliminate idealized models here, and we also show applications to tight security reductions in the RO model.  
%As far as we are aware, all positive results about it are in the  RO model~\cite{JC:PoiSte00,EC:AABN02,EC:FucPloSeu20,INDOCRYPT:BelDai20,C:RotSeg21} or the GGM~\cite{JMC:NevSmaWar09,C:CLMQ21,EC:BloLee22,EPRINT:Shoup23},  except that \emph{every} instantiation is secure wrt.~key-recovery under one-more discrete log~\cite{AC:PaiVer05}.  Even in the RO model, there is no known \emph{tight} reduction to DL, or even one that achieves the weaker goal of  avoiding the  ``forking lemma''~\cite{JC:PoiSte00}, except using an interactive assumption~\cite{INDOCRYPT:BelDai20}.


\heading{Prior work.}  Pointcheval and Stern~\cite{EC:PoiSte96} show the scheme is secure in the ROM assuming the discrete-logarithm (DL) problem is hard in the underlying group, but their reduction is quite loose, requiring 768-bit curve vs.~256-bit used in practice.
 Subsequent work \cite{EC:FucPloSeu20} showed tight security using additional idealized models such as the algebraic group model (AGM)~\cite{C:FucKilLos18} or  generic group model (GGM)~\cite{JMC:NevSmaWar09,EPRINT:Shoup23,C:CLMQ21}.
We want to avoid such idealized models.

Using an interactive strengthening of DL called multi-base DL (MBDL), Bellare and Dai (BD)~\cite{INDOCRYPT:BelDai20} show a reduction that loses only a factor of the number of hash queries.  (A subsequent result~\cite{C:RotSeg21} is  less tight  but uses a non-interactive assumption.) Here ``interactive'' means the adversary talks to an oracle. In the case of MBDL and related interactive assumptions like one-more DL (OMDL)~\cite{JC:BNPS03}, the oracle is moreover \emph{inefficient} (it computes DL), so these assumptions are ``non-falsifiable''~\cite{C:Naor03} in that an attack is hard to demonstrate.% Such assumptions defy cryptanalytic validation.
%We therefore try to stick to non-interactive and falsifiable assumptions as much as possible. 

Another sequence of works \cite{AC:PaiVer05,C:GarBhaLok08,EC:Seurin12,JC:FleJagSch19}  shows that there is \emph{no}  tight ``generic'' reduction in the ROM (even under a minimal formulation of security) for Schnorr signatures under \emph{any} ``representation-independent'' non-interactive assumption. ``Representation-independent'' means that a solution remains valid under change  of group representation. 
%Throughout, \emph{we will strive to stick to non-interactive and falsifiable assumptions whenever possible.} 
Thus, we ask:
 \begin{ques}  %Compared to simple assumptions like DL complex and interactive problems shown to hold in these models are more likely to fall prey to the reach of these results. 
Is there a \emph{tight} (non-rewinding/non-guessing) reduction in the ROM from security of Schnorr signatures to a non-interactive and falsifiable assumption?
 \end{ques}

% As most common assumptions are  representation-dependent, it seems hard to imagine \emph{a priori} what a candidate would look like.
 
%It turns out these questions are related, and the former is of indepedent interest.

%In particular, in programmable AI-GGM, Chen \emph{et al.}~\cite{C:CLMQ21} demonstrated \emph{non-cryptographic} hash functions work. We stress that their idea is \emph{not} to build a RO using the GGM but rather to use standard-model hashing. We will not limit ourselves to non-cryptographic hashing but ultimately want to drop idealized models.

%A natural question is whether CI hashing will play a role here.  While it is possible,  using CI in known ways~\cite{} would require a \emph{multi-input-output} form of CI for which we have no positive results and is impossible in general. We hope to lean on novel hardness assumptions on the underlying group as well, and we anticipate PGG and its variants will back our assumptions.  To our knowledge, this is the first time that UCE/PGG-style notions will be used for signature schemes.
%We propose a through study of the scheme using various models and assumptions, ultimately working our way towards results in the standard model. 

\subsubsection{Our Prior Results}
\heading{Tight reduction from Schnorr to CDL.}
At CRYPTO 2025~\cite{C:CFOS25}, we resolved the above question in the affirmative by  introducing the \emph{circular discrete-logarithm} (CDL) problem, a strengthening of DL that is non-interactive and falsifiable but crucially \emph{representation-dependent}. That is, we showed that if CDL holds in the underlying group, then Schnorr signatures are tightly secure in the ROM. The results extend to the multi-user setting~\cite{C:KilMasPan16}.

 %While our assumption inherits a peculiar structure of Schnorr, we stress that our results improve on prior work in two substantial directions, firstly by using a \emph{non-interactive and falsifiable} assumption and secondly by giving a \emph{completely tight reduction}.
% show further applicability of CDL by giving a tight proof of security for the recent three-round threshold-signing scheme $\Sparkle$~\cite{C:CriKomMal23,EPRINT:CriKomMal23} for Schnorr signatures. (Their full version~\cite{EPRINT:CriKomMal23} corrects an error in the proceedings version.) This is relevant as $\Sparkle$ has seen standardization efforts~\cite{nist:schnorr}. 
%Using carefully chosen idealized models that idealize different aspect of our assumption in appropriate elliptic-curve groups (\emph{e.g.}, twisted Edwards curves), we show that breaking CDL has essentially the same complexity as breaking DL. Consequently, if one believes in these models for analyzing DL/CDL in appropriate elliptic-curve groups, then the security level of Schnorr in such groups shown by our results exactly matches that indicated by decades of cryptanalysis.% Our results similarly establish that there is unlikely any faster attack on $\Sparkle$ than just solving DL.   %indicating that instantiating the scheme requires different techniques.
%We note  that CDL can likely be adapted to other $\Sigma$-protocols as well. %(Alternatively, we could replace $g_0$ with random generator $g$.)
%A natural question is whether CDL is related to PGG.  We have not been able to establish a relation. 
% Consider the following natural attempt at showing PGG implies CDL. Let $\advB$ be a $(\GG,f)$-CDL solver.  Consider the PGG adversary $\advA = (\advs,\advd)$ where $\advs$ on input $\params$ chooses $x \gets \Z_p$ and queries $\Exp$ with $f(x)$ to get $X$.  It returns $L = x \circ X$.  Then $\advd$ on inputs $g,L$ runs $\advB$ on $g,L$ to get $z,R$. It returns 1 iff $g^z = R \circ L^{f(R)}$.  Note that $\advs$ is a \emph{masking source} as defined in~\cite{TCC;BFHO22}.  Unfortunately, this idea does not work because $\advB$ may succeed regardless of the challenge bit for $\advA$.  
% As such, in addressing this problem, we will also consider variants of CDL and PGG.

To define CDL, let $\GG$ be as above, and let $f \colon \GG \to \Zp$ be an efficient function, which we call a \emph{conversion function} following the terminology regarding ECDSA~\cite{ecdsa}. We say that the \emph{circular discrete-logarithm} (CDL) problem holds in $\GG$ for $f$ if, given uniformly sampled $h \in \GG$, it is hard to find $(R,z)$ such that:
\begin{newmath}g^z = R \cdot h^{f(R)} \end{newmath}%
and $f(R) \neq 0$. One can think of $f$ as a very simple function that outputs some bits of its input, not necessarily a cryptographic hash function.  We denote CDL for $f$ in $\GG$ as $\CDL[\GG,f]$. The ``circularity'' in CDL is that in the solution equation $R$ occurs both in the base and the exponent.

In our result,  $f$ is used \emph{only in the proof}. Thus, our result also holds under a non-falsifiable but  weaker variant of CDL that asks that there just \emph{exists} some $f$ such that $\CDL[\GG,f]$ holds.  %An equivalent formulation is, given $h \getsr \GG$, find $a \in \Zp$, $b \in \Zp^*$ such that $f(g^a/h^b) = b$. Indeed, this immediately yields a CDL solution with $z = a$ and $R = g^a/h^b$.

%Note that CDL is representation-dependent. Indeed, \eqref{eq-cdl} may hold for one group representation but not another, as the value of $f(R)$ depends on the representation.
\heading{CDL in the EC-GGM.} Note that regardless of $\GG$, $\CDL[\GG,f]$ cannot hold for \emph{every} $f$. In particular, it is necessary that every non-zero value has a small preimage under $f$. We proved that this condition is also \emph{sufficient} for $\CDL[\GG,f]$ in the elliptic-curve generic group model (EC-GGM)~\cite{EC:GroSho22}. Our analysis provides strong evidence that the simple and well-studied, 2-to-1 \emph{ECDSA conversion function} $f_\mathsf{ecdsa}(x,y) \mapsto x \bmod p$ is a good choice for $f$ in  elliptic-curve groups.
For technical reasons, it is necessary to use the EC-GGM to analyze CDL rather than the standard GGM or AGM. 
%(In the GGM there is no notion of ``coordinates'' of a group element, and the AGM disallows mapping group elements to bits.)

We stress that the  upshot of CDL is that, compared to OMDL/MBDL, CDL is \emph{significantly easier to validate via cryptanalysis}. Even showing \emph{any} DL-hard group $\GG$ and \emph{any} conversion function $f$ with small preimage sets for which there is a attack on $\CDL[\GG,f]$  would be interesting. %This measuring stick provides   rigorous justification that Schnorr signatures as deployed face no better attacks than discrete-log itself.

%For example, the 2-to-1 \emph{ECDSA conversion function} $f_\mathsf{ecdsa}(x,y) \mapsto x \bmod p$ does not make sense in Shoup's GGM because a group element does not have an ``x-coordinate.'' 

\subsubsection{Analyzing CDL in the Standard Model} \label{sec-cdl-std}

To strengthen our confidence in CDL, we will conduct the first analysis of it in the \emph{standard}  model:
\begin{ques}
    Can we construct a group $\GG$ and conversion function $f \colon \GG \to \Z_p$ such that $\CDL[\GG,f]$ holds in the standard model under non-interactive and falsifiable assumptions? 
\end{ques}


%Besides corroborating CDL, answering this would be a step toward \emph{instantiating} Schnorr signatures, which we address later.
%Indeed, if we are able to instantiate Schnorr signatures with $H$ then we could solve this question, simply by taking $f$ to me $H$ hard-coded with the empty message $\varepsilon$.
Ideally, taking $\GG$ to be any DL-hard group, we would like to construct a conversion function $f$ such that $\CDL[\GG,f]$ holds under suitable assumptions.

\heading{Insufficiency of correlation intractability.}   For the knowledgeable reader, we discuss why it does not suffice to assume $f$ is  (an achievable form of) \emph{correlation intractable} (CI)~\cite{JACM:CanGolHal04} for CDL to hold. CI of a hash function $\hashH$ for a binary relation $\calR$ says that given the hash key $K$, it is hard to find $x$  such that $(x,\hashH(K,x)) \in \calR$. While general  CI  is impossible~\cite{JACM:CanGolHal04}, more recent work shows how to construct it for binary relations, see \emph{e.g.}~\cite{STOC:CCHLRRW19,EC:CCRR18,TCC:CanCheRey16}.

A canonical application of such CI hash functions is to instantiate Fiat-Shamir to get non-interactive zero-knowledge arguments from  $\Sigma$-protocols. As CDL can be  viewed  as a kind of  Schnorr instantiation  where there is no message and no signing queries, it  is tempting to conclude it suffices  for $f$ to be CI for  binary relations. (Here we  make $f$ keyed, which is not a problem for our results.) The issue is that we need CI  for a relation where a witness is a discrete-log, but since every element has a discrete-log there are no false instances. % $\calR_{\mathsf{Sch}}(X,Y)$ that returns 1 iff there  is a discrete-log of $X \cdot h^{Y}$. But such $\calR_{\mathsf{Sch}}$ is trivial, as all group elements have a discrete-log.
Thus,  there is no such CI hash function. 

%The core issue is that Schnorr's $\Sigma$-protocol doesn't satisfy meaningful soundness.

\heading{Using obfuscation.} Our approach will be to \emph{directly} analyze  \emph{specific}  CI hash functions as candidate $f$'s.  That is, we set $f$  to a ``high powered'' function to vet CDL. (Recall we just need a good $f$ to \emph{exist} and only use it in the proof.)  We will use an indistinguishability obfuscator~\cite{C:BGIRSVY01,SIAM:BGIRSVH16}  $\iO$ and puncturable pseudorandom function~\cite{AC:BonWat13,EC:BoyGilIsh15,CCS:KPTZ13}  $\PRF$. (While not falsifiable, iO is implied by falsifiable assumptions~\cite{STOC:JaiLinSah21,EC:JaiLinSah22}.) We aim to show:

\begin{conj}
    Let DL hold in $\GG$. Let $f := \mathsf{iO}(\PRF)$. Then $\CDL[\GG,f]$ holds.
\end{conj}

%Indeed, such $f$ satisfies many properties of a random oracle, including CI. % To show CDL, we may need assumptions on $\GG$ beyond DL. %Another possibility for constructing $f$ we will consider is extremely lossy functions~\cite{C:Zhandry16}.
Note that~\cite{TCC:CanCheRey16} show such $f$ satisfies  CI for certain relations.
Even though $f$ is not used in the scheme, using iO to build $f$    affects the tightness of our reduction because it depends on the time to run $f$. Therefore, we will seek  more efficient constructions of $f$.  % Other possibilities we will consider for building $f$ include extremely lossy functions~\cite{C:Zhandry16} and extractable functions~\cite{TCC:CanDak08}.


\iffalse
\heading{CDL as a stepping-stone.} In the proposed work, we will study:
\begin{ques}
For every group $\GG$ of order $p$ (or those of interest, say elliptic-curve groups) there is a standard-model function $f \colon \GG \to \Zp$ such that there is a reduction from $\CDL[\GG,f]$ to some more-standard assumption in $\GG$ like DL?    
\end{ques}

In particular, constructing $f$ for which there is a tight reduction from $\CDL[\GG,f]$ to DL in $\GG$ would, by composition, yield a tight reduction from EUF-CMA of $\SchScheme[\GG]$ to DL in $\GG$, resolving a decades-old open problem. Interestingly, this would \emph{not} contradict the above-mentioned impossibility results because our reduction is \emph{non-generic}  since it computes $f$.  Even a loose reduction here would corroborate CDL. %In general, a construction of $f$ could introduce other assumptions. We stress that Schnorr signatures themselves and their instantiation in practice using cryptographic hashing would remain unaffected.
\fi

%In Section~\ref{sec:ECGGM}, 
%We show that, when modeling $\GG$ as an \emph{elliptic-curve generic group} \cite{EC:GroSho22}, this assumption on $f$ is not only necessary but also sufficient for CDL to hold.
 %For an elliptic-curve group, a simple example of a function $f$ satisfying this assumption (as previously argued in~\cite{EC:GroSho22}) is the \emph{ECDSA conversion function}~\cite{ecdsa}, which maps a point $P=(x,y)$ on the curve to its reduced $x$-coordinate, \emph{i.e.}, $f\colon(x,y)\mapsto {x\bmod p}$.


\subsection{Tight Security of     Advanced Schnorr Functionalities}

We extend our treatment from basic Schnorr to advanced primitives built around it. In particular, we consider threshold Schnorr signatures, $\mathsf{Dilithium}$, and $\mathsf{Bulletproofs}$ --- core schemes driving NIST’s threshold and PQC standardization efforts as well as blockchain adoption.
%\begin{ques}
 %   Can we use CDL to show tight ROM security of  \emph{advanced} Schnorr functionalities?
%\end{ques}


%Such  functionalities include threshold signatures~\cite{SAC:KolGol20}, aggregate signatures~\cite{CTRSA:CGKN21}, blind signatures~\cite{EC:FucWol24}, multisignatures \cite{C:NicRufSeu21}, ring signatures~\cite{C:YELAD21}, and adaptor signatures~\cite{AC:AEEFHMMR21}. %We believe CDL will give us a unified way of obtaining \emph{tight} security proofs for such schemes under suitable strengthenings of DL. %Schnorr also has analogues in the lattice-based setting~\cite{} and is the basis for some succinct non-interactive argument systems like Bulletproofs~\cite{}. Below we discuss our preliminary work on some such primitives.

 \subsubsection{Tight Security of Threshold Schnorr Signatures}
 
%One focus will be \emph{threshold Schnorr signatures}. NIST has a second call for their standardization~\cite{}. 

\heading{Threshold Schnorr signatures.}
\iffalse
Threshold signatures (TS) originated in~\cite{C:Desmedt87,C:DesFra89}. In a TS scheme, there is a single verification key $\vkey$, a set of $n$ signers that each possess a share $\skey_i$ of the corresponding signing key $\skey$, and a threshold $t$. If at least $t$ signers get together they can produce a signature on a message of their choosing that verifies under $\vkey$. One distinction in security models for TS signatures is consideration of \emph{static} vs.~\emph{adaptive} corruptions.  In the former, the adversary must declare all the parties it controls at the outset of the game. In either case, the adversary can initiate concurrent executions of the protocol among subsets of parties of its choice.  Then, the adversary wins if it forges a signature on a message it never initiated the protocol on. %There are strengthenings to this definition we will also consider~\cite{}.

In practice, the most desirable TS schemes are threshold Schnorr signatures, meaning signatures produced are Schnorr ones. 
Threshold Schnorr signatures are central to both practice (\emph{e.g.}, blockchains, Bitcoin’s migration from ECDSA) and standards (NIST’s second call for threshold signatures). Yet their provable security status remains clouded across candidate schemes.
\fi
Threshold signatures allow any $t$ of $n$ parties to jointly produce a signature using their secret key shares that is valid under the group public key. (Public-key shares are usually published as well, or are implicitly revealed). The most relevant  schemes for practice are threshold \emph{Schnorr} signatures. Such schemes are  central to blockchains (\emph{e.g.}, motivating Bitcoin’s move from ECDSA to Schnorr) and targeted in NIST’s threshold signature efforts.

%In fact, one primary reason for Bitcoin's transition over to Schnorr   from ECDSA  is that  Schnorr signatures have a linear signing structure that lends itself more easily to  thresholdization. As a result of a surge of interest in such schemes, many candidates have recently been proposed~\cite{SAC:GolKom20,C:BCKMTZ22,EPRINT:CriKomMal23,C:CriKomMal23,C:CGRS23,EPRINT:BelTesZhu22,EPRINT:Lin22,EPRINT:Makriyannis22,CCS:RRJSS22,C:KatReiSak24,EPRINT:BDLR24}. %These schemes are more practical as compared to early works because they do not run expensive verifiable secret sharing every time a new signature should be generated but only once at setup.   %A main focus has been on reducing the number of rounds. In this regard, the best is the ``partially non-interactive'' scheme $\mathsf{FROST}$~\cite{SAC:GolKom20}, which has a number of variants that offer tradeoffs in performance with  robustness and network synchronicity assumptions~\cite{EPRINT:CriKomMal23,EPRINT:BelCheZhu22,CCS:RRJSS22,CRYPTO:CJRS22,EPRINT:KomGol24}. The scheme is  also specified in an IETF RFC~\cite{xxx} and seeing standarization efforts~\cite{xxx}.

\heading{Our focus.}
We analyze  the major standardization contenders $\mathsf{FROST}$~\cite{SAC:GolKom20,rfc9591} (two rounds) and $\mathsf{Sparkle}$~\cite{C:CriKomMal23}, $\mathsf{Gorgos}$~\cite{C:BDLR25} (three rounds), focusing on tightness under static and adaptive corruptions.  In plain terms, adaptive security means \emph{attackers cannot profit from corrupting signers after seeing the protocol in action}; it is explicitly desired by NIST. We  find:
\begin{quote}
  \emph{Circular security assumptions like CDL are a perfect match for analyzing adaptive security of threshold Schnorr signature schemes,  in particular $\Sparkle$.}
     \end{quote}
\noindent   
%\begin{ques}
 %   In what corruption models is tight  security  possible for these schemes, and under what assumptions?
%\end{ques} \label{ques-adap}
%Indeed, we will aim to resolve some fundamental questions that remain about even loose security proofs for such schemes.
%Indeed, tightness for threshold Schnorr signatures is as relevant as it is for basic Schnorr signatures, but it is rarely discussed in the literature and current proofs are loose.
 
  %$\mathsf{Gorgos}$~\cite{} is  not key-unique and enjoys loose adaptive security under DDH. 
 

%\heading{Tight static security of $\Sparkle$.}  %This result demonstrates the promise of CDL in the context of threshold Schnorr signatures.
$\Sparkle$ is a natural scheme explicitly designed for adaptive security. 
The proof of this in the algebraic group model (AGM) under algebraic one-more discrete-log (AOMDL) 
claimed in~\cite{C:CriKomMal23}, however, is flawed. As our main result here, 
we aim to recover a  proof under new assumptions.

\heading{Results for  $\mathsf{FROST}$.}  Bellare \emph{et al.}~\cite{C:BCKMTZ22}  give a loose  proof of static security for $\mathsf{FROST}$ under one-more DL (OMDL)~\cite{JC:BNPS03}, where the adversary makes DL queries. While one would like to avoid interactive assumptions for static security,  we  will show $\mathsf{FROST}$ is deficient here:
\begin{conj}
    There is no black-box reduction from  security of  $\mathsf{FROST}$ under static corruptions to any non-interactive (and perhaps falsifiable) assumption.
\end{conj}
%As a first step, we have proven this for the DL assumption in the case of ``canonical'' reductions that use the DL instance as the forger's public key.
Recent related work~\cite{EPRINT:CriKomMal25} shows there is no black-box proof of adaptive security for any ``key-unique'' scheme under a non-interactive assumption.
 Our desired result above is incomparable  because it  concerns  \emph{static} security but is specialized to  $\mathsf{FROST}$. 
 
In both $\mathsf{FROST}$ and $\Sparkle$, public-key shares are Shamir shares in the exponent. Recent work~\cite{C:CriSte25} introduces a ``Problem P'' that must be hard if any such scheme is to achieve \emph{adaptive} security, at least if we allow up to $t-1$ corruptions (which we assume below). Follow-on work~\cite{C:KKTZ25} shows  $\mathsf{FROST}$ is adaptively secure under  Low-Dimensional Vector Representation (LDVR) problem, which implies Problem P and is \emph{necessary} for adaptive security of $\mathsf{FROST}$ and $\mathsf{Sparkle+}$. They show   $\mathsf{FROST}$ is adaptively secure under AOMDL and LDVR in the AGM. %While we do not propose any new results on adaptive security of $\mathsf{FROST}$, this will be relevant for our analysis of of $\Sparkle$ below.  %As being two rounds is the main advantage of $\mathsf{FROST}$,  an interesting question we will consider is: \emph{Is there is a  modification to the scheme that preserves the efficiency (\emph{e.g.}, number of rounds) and  admits a proof of static security under  a non-interactive assumption like CDL?}

\heading{Results for  $\Sparkle$.} Crites \emph{et al.}~\cite{C:CriKomMal23} show static security of $\Sparkle$ in the ROM just under  DL.  At CRYPTO 2025~\cite{C:CFOS25}, we further showed its \emph{tight} static security under CDL.  %We stress that is not known how to reduce security of  such schemes that of basic Schnorr. %, so this result does not follow from our previously-mentioned result about basic Schnorr.

We aim to give the fist proof of  \emph{adaptive} security of $\Sparkle$, which is  even \emph{tight}, under an interactive  extension of CDL (+LDVR). 
We will define  the \emph{Vandermonde circular  discrete-logarithm} (VCDL) assumption. Let $\GG,f$ be as in CDL. The adversary gets group elements $(X_0,X_1,\ldots,X_t)$ where $X_i = g^{x_i}$, and it can make up to $t-1$ queries to an oracle that on $k \in \Z^*_p$ outputs $x_0 + \sum_{i=1}^t k^i \cdot x_i$. Think   of $x_1,\ldots,x_t$ coefficients of a  secret-sharing polynomial, and a query reply as revealing a share. The adversary must return a CDL solution wrt.~$X_0$, \emph{i.e.}~$(R,z)$ such that \smash{$g^z = R\cdot X_0^{f(R)}$}. 


Despite its interactivity, VCDL remains amenable to cryptanalytic validation not only 
because the assumption is falsifiable, but also because the adversary’s queries must take 
very simple \emph{Vandermonde} form rather than arbitrary representations as in AOMDL. 
Indeed, we found that if arbitrary representations were allowed, VCDL would be false.  At a high level, we want to prove: %So VCDL still gives a good handle on security of the higher-level scheme being analyzed. 
\begin{conj}  \label{conj-spar}
(1) VCDL implies LDVR, establishing a  \emph{hierarchy}: VCDL implies LDVR implies Problem P. (2) As a partial converse, LDVR implies  VCDL in the elliptic-curve generic group model. (3) $\Sparkle$ meets \emph{tight  adaptive security} in the ROM under VCDL. %The intuition is that VCDL queries are used to handle the forger's corruption queries.
\end{conj}
%To explain what we mean, note that  LDVR can be seen as a problem over $\Z_p$ that  matches $\mathsf{FROST}$ in that it captures such a  problem that must be solved if their proposed reduction to OMDL fails in the AGM. We want to do the same for $\Sparkle$, but will actually do so  directly for COMDL in the EC-GGM rather than the full-fledged scheme. Figuring out the right Problem X will be a lynchpin of the proposed work and give actionable insight into the security of $\Sparkle$.

%The most challenging step is (2). To gain intuition, we will study how LDVR is used in the proof of~\cite{} that OMDL + LDVR implies adaptive security of $\mathsf{FROST}$ in the AGM, and why this doesn't work for $\mathsf{Sparkle+}$. We may also need to modify LDVR for our results, as OMCDL doesn't  %The formulation of COMDL is thus delicate. % We also believe that the natural relation to Feldman's VSS makes COMDL attractive to the cryptanalyst.
%In other words, these results will show that adaptive security of $\mathsf{Sparkle+}$ is fundamentally related to adaptive security of Feldman's VSS and LDVR.
%A notable a feature of Problem P and LDVR is they do not involve a group, just $\Z_p$. This is not the case for COMDL, which explains why we will need an additional assumption to prove COMDL  in the EC-GGM.  %While it is not as common to prove such conditional results in the GGM/EC-GGM, it was done in the case of \emph{e.g.} on Schnorr's blind signature scheme~\cite{}

%While it may seem odd to have a conditional proof in the EC-GGM, it makes sense here because LDVR is a problem that only deals with the underlying field $\Z_p$. %such analyses have been done before, \emph{e.g.}~the classic proof that blind Schnorr signatures are secure in the GGM assuming ROS. (Although ROS turned out to be false.)
The intuition for (3) is that oracle queries allow the reduction to simulate the adversary’s 
corruption queries. We use LDVR from~\cite{C:KKTZ25}, but unlike their result on 
$\mathsf{FROST}$ we shift the idealization beyond the ROM to the analysis of the 
\emph{underlying} assumption rather than the scheme.

A key point is that the ``circularity'' in VCDL is not just required for \emph{tight} adaptive 
security --- \emph{it is essential for proving adaptive security at all}. The central hurdle in analyzing 
any ``key-unique'' threshold Schnorr scheme is \emph{rewinding}: across two executions the 
forger may corrupt \emph{all} parties. Circular assumptions avoid rewinding and thus resolve 
this issue \emph{in addition} to enabling tightness. This works for $\mathsf{Sparkle+}$ 
but \emph{not} to $\mathsf{FROST}$ because the reduction would need to make queries to answer not just the adversary's
corruption but \emph{signing} queries, which may not be Vandermonde. %Thus, one would need COMDL with unrestricted oracle queries here, which as we said makes the assumption false. In sum, our projected results will bear out the intuition that $\Sparkle$ has better security than $\mathsf{FROST}$.

\heading{Results for $\mathsf{Gorgos}$.} Finally, we will  prove \emph{tight}  adaptive security of $\mathsf{Gorgos}$ --- which is not key-unique  --- under DDH + CDL. In the existing  proof, key shares are ‘`rigged'’ via DDH, followed by a loose, rewinding reduction to DL. We will replace the latter with a tight reduction to  CDL. %This serves to underscore its superior security properties relative to the previously mentioned schemes.

 
 
\heading{Summary of the results.} % Our work will  directly address NIST’s  criterion of adaptive security, push tightness as an additional criterion, and help practitioners make the best choice among tradeoffs for their application.
In summary, our planned results are:

\begin{table}[H]
\centering
\renewcommand{\arraystretch}{1.2}
\begin{tabular}{|l|l|l|l|}
\hline
\textbf{Scheme} & \textbf{Security Type} & \textbf{Assumption / Model} & \textbf{Source} \\ \hline
$\mathsf{FROST}$ & Loose static & One-More DL (OMDL) & \cite{C:BCKMTZ22} \\
                 & Loose static & Non-interactive (\emph{impossible}) & \emph{Our work} \\
                 & Adaptive & Algebraic OMDL + LDVR + AGM & \cite{C:KKTZ25} \\ \hline
$\Sparkle$       & Tight static & Circular DL & \emph{Our work} \\
                 & Tight adaptive & Vandermonde Circular DL & \emph{Our work} \\ \hline
$\mathsf{Gorgos}$ & Tight adaptive & DDH + Circular DL & \emph{Our work} \\ \hline
\end{tabular}
\end{table}
\noindent As the schemes in differ in efficiency, our results allow security vs.~efficiency tradeoff selections.
%In our view, our planned results make $\Sparkle$, which is less efficient than $\mathsf{FROST}$ but more efficient than $\mathsf{Gorgos}$, an attractive middle ground  if three rounds is acceptable and one buys LDVR.
\iffalse
\begin{enumerate}
\item $\mathsf{FROST}$: Loose static security  under OMDL (prior work) and impossible under non-interactive assumptions (our proposed work);  adaptive security  under AOMDL + LDVR in the AGM (prior work).
\item $\Sparkle$: Tight static security  under CDL (our CRYPTO 2025 work);  tight adaptive security  under VCDL (our proposed work).
\item $\mathsf{Gorgos}$: Tight adaptive security  under DDH + CDL (our proposed work).
\end{enumerate}
\fi

%In our view, this clarifies that $\Sparkle$, which is roughly twice as efficient as $\mathsf{Gorgos}$, is an attractive choice in terms of efficiency-security tradeoffs when three rounds can be tolerated.
%Thus, our results will bear out an ordering in terms of increasing security: $\mathsf{FROST}$, $\Sparkle$, then $\mathsf{Gorgos}$. %This will underscore the strong provable security support for the scheme when using practical parameters.

\iffalse
\heading{Main takeaway.} To summarize, we will show $\mathsf{FROST}$ cannot be proven secure under static corruptions using any non-interactive assumption, while $\mathsf{Sparkle+}$ can be proven tightly secure under adaptive corruptions using COMDL + LDVR. Finally, we will show $\mathsf{Gorgos}$ can be proven tightly secure under adaptive corruptions using  DDH + CDL.
\fi

% The fact that this is not the case for $\Sparkle$ reflects its design rationale for achieving adaptive security as mentioned above.

% In the proposed work, we will consider tight security of additional threshold Schnorr signatures beyond $\Sparkle$ like $\mathsf{FROST}$~\cite{} and $\mathsf{Gorgos}$~\cite{}, as well as the case of \emph{adaptive} security.  %In the proposed work, we will extend the treatment to other high-profile threshold Schnorr signature schemes as well as consider adaptive security.

%While the main objective is to obtain \emph{positive} results,  to understand the problem better we will look at impossibility results as well.  %The only other tight security proof for a Schnorr signature scheme we are aware of %Indeed, while some threshold Schnorr signatures have been designed to have a tight security proof, we are interested in analyzing security of the most prominent threshold Schnorr signature schemes, which do not. In part, this is because we view it as unlikely that practitioners will switch to a scheme just for a tight security proof. Indeed, these tightly secure threshold Schnorr signatures are often actually  less efficient, either in the number of rounds, or in computation when non-tight schemes have their  parameters set as if they were tight. We thus ask:
\iffalse
\begin{ques}
    Can we show tight security in the ROM of popular threshold Schnorr signature schemes assuming CDL or suitable extensions to it?
\end{ques}

\heading{Tight reduction from  $\Sparkle$ to CDL.}
In the work introducing CDL~\cite{}, the PI took an initial step and proved  that $\Sparkle$~\cite{EPRINT:CriKomMal23}, a three-round threshold Schnorr signature scheme, is tightly secure in the ROM under static corruptions assuming CDL. This  demonstrates the relevance of CDL beyond basic Schnorr signatures and specifically to threshold Schnorr signatures. We plan to build off of this result in the proposed work. 
\fi
%$\Sparkle$ is natural and efficient, combining Shamir secret sharing of the secret key in the exponent with an initial round of hash-based commitment to the signers' randomness. To start with, we consider a \emph{static} corruption

%Highly efficient such schemes have been introduced in recent years, including~\cite{SAC:GolKom20,C:BCKMTZ22,EPRINT:CriKomMal23,C:CriKomMal23,C:CGRS23,EPRINT:BelTesZhu22,EPRINT:Lin22,EPRINT:Makriyannis22,CCS:RRJSS22,C:KatReiSak24,EPRINT:BDLR24}. A main focus has been on reducing the number of rounds. In this regard, we focus on the ``partially non-interactive'' scheme $\mathsf{FROST}$~\cite{SAC:GolKom20}, which has a number of variants that offer tradeoffs in performance with other properties like robustness and network synchronicity assumptions~\cite{EPRINT:CriKomMal23,EPRINT:BelCheZhu22,CCS:RRJSS22,CRYPTO:CJRS22,EPRINT:KomGol24}. This scheme is   important to analyze because it is specified in an IETF RFC~\cite{xxx} and seeing standarization efforts~\cite{xxx}.

\iffalse
In preliminary work we have focused  on  $\mathsf{FROST}$, a  two-round threshold Schnorr signature scheme  specified in an IETF RFC~\cite{} and seeing standardization efforts~\cite{}. There are actually multiple versions: $\mathsf{FROST1}$~\cite{}, $\mathsf{FROST2}$~\cite{},  and $\mathsf{FROST3}$~\cite{}, which we refer to generically as $\mathsf{FROST}$. %We first discuss the scheme and what's known about it security. % As a main result here, we will show a simple modification to $\mathsf{FROST}$ such that, while the first round is no longer message independent, (potentially tight) static and adaptive security are provable under \emph{non-interactive} assumptions. Moreover, the scheme is more bandwidth-efficient than $\mathsf{FROST}$.


\iffalse
\heading{The scheme.}
 In $\mathsf{FROST}$, the Schnorr secret key $x \in \Z_p$ is initially shared via $t$-out-of-$n$ Shamir secret sharing, so that party $i$ gets share $x_i \in  \Z_p$ for all $i\in [n]$. (The $g^{x_i}$ are typically also published \emph{e.g.}~to achieve identifiable abort.) To sign, a set $\calS$ of at least $t$ signers first engage in a message-independent round of preprocessing, whereby each signer $i\in \calS$ sends a ``presignature'' $\psig_k = (g^{r_k},g^{s_,k})$ for random $r_k,s_k \in \Z_p$, retaining $(r_k,s_k)$ in its state.  The execution of the protocol is coordinated by a ``leader.'' The leader sends $(\calS,m, \{\psig_j\}_{j \in \calS})$ to each signer $m \in \calS$. In the second round, signer $k\in\calS$ sends $(R,z)$ where
$$R \gets \prod_{i \in \calS} R_k \cdot S_k^{d_k}~~~\text{and}~~~z_k \gets (r_k + s_k \cdot d_k + \lambda_k^\calS \cdot \skey_k \cdot c) \bmod p\;,$$ where $d_k, c$ are derived via ROs $H_1,H_2$, \emph{e.g.} in $\mathsf{FROST1}$ we have $d_k \gets H_1(\vk,m,\calS,\{R_i,S_i\}_{i \in \calS}, i)$ and $c \gets H_2(\vk,m,R)$.  To increase throughput,  each party will actually initially  execute the first round many times and publishes a list of  presignatures. In an online  round, the leader chooses a presignature from each signer's list; each such presignature must be only used once.
\fi 

% Thus, using this technique, $\mathsf{FROST1}$ and $\mathsf{FROST2}$ need only a single round of online communication once the message becomes known. % The ephermeral public keys are cleverly designed to prevent attacks in this case like the ROS attack~\cite{xxx}.

\heading{Existing analyses.}  Bellare \emph{et al.}~\cite{EPRINT:BelTesZhu22} show that $\mathsf{FROST}$ is secure under \emph{static} corruptions in the ROM using the one-more discrete logarithm (OMDL) assumption~\cite{xxx}. In OMDL, the adversary gets \emph{many} DL-challenges $(X_1,\ldots,X_\ell)$ where $x_i \gets \Z_p$, and must output $(x_1,\ldots,x_t)$ while making at most $\ell-1$ queries to a DL-oracle that on input a group element $Y$ returns the DL of $Y$.
 %Thus, in addition to being interactive, OMDL is \emph{non-falsifiable}~\cite{C:Naor03} because the game is not  efficient.
 Crites \emph{et al.}~\cite{} show that $\mathsf{FROST}$ is (loosely) secure in the ROM under OMDL for up to $t/2$ \emph{adaptive} corruptions ($t$ is the signing threshold). For  higher corruption thresholds they require the AGM and a new assumption they call the  low-dimensional vector representation problem (LDVR). %Introducing LDVR problem is  necessary because of an ``attack'' under adaptive security on threshold Schnorr signatures when partial verification keys are published as in $\mathsf{FROST}$. %(It is not known if the attack is efficient in general.) The security loss in the above results is about the same as in the classic results for basic Schnorr.

% (Indeed, we cannot use AGM with OMCDL in the first place, because in the AGM one cannot map group elements to bits like the conversion function in CDL does.) 
%We say ``some version of'' above because, as we will see, it is not obvious how to define OMDL so that it is both useful and achievable. 

%Second, Shoup~\cite{DBLP:journals/cic/Shoup25} shows breaking $\mathsf{FROST}$ is  as hard as DL in the ROM + \emph{elliptic-curve generic group model} (EC-GGM)~\cite{EC:GroSho22}, which as previously discussed we want to avoid. %like $\mathsf{FROST}$, \emph{i.e.}, it would be preferable only to analyze a simple assumption in the EC-GGM.

\heading{Impossibility of loose security under DL.} Reductions in the above results are loose and rely on  interactive, non-falsifiable assumptions.  To better understand the space of possible assumptions one can use here, we ask whether using  interactive assumptions to prove static security is necessary. To start with:
\begin{ques}
Can $\mathsf{FROST}$ be proven (even loosely) secure under static corruptions using DL? 
\end{ques} 


Here we will aim to give \emph{metareduction} ruling out such algebraic reductions. Note that while security of a threshold Schnorr signature scheme implies  security of basic Schnorr, the desired result here does \emph{not} follow from existing metareductions for tight security of basic Schnorr. This because we want to rule out even \emph{loose} reductions in  our case.

%\begin{ques}
 %   Can $\mathsf{FROST}$ be proven tightly secure under static corruptions using OMCDL?
%\end{ques}

\begin{ques}
    Does $\mathsf{FROST}$ admit a simple and efficiency-preserving modification for which we can prove (even loose) security under static corruptions using CDL?
\end{ques}
\fi


\iffalse
\heading{A modification to $\mathsf{FROST}$.}
In light of such a result, we  ask: 
\begin{ques}
Can we modify $\mathsf{FROST}$ so that we can prove static or even adaptive security of the resulting under a \emph{non-interactive} assumption? 
\end{ques} 
Indeed, it is early in the adoption of threshold Schnorr signatures in practice, giving us more flexibility in introducing  such modifications.  We propose  a modification that, while it makes the first round message-dependent (so presignatures can no longer be preprocessed), the presignatures are smaller and security could be proven under CDL. The idea is that we require   the set of signers $\calS$ and the message $m$ to be sent by the leader in the initial message, and that a presignature includes $\calS$ and $m$.  But we can also eliminate one group element in a presignature, so that in our modified $\mathf{FROST}$ a presignatures looks like $\psig_k = (S= g^{s_k},\calS,m)$. Naturally, such a presignature can then only be used to sign the particular message $m$ with the set of signers $\calS$. 

Interestingly, we sketch a ``semi-tight'' rather than tight reduction from static security of our modified $\mathsf{FROST}$ to CDL. To explain the terminology here,  recall that given a forger against the Schnorr signature scheme in the ROM with success probability $p_{\mathsf{succ}}$, the PS result~\cite{}  says there is an adversary solving DL with similar running time and having success probability $\Theta(p_{\mathsf{succ}}^2/q_H)$, where $q_H$ is the number of RO queries made by the forger. We call an improvement ``semi-tight'' if the adversary solving DL in the underlying group has similar running-time but success probability $\Theta(\mathsf{succ}/q_H)$. The idea is that when simulating a presignature of an honest signer we can program the outputs of $H_1$ and $H_2$, namely $d_k$ and $c$
\fi
%\begin{ques}
 %   Can $\mathsf{FROST}$ be proven \emph{tightly} secure using a notion of \emph{one-more CDL}?
%\end{ques}

\iffalse
\begin{ques}
    Can $\mathsf{FROST}$ be proven adaptively security under \emph{one-more CDL}?
\end{ques}

\begin{ques}
    Can $\mathsf{FROST}$ be proven tightly secure under adaptive  corruptions using a notion of \emph{one-more CDL} (also assuming LDVR for high corruption threshold)? 
\end{ques}

To get some intuition, recall that in basic Schnorr signatures, $R$ in a signature $(R,z)$ is chosen at random \emph{once the message to be signed is already determined}. This is crucial to answering a signing query in the proof: the reduction chooses $R := R'/\vk^c$ where $R'$ is a random element whose DL it knows.  In $\mathsf{FROST}$, the analogue of $R$ is a product of the $R_i \cdot S_i^d$ terms chosen \emph{before} the message to be signed. As such, the forgoing proof methodology breaks down and OMDL is used to simulate a partial signature of honest party $i$ by making DL query $R_iS_i^{d_i} vk_i^{c \lambda_i}$. Obtaining an additional DL solutions of all $(X_1,\ldots,X_\ell)$ uses forking.

To adapt this to a tight proof under OMCDL, we must first have to ask how to define  OMCDL. In preliminary work, we have found it non-trivial to define OMCDL so that it is both useful and achievable.  One possibility is that the adversary gets one challenge $X = g^x$ and it can query a CDL oracle that on input $R$ returns DL of $R \cdot X^{f(R)}$. Say that the adversary wins if it can produce a CDL solution for $X$ different from any of its oracle responses. But making the query $R=X$ the adversary can easily learn $X$ and win the game.  
\fi

\iffalse
 \paragraph*{The quest for tight reductions.} Tightness of a reduction refers to its computational overhead and   determines the security level of a construction for a given key-length.  It is widely conjectured that breaking Schnorr signatures is as hard as solving DL.
Unfortunately, current RO-model reductions to DL are \emph{lossy}; a tight reduction requires the AGM~\cite{AC:BauFucPlo21,TCC:LipParSii23}. Moreover, work culminating with Fleischhacker \emph{et al.}~\cite{AC:FleJagSch14} shows there is no tight, ``generic'' reduction to a ``generic'' non-interactive problem, meaning invariant under isomorphism.   As mentioned above, Bellare and Dai~\cite{INDOCRYPT:BelDai20} show how to avoid the  ``forking lemma''~\cite{JC:PoiSte00} under an \emph{interactive} assumption, but still get a reduction  losing a factor of the adversary's number of RO queries. % A multi-user security analysis of Kiltz \emph{et al.}~\cite{C:KilMasPan16}, while independent of the number of users,  does \emph{not} avoid this  lemma. %We want to eliminate both downsides of this result. 
\fi



\iffalse
\paragraph*{Circular discrete-log in $\Z_p^*$.}
To illustrate our approach, we introduce a problem in $\Z_p^*$  we dub \emph{circular discrete-log} (CDL): Given  $X \getsr \Z_p^*$, output  $z,R \in  \Z_p^*$ such that \smash{$g^z = R \circ X^{R}$};  equivalently, output $a,b \in \Z_p^*$ such that $g^a/X^b = b$.
One can think of CDL as simplifying Schnorr to: (i) empty message space, (ii) identity hash function, (iii) no signing queries. While   impossibility results do apply in general to such a weak setting, they do \emph{not} apply to CDL because it is representation-specific.
In preliminary work, we give a  tight reduction from Schnorr in $\Z_p^*$ to  CDL. 
Unfortunately, we also find an attack on CDL in the auxiliary-input setting. (It is important to consider auxiliary-input (aka.~preprocessing) in this context, as shown by~\cite{C:CLMQ21}.) Namely, in the preprocessing phase, the adversary computes \smash{$a \gets \mathsf{DLog}_{\GG,g}((p-1)/2)$} and returns $a$.  In the online phase, on inputs $X,a$ the adversary sets $z \gets a$ and $R \gets g^a/X^b$ where $b = (p-1)/2$, then returns $(z,R)$. 
\fi

%In $\Z_p^*$, we will have $g$ randomly chosen per-signer and included in its public-key. With such a modification, we conjecture CDL  holds  in the auxiliary-input setting. 
 %Assuming $p = 2p' -1$ for a prime $p'$, we work in $\mathbb{QR}_p^*$, where we also conjecture CDL holds, even for fixed $g$, in the auxiliary-input setting. 
\iffalse 
\paragraph*{Weak instantiability to tight reducibility.} In preliminary work, we give a \emph{completely tight} reduction in the ROM from (strong) UF-CMA of  $\mathsf{Sch}[\GG,H]$ to $(\GG,f)$-CDL for \emph{any} $\GG,f$ where zero does not have  too large a preimage-set under $f$ (though for CDL to hold, no preimage-set can be too large). % an assumption needed for CDL to hold in the first place. 
We should get a tight reduction  in the multi-user setting as well, following~\cite{C:KilMasPan16}.  Note that $f$ is \emph{not} used to instantiate $H$, and our result does not contradict~\cite{AC:FleJagSch14} because $f$ can be representation-dependent. Moreover, $(\GG,f)$-CDL is not only \emph{non-interactive}  but for  fixed $\GG,f$ is also \emph{falsifiable}~\cite{C:Naor03}. Alternatively, we can base our results on the non-falsifiable (but weaker) assumption that for every $\GG$ there exists $f$ such that $(\GG,f)$-CDL holds. We don't even necessarily need to know an explicit such $f$, as $f$ occurs \emph{only in the proof}.
%We fur  thermore stress  that the conversion function $f$  is \emph{not} the same as the hash function ultimately  used to instantiate the  scheme.  %For example, $f$ might be the ECDSA conversion function (see below) while the hash function might be $\mathsf{SHA256}$. %The idea for our reduction is to handle signing queries as in  prior PROM proofs, and  respond to the $i$-th hash query $R_i \| m_i$ with $f(B_i \circ R_i^{a_i})/a_i$
%for $B_i \getsr \GG$ and $a_i \getsr \Z_p^*$. Given forgery $(m^*, R^*,z^*)$, assuming wlog hash query $m^* \| R^*$ was previously made, suitable ``post-processing'' of the forgery allows extracting a CDL solution. Note that the reduction does \emph{not} have a $q_H$ factor.

%While our reduction works for \emph{any} $\GG, f$ such that $(\GG,f)$-CDL holds,
%\begin{task}
%     Characterize $\GG,f$ such that $(\GG,f)$-CDL holds.
%\end{task}
%We conjecture this is the case when the \emph{converse} direction holds: 
%\begin{task}
%     Characterize $\group,f$ such that if $(\GG, f)$-CDL is easy, Schnorr signatures in $\group$ are insecure in the RO model.
%\end{task}
%We call such $f$  the \emph{hardcore} conversion functions of Schnorr in $\group$.
\paragraph*{Using ECDSA's conversion function.} For   $\GG = \mathsf{EC}_p$, we conjecture  the ECDSA conversion function  $f_{\mathsf{ecdsa}}$~\cite{ecdsa}, which outputs the $x$-coordinate of the input  modulo $p$, is  suitable,  connecting security of Schnorr to ECDSA.  In fact, 
  $(\GG,f)$-CDL is  a special case of the \emph{semi-logarithm problem} (SLP)~\cite[Definition 6]{TCC:FerKilPoe17}. % While \cite[Theorem 1]{TCC:FerKilPoe17} gives a reduction from an ElGamal-based signature to the empty-message, no-signing-query setting, compared to our preliminary result for Schnorr:  (1) their reduction is loose, (2) their security notion for the ElGamal-based signature allows only one signature per message, (3) they work in the bijective RO model and model $f_{\mathsf{ecdsa}}$ as a composition of three functions, and (4) their proof techniques are completely different. 
For  $\GG = \mathsf{EC}_p$ and $f = f_{\mathsf{ecdsa}}$~\cite{TCC:FerKilPoe17}  give a (loose)  reduction from SDL to DL in the bijective ROM. %For CDL, we can show a tight such reduction for algebraic adversaries. 
We wish to get better bounds and avoid idealizing $f$. Unfortunately,  the AGM doesn't allow mapping group elements to bits. So, we will focus on the \emph{elliptic-curve} GGM (EC-GGM)~\cite{EC:GroSho22}, which was precisely   introduced to analyze $f_{\mathsf{ecdsa}}$:
\begin{task}
   Show that $(\mathsf{EC}_p,f_{\mathsf{ecdsa}})$-CDL holds in the EC-GGM with the same complexity as DL.
\end{task}


In the EC-GGM, the group encoding is a random  $\sigma \colon \Z_p \to \mathsf{EC}_p$ respecting ``trivial'' relations. Thus, the EC-GGM will  allow us to implement $f_{\mathsf{ecdsa}}$. Compared to the  GGM, using the EC-GGM is also important  because there are secure schemes in the GGM that are \emph{insecure}  in $\mathsf{EC}_p$.  %EC-GGM also lets the adversary ``obliviously sample'' elements.% For example, ECDSA is strongly unforgeable in the GGM but in reality is malleable. 

 \paragraph*{Reducing to CDL to DL.} We believe CDL could  provide a new path to proving tight security under DL. %A \emph{fixed} $f$  such that for \emph{every} DL-hard group $\GG$, $(\GG,f)$-CDL tightly reduces to DL would contradict~\cite{AC:FleJagSch14}. 
 The ``dream result'' we aim for is: for any DL-hard $\GG$, there exists $f$ (that depends on $\GG$ to avoid impossibility results) such that $(\GG,f)$-CDL holds. %(\emph{E.g.}, $f = f_{\mathsf{ecdsa}}$ if $\GG = \mathsf{EC}_p$).   
The construction of $f$ could use additional assumptions like program obfuscation (cf.~Section~\ref{sec-feas}) as long as the reduction is tight, and such primitives would occur \emph{only in the proof}. In preliminary work, we show any CI hash function for ``sparse relations''  satisfies   \emph{weak} $(\GG,f)$-CDL in any DL-hard $\GG$, which is defined as: Given $h,g \getsr \GG$, output $R \in \GG$ such that \smash{$h = R \cdot g^{f(R)}$}.  So, it is plausible that a construction of such CI hash function, like an obfuscated puncturable PRF~\cite{EC:CCRR18}, can be shown $(\GG,f)$-CDL via  a direct analysis.  We could use a group-based PRF that depends on $\GG$ via the group operation. %Task~\ref{task-cdl} will be a special case of this problem.   %We will further investigate feasibility in Section~\ref{sec-beyond}. %; they would \emph{not} be used to instantiate Schnorr in practice.
%Txshe   concrete security bound will allow us to further justify key-length in practice. %Concretely, we aim to show that  $(\mathsf{EC}_p,f_{\mathsf{ecdsa}})$-CDL is basically secure as DL in the EC-GGM. %s a first step toward our dream result:
% \begin{task}
%Give a tight reduction from $(\mathsf{EC}_p,f_{\mathsf{ecdsa}})$-CDL to DL modeling $f_{\mathsf{ecdsa}}$ as in~\cite{TCC:FerKilPoe17}.
% \end{task}
% Another motivation here is analogous to that for showing a new assumption holds in the GGM with the same bounds as for DL. Indeed, Bellare and Dai~\cite{BelDai} show their new interactive assumption is as hard as DL in the GGM. %Unfortunately, using the GGM does not seem to be possible in our case because the conversion function is representation-dependent.



 
%In particular, we aim to resolve the above question affirmatively in the ``elliptic-curve GGM'' of Groth and Shoup~\cite{EC:GroSho22}, by drawing on techniques from their analysis of EC-DSA.  In particular, their is the first analysis of ECDSA that does not idealize the conversion function.
%AAs a first step, we would like to analyze the assumption in  idealized models.  To do so, it makes sense to look  at prior analyses of ECDSA.  Unfortunately, all such analyses~\cite{} either explicitly or implicitly idealize the conversion function, and a recent paper of Hartmann and Kiltz~\cite{TCC:HarKil23} shows that  programmability of the conversion function is in some sense necessary. We do \emph{not} want to idealize the conversion function.  %However, it may be possible to prove that \emph{if} EC-DSA is secure \emph{then} $(\mathsf{EC}_p,\calE_{\mathsf{ecdsa}})$-CDL holds. % Of idealized models used  used for analyzing ECDSA, the EC-GGM of Gorth and Shoup~\cite{EC:GroSHo22} seems the most realistic, and we will also try to analyze $(\mathsf{EC}_p,\calE_{\mathsf{ecdsa}})$-CDL in it. 
%If this is infeasible, we may work in the AGM~\cite{C:FucKilLos18,C:Zhandry22b} to start with. 
%Finally:

\iffalse
\begin{task}
    Give a tight reduction from  $(\mathsf{EC}_p,f_{\mathsf{ecdsa}})$-CDL to DL in the \emph{standard model}.
\end{task}
In particular, this would yield the first \emph{tight} security proof for Schnorr under DL for appropriate elliptic-curve groups.
\fi


%\paragraph*{Tight reductions in the AI-GGM.}  Chen \emph{et al.}~\cite{C:CLMQ21} do not consider concrete security.  We will investigate whether concrete security of Schnorr given by their proof in the AI-GGM matches the bound for DL.  If not, we will devise a new proof that gives better concrete security.  


\iffalse
Let
$\advB$ be weak $f$-CDL solver. We construct a PGG adversary $\advA = (\advs,\advd)$ as follows.  Source $\advs$ chooses $x \getsr \Z_p$ and queries $x$ to get $X$. It returns $L = \calE^{-1}(x) \circ X$.  On input $(g,L)$, distinguisher $\advd$ runs $\advB$ on $g,L$ to get $R$. It returns 1 iff $R \circ g^{R} = L$. Note that for simple $\calE$, $\advs$ is a \emph{masking source} as defined in~\cite{TCC:BFHO22}.  For the analysis, we will need to estimate the density of elements in $\GG$ that have a CDL solution, and how many solutions they have. If the approach works, we aim to show PGG implies \emph{strong} $\calE$-CDL. 
\fi


%Unfortunately, it is not possible to evaluate (strong) CDL in the GGM because $\Z_p^*$ is a field.
%Nevertheless, we first propose to investigate whether weak and strong CDL are equivalent.  Second, morally speaking, we conjecture that PGG implies weak CDL or even strong CDL. The  problem is that $\Z_p^*$ is \emph{not} PGG, \emph{e.g.}, DDH does not hold. Thus, we will switch to the subgroup $\mathbb{QR}_p$ of quadratic residues, where $(p-1)/2$ prime.  



 \iffalse
\begin{ques}
    Can we extend (strong) CDL and above results to other groups, \emph{e.g.}, $\GG = \mathsf{EC}_p$?
\end{ques}
\fi

\iffalse
\paragraph*{Secure embeddings: Tight security in the general case.}
For a group $\GG$ of order $p$ generated by $g \in \GG$, fix an embedding  $\calE \colon \GG \to \Z_p$. Call an embedding secure if it is hard for an adversary to win the  following game: Given $h \getsr \GG$, output  $z\in \Z_p^*, R \in \GG$ such that \smash{$g^z = R \circ h^{\mathcal{E}(R)}$}; equivalently, output $a,b \in \Z_p$ such that $\mathcal{E}(g^a/h^b) = b$.  
%Call such $\mathsf{E}$ a \emph{secure embedding} for $\GG$. In fact, it suffices for our proof of unforgeability of Schnorr based on strong CDL that \emph{any} such embedding exists.  To analyze embeddings for $\mathsf{EC}_p$, it is attractive to work in the GGM. Actually, Chen \emph{et al.} observe that using a simple hash function $\mathsf{E}_{pos}$ that on input $(x,y)$ outputs whichever of the $x$-coordinate or $y$-coordinate is positive, Schnorr is  secure in the GGM but not in reality.  As a result, they advocate using GGM-with-preprocessing (AI-GGM)~\cite{C:Unruh07}.  An alternative we also point out is the recent EC-GGM model of Groth and Shoup~\cite{EC:GroSho22}.
In the PROM, we can tightly reduce security of Schnorr signatures in $\GG$ to security of the embedding in the above sense. 
In preliminary work, we show that a random bijection $\calE$ (where the entire truth table is given to the adversary) is a secure embedding in the GGM. We propose to investigate:
\fi

%In fact, we conjecture that for $\mathsf{EC}_p$, mapping a point on the curve $(x,y)$ to whichever coordinate is positive works.  Interestingly, Chen \emph{et al.}~\cite{C:CLMQ21} show using this as the hash function in Schnorr is \emph{not} secure. This underscores the point that requiring a secure embedding is much weaker than requiring a hash function that makes Schnorr unforgeable.  In particular, their attack requires a signing oracle.  

%In particular, the result of Chen \emph{et al.}~\cite{} uses the \emph{auxiliary-input} GGM, where the adversary gets an unbounded-time preprocessing phase before it plays the game, during which it can examine the entire truth-table for the group and produce a short ``advice'' string passed to the adversary in the online phase.  While this certainly captures in important class of attacks (dating back to Hellman tables~\cite{}), to us it seems like overkill for an initial study of Schnorr in the GGM.  So why do Chen \emph{et al.} use it? Interestingly, they point out attacks on Schnorr using simple hash functions that are not captured by the standard GGM, but are captured by the AI

%\paragraph*{Results in the programmable GGM.}
%We will first revisit Schnorr signatures in the (programmable) GGM. Indeed, Chen \emph{et al.} point out that for certain hash functions  there is an attack over $\mathsf{EC}_p$, but this attack is \emph{not} captured by the GGM.  We aim to prove a revised result, slightly changing the requirements on the hash function, in the recently-proposed EC-GGM model of Groth and Shoup~\cite{EC:GroSho22}. This more-realistic model capture the aforementioned attack by enforcing  some  simple structure on the group encoding. The next problem is that the results of Chen \emph{et al.} assume a sufficiently small message-space. To drop this assumption, we will make stronger assumptions on the hash function. Indeed, we observe that in the setting of Chen \emph{et al.} the \emph{identity} hash function suffices. To build intuition, we start with this special case.




\iffalse
\paragraph*{Non-programmable AI-GGM.} We will also  consider Schnorr in a new \emph{non-programmable} AI-GGM. In particular, we will investigate if analogous results hold there. 
\fi


\iffalse 
 \subsubsection{Results in the (Augmented) AGM} 
 As the next step, we would like to model the group with something in between the GGM  and standard model, which using a standard-model hash.  The AGM~\cite{C:FucKilLos18}, which was used to give a tight proof for Schnorr in the ROM~\cite{EC:FucPloSeu20},  seems like one possibility.  Unfortunately, in  the AGM,  Schnorr signatures using a standard-model hash \emph{do not  even compile}. In particular, in the AGM one is not allowed to map group elements to bits.  This begs the question: 
\begin{ques}
    Can we \emph{augment} the AGM in some way to allow hashing group elements to bits?
\end{ques}
Our initial idea is to augment the model  with a hashing operation  that is  ``one-way'' in the sense that  the bits cannot be converted back into group elements.  We will also require that   for any  hash value the adversary  outputs it must output an explanation of the input used to  compute it.  If successful, we will try to prove Schnorr secure in such a model.  We will see what assumptions on the hash function suffices compared to the standard-model case below. 

\fi
% Hashing will use a random key that is not input to any other part of the adversary's computation.

\subsubsection{Instantiability Results} \label{sec-inst}


Impossibility results~\cite[Tables 1 and 2]{FukumitsuH21} rule out certain reductions  in the \emph{non-programmable} ROM (where the reduction doesn't handle the adversary's RO queries) based on, \emph{e.g.}, OMDL; it appears they  also apply in the standard model.  Still, they don't cover \emph{non-blackbox} techniques. %But, to our knowledge, such results do not  preclude a proof of security for Schnorr in the \emph{standard model}. 
%This seems odd as the NPROM is ``stronger'' than the standard model, but the impossibility proofs are meta-reductions that simulate a particular forger for the reduction. As this forger \emph{makes calls to the RO}, it is unclear how to simulate a standard-model forger instead.
We first consider security against \emph{key-recovery} which, following~\cite{AC:PaiVer05}, we call BK-CMA. That is, in BK-CMA the adversary makes signing queries and wins only if it recovers the signing key.  %Paillier and Vergnaud showed Schnorr is BK-CMA for \emph{any} hash function $H$ if  OMDL holds for $\GG$. %Roughly, OMDL asks: given $n$ DL instances $X_1,\ldots,X_n \in \GG$ and up to $n-1$ calls to a DL oracle that on input $Y$ returns $\mathsf{DLog}_{\GG,g_0}(Y)$, outpu $\mathsf{DLog}_{\GG,g_0}(X_1), \ldots, \mathsf{DLog}_{\GG,g_0}(X_n)$.

In preliminary work, we show BK-CMA assuming $\GG$ is (plain) DL-hard and $H$ is an \emph{extremely lossy function} (ELF)~\cite{C:Zhandry16}.  ELFs can be built under exponential-DDH, which is a standard, non-interactive assumption as compared to OMDL. An ELF is a keyed function whose key $K$ (which we will put in Schnorr's verification key) can be generated in ``injective mode'' or ``lossy mode.''  In injective mode, $\mathsf{ELF}(K,\cdot)$ is injective whereas in lossy mode $|\mathsf{ELF}(K,\cdot)|$ is polynomial.  Furthermore, any distinguisher between the two  keys has advantage less  than any inverse-polynomial for large enough,  but polynomial, lossy image-size \emph{depending on the distinguisher's running-time}, which is  \emph{non-blackbox}.  We also use that the image of $\mathsf{ELF}(K,\cdot)$ in lossy mode can be efficiently enumerated.

We sketch the proof that $\mathsf{Sch}[\GG,\mathsf{ELF}]$ is BK-CMA assuming DL in $\GG$.  Let  $\advA$ be a BK-CMA adversary. We  switch $H$ to lossy mode.  Now consider the DL-adversary $\advB$ that on input $X$ runs $\advA$, setting $X$ to the verification key and answering signing queries as follows: On query $m$, $B$ does: $z \getsr \Z_p^*$, $c \getsr [\mathsf{ELF}(K,\cdot)]$ until $H(m || g^z/X^c) = c$, where $[\mathsf{ELF}(K,\cdot)]$ is the image of $\mathsf{ELF}(K,\cdot)$. It returns $(R,z)$ where $R = g^z/X^c$. At the end, $\advB$ returns $\advA$'s output. For simplicity, suppose $\mathsf{ELF}(K,\cdot)$ has equal-size preimage sets.   Then it is easy to see simulating a signing query takes a polynomial number of tries on average. Unfortunately, ELFs do not give a concrete-security bound. % Unfortunately, ELFs do not  give a concrete security bound. %Note that $\advA$'s runtime is long enough to break the ELF, which is not a contradiction. 

\begin{task} \label{task-inst}
    Extend the above instantiability result  from BK-CMA to UF-CMA of Schnorr signatures.
\end{task}

%Crites \emph{et al.}~\cite{EPRINT:CriKomMal21} turn existence of reduction from UF-CMA to BK-CMA for Schnorr into a knowledge assumption in its own right and prove it in the AGM.  Instead,

To do this, we  plan to  build an $H$ that instantiates Schnorr under UF-CMA by putting an ELF along with an  $f$ satisfying ($\GG$,$f$)-CDL inside an obfuscated program. In the proof, we want $H$ to ``hit'' $f$ during the forgery verification but the ELF during signing oracle simulation. This may be possible as the reduction has the flexibility to choose $R$  in the signing oracle simulation.  %Conceptually, this approach could be seen as strengthening the above-mentioned construction of CI hash function by using ELFs, similar to PI's prior work~\cite{AC:MurONeZah22}.
To start, we will consider a relaxation  of the problem to  \emph{composite-order} groups, as in a prior framework developed by the PI~\cite{AC:GLOW12} for proving security of many signature schemes. %To our knowledge, no one has studied Schnorr signatures in this setting. %We will also look for other techniques that do not use ELFs. %, as a cryptographic hash function is unlikely to be one (or even yield a construction of one)~\cite{TCC:FisRoh23}.  

%%While our techniques above depart from using CI hash functions, we will also try using constructions of CI hash functions directly and combining them with ELFs. 

\fi
  
  
%Fix $\GG,f$ as  in CDL. A straightfoward analogue of OMDL for CDL would be that the adversary gets $(X_1,\ldots,X_\ell)$ where $x_i \gets \Z_p$ for $i \in [\ell]$ and needs to output CDL solutions $(z_i,R_i)$ such that $g^{z_i} = R_i \cdot X_i^{f(R_i)}$ and $f(R_i) \neq 0$, for all $i \in [\ell]$. To this end, it can make at most $\ell-1$ queries to a CDL oracle that on input $(i,R)$ returns the DL of $R\cdot X_i^{f(R)}$.
%Note that the adversary may make many CDL oracle queries with the same $i$. Unf

%An example requirement is that the adversary must not be allowed to output a response from a CDL-query. d.
 %As such, we want to either find a modification of COMDL that tightly implies static security of $\mathsf{FROST}$, or  a modification to $\mathsf{FROST}$ for which we can get such a result. 

 
%We will study \emph{multisignatures with key aggregation},  starting with  Schnorr-based $\mathsf{HBMS}$~\cite{AC:BelDai21} and $\mathsf{MuSig2}$~\cite{C:NicRufSeu21}. Finally, we will treat \emph{blind signature

%%\begin{ques} \label{ques-thresh}
 %   How does the \emph{threshold} setting affect tightness and instantiability?
%\end{ques} 

%We will next consider the \emph{threshold} setting undergoing standardization~\cite{nistmpc,nist:schnorr}. %Our starting point will be enhanced security notions for   Schnorr allowing to modularly   ``thresholdize'' it  given by Shoup~\cite{EPRINT:Shoup23}.  Shoup proves them using the EC-GGM \emph{and} RO models. We'd like to drop the EC-GGM and give tight RO-model reductions  under CDL and  knowledge assumptions from Section~\ref{sec-know}.    %]
%We will start with  the Schnorr-based threshold signatures in~\cite[Section 4.2]{nist:schnorr}, as well as   $\mathsf{Sparkle+}$~\cite{C:CriKomMal23}, which is similar and  based on  DL. We will aim for tight RO-model security under CDL, as well as instantiations using CDL+ELFs. 
%We will similarly study Schnorr-based \emph{multi}-signatures, \emph{e.g.}~\cite{DBLP:journals/dcc/MaxwellPSW19,C:NicRufSeu21}.
%If this  fails, we may  use a high-moment forking lemma~\cite{C:RotSeg21,cryptoeprint:2024/934} instead. Adaptive security of $\mathsf{Sparkle+}$ requires the AGM. (Adaptive security of $\mathsf{Twinkle}$~\cite{DBLP:conf/eurocrypt/BachoLTWZ24} does not, but the scheme isn't similar to~\cite[Section 4.2]{nist:schnorr}.) A ``dream result'' wouuld be to prove an instantiability result here by combining the above techniques with  our  framework from Section~\ref{sec-know}.  %Adaptive security requires the AGM; to remove it we plan to incorporate our framework from Section~\ref{sec-know}. % We will also look for low-cost modifications to  the schemes that are beneficial for tightness and instantiability. 
 %Indeed, for optimizations or tight security, they again require the AGM + RO model. %  We will also consider the scheme of~\cite{AC:BonDriNev18} which, like Schnorr,  hashes a base to the exponent (here, to prevent rogue key attacks).

%\subsubsection{Tight Adaptive Security using DDH+CDL}

%x\subsubsection{Concurrent Security of  Anonymous Credentials Light}

\iffalse
Anonymous Credentials Light ($\mathsf{ACL}$)~\cite{CCS:BalLys13} is an easily-deployable anonymous credentials system that are basically blind-Schnorr signatures on user-committed attributes. Despite its popularity, a  worry in practice is \emph{concurrent} security, where an adversary may interleave many signing sessions. In fact, concrete attacks against $\mathsf{ACL}$ were proposed in this setting~\cite{EC:BenhamoudaEtAl21}.  However, more recent work~\cite{CCS:KasLosRen23} showed  issues with the attacks and  gave the first concurrent-security proof for $\mathsf{ACL}$ in the \emph{AGM + ROM}~\cite{CCS:KasLosRen23}. As usual we, would like to drop the AGM here. As the issues with giving a proof in the ROM alone are rewinding-based and CDL can eliminate rewinding, we ask:
\fi

%\begin{ques}
%    Is $\mathsf{ACL}$ concurrently-secure under CDL in the ROM?
%\end{ques}

%We stress that while a \emph{tight} CDL-based proof of concurrent security would be ideal, even a \emph{loose} ROM proof under CDL (\emph{e.g.}, via guessing a  query) would be a significant advance.
%As the schemes also differ in efficiency (round complexity, communication, and computation),   our results will help  practitioners  choose security vs.~efficiency tradeoffs for their applications.

\subsection{Tight Security of Dilithium}

NIST-selected $\mathsf{Dilithium}$~\cite{Dilithium} is the lattice analogue of Schnorr signatures. 
Its current tight ROM proof relies on three assumptions, including the artificial 
``Self-Target Module Short Integer Solution'' (ST-MSIS), which is merely unforgeability of 
$\mathsf{Dilithium}$ without signing queries. This makes ST-MSIS an overly strong assumption and a poor guide for cryptanalysis.

To address this, we  will introduce \emph{circular-MSIS}, a  novellattice analogue of CDL. Unlike ST-MSIS, 
it avoids adversarially-chosen inputs to $\hashH$ and treats $\hashH$ as a standard-model 
function rather than a random oracle. This will bring the key benefits of CDL to the lattice setting. 
Adapting our CDL-based techniques to the lattice setting and to the \emph{quantum} ROM (QROM) as in~\cite{Dilithium} --- 
and using standard MLWE and MSIS --- we will show:

\begin{conj} \label{conj-dil}
$\mathsf{Dilithium}$ is tightly secure in the QROM under MLWE, MSIS, and circular-MSIS.
\end{conj}

Since there is no established idealized model for lattices, we will justify circular-MSIS 
by constructing a specific conversion function and proving its validity under a standard 
lattice assumption, perhaps using ideas from a related work~\cite{C:CLMQ21}, and dovetailing on Section~\ref{sec-cdl-std}. 
If proving Conjecture~\ref{conj-dil} in full generality proves infeasible, we will first 
target relaxed notions such as selective security, or pivot to other post-quantum signatures 
such as $\mathsf{Falcon}$~\cite{falcon}.
%\noindent Our final proof will actually be in the \emph{quantum} ROM following~\cite{Dilithium}.
% If successful, we will also consider how to construct a standard-model conversion function that makes circular-MSIS hard. 
 
%A positive answer here will give more confidence in the security of $\mathsf{Dilithium}$, which is expected to  become widely used.

%The main  challenge  will be to adapt the algebra used in our tight proof of Schnorr signatures under CDL to the lattice setting. %If successful, we will look at candidate $f$ for circiular MSIS, inspired by~\cite{}.
 %If we are successful, we will also consider how to construct such a conversion function in the standard model here.


%\heading{Trading decision for search assumptions.}
%For \emph{lossy} $\Sigma$-protocols, Fiat-Shamir yields  signature schemes  with a \emph{tight} ROM  proofs. However, an important foundational issue is that  \emph{decisional assumptions} are then required, whereas search assumptions are used in the original proofs. We ask: \emph{Can we trade decisional for search assumptions while keeping reductions tight?}  To this end, we will use CDL or analogues to give tight ROM proofs for signature schemes obtained from  $\Sigma$-protocols like Okamoto~\cite{} and Guillou-Quisquater~\cite{}, \emph{without} going through lossiness. 

%In preliminary work, we have done this for Guillou-Quisquater~\cite{}, which is lossy under  $\Phi$-Hiding~\cite{}. Ideally, we will work at the level of a abstract $\Sigma$-protocol and define a CDL-style assumption and the properties needed for it.
%Note that it may be possible to show circular MSIS implies MSIS, in the same spirit that CDL implies DL.



%Finally, we will study tightness of  Fiat-Shamir for proof systems called \emph{succinct non-interactive arguments of knowledge} (SNARKs). In a SNARK for a relation $\calR$, a prover and verifier share an instance $x$ but only the prover knows the corresponding witness $w$ such that $x$ is in the language. Now, the prover wishes to convince the verifier with a  \emph{one-shot}  proof $\pi$ that $x$ is in the language, without revealing any information about $w$. The  goal is to minimize $|\pi|$ as a function of $|w|$.
\subsection{Tight Soundness of Bulletproofs-based Range Proofs} 

Broadening beyond signature schemes, we will study the tightness of related Fiat–Shamir-based proof systems known as
\emph{succinct non-interactive arguments} (SNARKs), where a prover 
produces a short proof $\pi$ that convinces a verifier an instance $x \in \calR$
without revealing the witness $w$. 

Our focus is $\mathsf{Bulletproofs}$~\cite{Bulletproofs}, a powerful generalization of Schnorr. It is built around an \emph{inner-product relation}, which enables  proofs of
correctness for more general computations. While $\mathsf{Bulletproofs}$ satisfies knowledge 
soundness under DL, the reduction is loose, giving only about 22-bits of  
security for typical parameters. As for Schnorr, tighter bounds are known in the algebraic group model~\cite{C:GhoTes21}, 
but we want to avoid such models. Therefore, we ask:

\begin{ques}
Does $\mathsf{Bulletproofs}$ admit tight \emph{knowledge} soundness in the ROM under CDL?
\end{ques}

Interestingly, CDL seems more suited to \emph{plain soundness} than knowledge soundness here. 
For soundness, one needs the relation to have  false statements, but the inner-product relation is \emph{trivial}: 
every instance has a witness. Thus, CDL gives no meaningful guarantee for $\mathsf{Bulletproofs}$ itself. 

Our solution will be to \emph{directly prove tight ROM soundness for protocols that apply 
$\mathsf{Bulletproofs}$ to non-trivial relations}, where false statements exist. We will use variants of CDL for this, which we will analyze as usual in the EC-GGM. Our main case of interest is 
$\mathsf{Bulletproofs}$-based \emph{range proofs}, as deployed by Monero in 2017~\cite{monero1}. The main challenge is that the protocol is recursive, and the prover obtains multiple 
Fiat–Shamir challenges per round. We will analyze the protocol round-by-round. If this proves too difficult, we can pivot to related applications of $\mathsf{Bulletproofs}$.
\iffalse
\heading{Kilian’s protocol.}
Chiesa \emph{et al.}~\cite{} showed improved upper-bounds on  soundness of Killian's protocol (succinct argument) as well as a lower-bound  showing that further improvement entails proving tighter security of Schnorr's $\Sigma$-protocol. 
Motivated by this connection, we will 
investigate whether  our techniques generalize to give better upper-bounds on soundness of Killian's protocol.
\fi
\iffalse
\subsection{Thrust Summary}

This thrust  will establish that Schnorr signatures and many of its extensions/analogues (threshold signatures, lattice-based signatures, SNARKs, \emph{etc.}) admit tight ROM proofs under new ``circular'' assumptions. In particular, we will give the first  proof (and even a tight one) for the influential threshold Schnorr signature scheme $\Sparkle$ under our new VCDL assumption, directly addressing NIST’s call. More generally, our planned results will address long-standing problems and inform NIST threshold and post-quantum signature standards, as well as blockchain deployments.
\fi

%\subsection{Lattice-Based Fiat-Shamir Signatures}
%Although the CDL problem may seem tailored to Schnorr, we believe our approach is actually quite general  and we will be  able to adapt it  to  signature schemes obtained via FS applied to  other ``special-sound'' $\Sigma$-protocols such as GQ~\cite{C:GuiQui88} and Okamoto~\cite{C:Okamoto92}. We will also focus on instantiability of FS when applied to the Chaum-Pedersen (CP) $\Sigma$-protocol~\cite{C:ChaPed92}, which is statistically-sound and does not require the forking lemma.  Applying FS to CP yields the verifiable random function (which can be viewed as a unique signature scheme) of~\cite{cryptoeprint:2017/099}, standardized by the IETF and used in iMessage and WhatsApp.  It is also deployed on the web in the form of Privacy Pass~\cite{DBLP:journals/popets/DavidsonGSTV18}.  Finally, we will  study instantiability of FS-based signature schemes  submitted to the  recent NIST PQC competition~\cite{pqc}, \emph{e.g.},  those based on MPC-in-the-head~\cite{STOC:IKOS07}.  
%Moving outside of signatures, we will extend our techniques to instantiate FS-based \emph{proof carrying data} schemes \emph{low-degree RO model .
%Finally, we will look at instantiability of alternative transforms to Fiat-Shamir~\cite{AC:BelPoeSte16} as  points of comparison.
 %They are similarly used to build an oblivious Privacy Pass protocol~\cite{DBLP:journals/popets/DavidsonGSTV18}. %Finally, besides signatures, we will also consider instantiability of FS for Bulletproofs~\cite{SP:BBBPWM18} and zkSNARKs. %such as $\mathsf{Plonk}$~\cite{EPRINT:GabWilCio19,cryptoeprint:2024/994}. 
 % barriers~\cite{STOC:GenWic11,TCC:BBHMR19} will guide our approach.
 %such as $\mathsf{Plonk}$~\cite{EPRINT:GabWilCio19,cryptoeprint:2024/994}, which are secure in the RO-model under falsifiable assumptions. (So, to instantiate them we will need non-falsifiable assumptions~\cite{STOC:GenWic11}; other related negative results include~\cite{TCC:BBHMR19}.) 
%One reason instantiations are especially  important for zkSNARKs is their use in \emph{recursive proofs}  (\emph{e.g.}, in  ZKrollups), which cannot be analyzed in the RO model in the first place. 
%Known barriers~\cite{STOC:GenWic11,TCC:BBHMR19} will guide our approach. %Intuitively, these schemes output ordinary  Schnorr signatures u nder  aggregated keys. %Note that Schnorr-based \emph{blind} signatures~\cite{EC:FucWol24} use ordinary Schnorr signatures as a black-box, so our analyses  will apply there. %associated to the group parameters could permit leveraging a standard-model notion of programmability analogous to the notion of Hofheinz and Kiltz~\cite{C:HofKil08} for hash functions.  
%For example, trapdoor groups can be used to instantiate the AGM~\cite{}; however, the proof for Schnorr in the AGM requires a  (programmable)  RO. 
%The second relaxation is using \emph{composite-order} groups. % One issue to here is that the domain of $f$ is $\GG$, not $\GG \times \bits^*$.

%To get a better handle on the problem, we will start with two relaxations.  The first is using a \emph{trapdoor} group or \emph{trapdoor} hash function, where the simulator knows some secret information associated with the group parameters or hash key (or both). %Having a trapdoor associated to the group parameters could permit leveraging a standard-model notion of programmability analogous to the notion of Hofheinz and Kiltz~\cite{C:HofKil08} for hash functions.  
%For example, trapdoor groups can be used to instantiate the AGM~\cite{}; however, the proof for Schnorr in the AGM requires a  (programmable)  RO. 
%The second relaxation is using \emph{composite-order} groups.  Indeed, The PI previously developed a framework of ``dual-form signatures''~\cite{AC:GLOW12} for proving signature schemes secure in composite-order groups when their proofs in prime-order groups require the GGM.

%Impossibility results in the NPROM often do not carry over to the standard model.  In the present context, the reason is that in a the forger simulated by the ``meta-reduction'' makes calls to the RO.  To show an impossibility result in the standard model, it would need to simulate a standard-model forger.  
%We note that Pass~\cite{STOC:Pass11} shows that there is does not exist a black-box proof for Schnorr's identification scheme under \emph{active} attacks, however, such attacks is not relevant here.

%\subsection{Signature  Schemes with Advanced Functionalities}

%Schnorr signatures are the basis for numerous signature schemes with advanced functionalities like threshold signatures~\cite{}, ring signatures~\cite{} and multisignatures~\cite{}. Such threshold signatures are currently in the process of being standardized by NIST~\cite{}.
%We will aim to extend both our  tight-security and instantiability results for Schnorr signatures to such  advanced primitives as well.

%Building on our proposed work in the NPROM, we will consider relaxations to trapdoor groups or hash functions, as well as to composite-order groups, While it does not seem possible to use the AGM as it does not allow mapping group elements to bits,  the right knowledge assumptions could ``simuluate'' it for us while while still allowing such mappings.  While our assumptions will probably be new and untested, we do aim to give the first positive results about Schnorr in the standard model more than \emph{30 years} after it was introduced. 
 
\iffalse
\paragraph*{Generalizing to other $\Sigma$-protocols.} There are numerous other $\Sigma$-protocols to which FS is commonly applied, such as...  We propose to generalize our results to signature schemes obtained from these $\Sigma$-protocols along the following lines. Given such a $\Sigma$-protocol,  fix a mapping from the commitment-space to the challenge-space, and consider an empty message-space and no signing queries. Then ``forging'' corresponds to an analogue of strong CDL in this setting, which we can use to study instantiability of the corresponding signature scheme.  We expect that extensions of PGG to the respective settings will be critical in justifying  the new assumptions. 
\fi


\iffalse
\subsection{Instantiability of Alternative Transforms}



\paragraph*{Fischlin's transform.} 
 Fischlin's RO-model transform~\cite{C:Fischlin05} is an alternative to FS that offers a \emph{tight} security reduction an underlying `Fiat-Shamir' proof of knowledge. Compared to transforms of 
Bellare \emph{et al.}~\cite{AC:BelPoeSte16} that also have tight reductions, Fischlin's only requires standard special-soundness of the underlying proof system and no trapdoor (so applies to Schnorr).
 
 Performance analysis of  Fischlin-based Okamoto signatures~\cite{AFRICACRYPT:DagVen14} shows they offer faster signing and verification times than FS-based ones. There are also additional benefits, such as that the verifier can start  verification before the entire signature is received, and the scheme achieves leakage-resilience with better leakage. Note that leakage resilience is why they  focus on Okamoto; however, similar performance gains should also apply in other cases.

Due to the above benefits, we are interested in \emph{instantiability} of Fichlin's transform. Specifically: 
\begin{ques}
   How does instantiability of Fischlin's transform compare to   to Fiat-Shamir?
\end{ques}

Unlike FS, Ananth \emph{et al.}~\cite{TCC:ABGR13} show that Fischlin's transform can yield an unsound non-interactive argument of knowledge (much less a signature scheme) when applied to suitable `Fiat-Shamir' proof of knowledge. In this sense, instantiability of Fischlin's transform seems  \emph{worse} than FS. 

\paragraph*{Fischlin+ transform.} Upon closer examination, we find the counter-example of~\cite{TCC:ABGR13} relies on a lack of \emph{domain separation} between the underlying proof system and the transform. This can be  fixed by appending a super-logarithmic length random string to the parameters of the transform and using it as a prefix to all hash queries. 
Call the resulting transform Fischlin+.

\begin{ques}
    Is the Fischlin+ transform instantiable?
\end{ques}

% While instantiability of the RO used by FS to obtain NIZKs is now well-understood, the same is not true of Fischlin's transform.  Hence, we will start by studying  this case, which could have significant implications on its own.  We are particularly interested in the following questions:
%\begin{enumerate}
  %  \item What assumptions are sufficient to instantiate the RO in Fischlin's transform for NIZKs?
   % \item To instantiate Fischiln's transform for NIZKs, is a form of CI necessary?
 %   \item Does instantiating Fischlin's transform for NIZKs require  milder assumptions than FS?
%\end{enumerate}

Fischiln gives a thoughtful explanation  of how the RO is used in his transform.
In particular, he uses a RO for a Nakamoto-style proof-of-work, as well as for ensuring that hashing  distinct inputs results in  independent outputs.  Interestingly,  Canetti \emph{et al.}~\cite{EC:CCRR18} observe such a  proof-of-work can be instantiated based on an achievable form of  CI hashing.
We aim to show instantiability of Fischlin+ on \emph{any} starting `Fiat-Shamir' $\Sigma$-protocol.
\fi 



\section{ROM Instantiability via Extremely Lossy Functions} \label{sec-pgg}
In our second thrust, we will instantiate Schnorr signatures~\cite{C:Schnorr89} and 
the Fujisaki–Okamoto hybrid-encryption transform~\cite{JC:FujOka13}. The connection is that our approach for both 
builds on new techniques for \emph{extremely lossy functions} (ELFs)~\cite{C:Zhandry16} 
and related tools. Whereas prior work treated ELFs largely as a theoretical construct, 
our techniques will make them directly relevant to standardized and widely deployed schemes. We aim not only to prove instantiability, but to clarify which hash-function and PKE properties deployed standards implicitly rely on, ensuring their security rests on transparent, scientifically-vetted criteria. %Formalizing such criteria 
%provides NIST and other standards bodies with concrete targets for validation.
  %ELFs  can  realized from the ``exponential'' Decision Diffie-Hellman (DDH) assumption (DDH against exponential-time adversaries), a standard assumption on elliptic curves. In the proposed work,  we will investigate sufficient relaxations that could be realizable from cryptographic hashing. %Thile not a true falsifiable assumption in the sense of Naor~\cite{C:Naor03} (because the adversary runs in exponential time), exponential DDH fits into a related framework of Goldwasser and Kalai~\cite{}. 
%For%Schnorr signatures,  we will show the first instantiability results. Then, building on our work at ASIACRYPT 2022, we will show instantiations of the celebrated  Fujisaki-Okomoto ($\mathsf{FO}$) encryption transform~\cite{} that are \emph{simple and post-quantum}. %A connection between Schnorr signatures and ROM encryption transforms here is that we will develop new instantiations for both primitives using ELFs. 

%More generally, we believe our techniques surrounding ELFs we will also be useful to show instantiability of extensions to Schnorr signatures such as threshold Schnorr and to the lattice-based analogoue $\mathsf{Dilithium}$, considered  in our previous thrust. Moreover, for the part of the proposed work on instantiating encryption transforms, we will develop new notions of ``adaptivity'' concerning  cryptographic primitives that consider auxiliary input, which will also be of independent interest. %We pursue such results as part of the proposed research.
%Accordingly, one of  focus on their post-quantum instantiability. The connection of Schnorr signatures is that they can also be instantiated using related techniques based around ELFs. %Post-quantum instantiability is not straightforward for us because the exponential DDH assumption is pre-quantum.

%When we implement and deploy a scheme designed and analyzed in an idealized model, of course we do not have an ``idealized object'' such as a random oracle accessible to all parties.  The idea is to replace this idealized object with a concrete one, such as $\mathsf{SHA256}$, whose code the all parties know and can locally execute. However, we do not have any rigorous evidence that this   ``instantiated'' scheme using $\mathsf{SHA256}$ will actually be secure. To remedy this, we will look at new approaches to instantiation result. Our main cases studies here include Schnorr signatures, the Optimal Asymmetric Padding (OAEP) and Fujisaki-Okamoto encryption transforms, all of which are set in the ROM. To broaden our study we then look at  ``cross-cutting'' approaches to instantiability that apply to many idealized models and give some applications.  %We will then study a unified paradigm for instantiating popular ROM-based encryption transforms like OAEP and Fujisaki-Okamoto, which are now relevant in the post-quantum setting. Finally, we will study ``cross-cutting'' approaches to instantiability that apply to many idealized models and give some applications.


\subsection{Instantiating Schnorr Signatures via ELFs} 

 For the knowledgeable reader, we note that Schnorr signatures cannot be instantiated using (achievable versions of) \emph{correlation-intractable} (CI) hashing~\cite{JACM:CanGolHal04} for analogous reasons to why it cannot be used for CDL.
Indeed, all positive results  are in the  ROM (+ AGM)~\cite{JC:PoiSte00,EC:AABN02,EC:FucPloSeu20,INDOCRYPT:BelDai20,C:RotSeg21} or the GGM~\cite{JMC:NevSmaWar09,C:CLMQ21,EC:BloLee22,EPRINT:Shoup23},  except  that the scheme is secure wrt.~key-recovery under chosen-message attack  for \emph{any} hash function assuming OMDL (one-more discrete-log)~\cite{AC:PaiVer05}, an interactive assumption where the adversary makes DL queries. 
%Moreover, impossibility results~\cite[Tables 1 and 2]{FukumitsuH21} rule out  proofs  in the \emph{non-programmable} ROM but don't cover \emph{non-blackbox} techniques, which we will employ via  ELFs. %But, to our knowledge, such results do not  preclude a proof of security for Schnorr in the \emph{standard model}. 
%This seems odd as the NPROM is ``stronger'' than the standard model, but the impossibility proofs are meta-reductions that simulate a particular forger for the reduction. As this forger \emph{makes calls to the RO}, it is unclear how to simulate a standard-model forger instead.

\subsubsection{Instantiation    against Key-Recovery under Chosen-Message Attack} \label{sec-kr}

We first instantiate Schnorr signatures under \emph{key-recovery} against chosen-message attack which, following~\cite{AC:PaiVer05}, we call BK-CMA. In  BK-CMA,  the adversary can make signing queries  as usual but has to  recover the signing key to win. We will show: %Paillier and Verngard~\cite{AC:PaiVer05} showed an instantiation of Schnorr signatures under this notion assuming OMDL, which is interactive and non-falsifiable. (OMDL is assumed on the underlying group; the hash function can be anything.) 
%o improve on the result %of~\cite{AC:PaiVer05}, we ask:
%\begin{ques}
%    Can we instantiate Schnorr signatures against key-recovery under chosen-message attack using a  \emph{non-interactive}  and \emph{falsifiable}?
%\end{ques}
%Paillier and Vergnaud showed Schnorr is BK-CMA for \emph{any} hash function $H$ if  OMDL holds for $\GG$. %Roughly, OMDL asks: given $n$ DL instances $X_1,\ldots,X_n \in \GG$ and up to $n-1$ calls to a DL oracle that on input $Y$ returns $\mathsf{DLog}_{\GG,g_0}(Y)$, outpu $\mathsf{DLog}_{\GG,g_0}(X_1), \ldots, \mathsf{DLog}_{\GG,g_0}(X_n)$.

  
  \begin{conj}
      Let $\hashH := \ELF$ and suppose DL holds in $\GG$. Then Schnorr signatures $\mathsf{Sch}[\GG,\hashH]$ are BK-CMA in the standard model.
      \end{conj}
  %We later show how to ``balance'' the assumptions here. %We  explain the idea of our prelimiminary result here.
 An ELF is a keyed function whose key $K$ (which we put in Schnorr's public key) can be generated in ``injective mode'' or ``lossy mode.''  In injective mode $\mathsf{ELF}(K,\cdot)$ is injective whereas in lossy mode, which unlike in standard lossiness~\cite{JACM:PeiWat11} \emph{depends on the adversary}, allowing  $|\mathsf{ELF}(K,\cdot)|$  to be polynomial-size  %Namely, for any adversary $A$ and polynomial $p$ we can generate keys in a way that $A$  can distinguish the modes with advantage at most $1/p$ and $|\mathsf{ELF}(K,\cdot)|$ is polynomial when $K$ is generated in lossy mode. %In particular, for any inverse polynomial $1/p$, any distinguisher between the two  key generation modes has advantage less than $1/p$ for a large enough  but polynomial-size lossy image-size. 
 and the  image efficiently enumerable. ELFs can be constructed under  exponential DDH~\cite{C:Zhandry16}, which, while not falsifiable, is standard for elliptic curves and preferable to OMDL. Note that non-falsifiable assumptions can still be  invalidated~\cite{C:BelPal04}. 
%We sketch the proof that $\mathsf{Sch}[\GG,\mathsf{ELF}]$ is BK-CMA assuming DL in $\GG$.  
Our key idea is that:
\begin{quote}
\emph{ELFs allow answering signing queries in a standard-model proof of unforgeability for Schnorr signatures by simulating ``programmability'' used in the classic ROM proof.} 
\end{quote}
To understand this, we sketch our proof idea. Let  $\advA$ be a BK-CMA adversary. We  first switch $\hashH$ to lossy mode.  Consider the DL-adversary $\advB$ that on input $X$ runs $\advA$ on the same input $X$  and answers signing queries as follows: On query $m$, $B$ does: $z \getsr \Z_p^*$, $c \getsr [\mathsf{ELF}(K,\cdot)]$ until $\hashH(m || g^z/X^c) = c$, where $[\mathsf{ELF}(K,\cdot)]$ is the image of $\mathsf{ELF}(K,\cdot)$. (In the ROM, we would \emph{program} $\hashH(m || g^z/X^c)$ to be $c$.) It returns $(R,z)$ where $R = g^z/X^c$. Eventually, $\advB$ returns $\advA$'s output.  %For simplicity, suppose $\mathsf{ELF}(K,\cdot)$ has equal-size preimage sets.  
%This loop  executes a polynomial number of times on average. %Unfortunately, ELFs do not give a concrete-security bound. % Unfortunately, ELFs do not  give a concrete security bound.  %We will conduct a detailed analysis  to the proposed work.

    ELFs can be replaced by \emph{tunable lossy functions} (TLFs) defined in Section~\ref{sec-tools}, lowering the assumptions down to sub-exponential DDH. With respect to cryptographic hashing, however, neither ELFs nor TLFs can be constructed in the ROM~\cite{TCC:FisRoh23}, and real hash functions are anyway not injective. Thus, ELFs/TLFs do not yield a useful design criterion for them. To address this, we will explore whether further relaxations are sufficient, such as ``targeted lossiness''~\cite{C:QuaWatWic21}.
 
\iffalse
\heading{Balancing the  assumptions.} Zhandry's ELF construction~\cite{Zhandry16} relies on exponential DDH, and, since ELFs are ``complexity absorbing,'' in our  proof we also assume  polynomial hardness of DL in the underlying group. To ``balance'' hardness of DDH and DL we can instead rely on sub-exponential DDH. In that case, we can adapt Zhandry's ELF to have sub-exponential size. Moreover, as the DL adversary will need to enumerate the lossy ELF image, sub-exponential DL should hold, which is indeed implied by sub-exponential DDH.
\fi
%\heading{Relaxations of ELFs.} ELFs have the curious property that a random oracle (or even any hash function build from one) is \emph{not} an ELF. In other words, we are instantiating a RO using a property of a hash function that is not satisfied by a RO. We will therefore investigate novel relaxations of ELFs that \emph{can} be achieved in the ROM and suffice for our applications.

\subsubsection{Instantiation against Unforgeability under Chosen-Message Attack}



%\begin{ques}
%    Can we instantiate Schnorr signatures against \emph{unforgeability} under chosen-message attack (UF-CMA) using non-interactive assumptions?
%\end{ques}
We  will  next substantially strengthen  the above  result 
by leveraging our  \emph{circular discrete-logarithm} (CDL) assumption from Thrust 1. We assume CDL in the underlying group. Let $f$ be a good conversion function (cf.~Section~\ref{sec-cdl-std}).  We will show:

\begin{conj}
        Let $\hashH := \mathsf{iO} (f \circ \PRF)$ and suppose $\CDL[\GG,f]$ holds. Then Schnorr signatures $\mathsf{Sch}[\GG,\hashH]$ are UF-CMA in the standard model.
\end{conj}
Above, $\iO$ denotes an indistinguishability obfuscator and $\PRF$ a puncturable pseudorandom function. Here, we will use ELFs, rather than $f$ as in the previous thrust, only within the proof. Our rough idea for the proof is to undetectably switch $\hashH$ so that it has two execution branches: an $\ELF$ branch and an $f$ branch. In the reduction, the $\ELF$ branch is used to answer the adversary's signing queries, while the $f$ branch enables the reduction to ``profit'' from the adversary's forgery.
%The idea is that $\ELF$ will allow us to handle the signing queries, while $f$ plays the orthogonal role of allowing us to ``profit'' from the adversary's forgery. Roughly speaking, we instantiate $\hashH$ using an  \emph{obfuscation} of $f$, \emph{i.e.}~we use indistinguishability obfuscation (iO)~\cite{C:BGIRSVY01,SIAM:BGIRSVH16} here. Furthermore, We will incorporate $\ELF$ into $\hashH$ but only in the proof.
%For simplicity, we   ignore the puncturable PRF here. Consider a forger $\advA$ against Schnorr signatures instantiated as above. We will first change $\hashH$ to have an  ``trapdoor branch'' on which it evaluates $\ELF$ and not $f$.  %Intuitively, we will use the trapdoor branch to answer signing queries in the subsequent reduction to CDL. 
%We will then change $\ELF$ to lossy. Now, we sketch a reduction to CDL. The CDL adversary $\advB$ runs the $\advA$ on its input, and answers signing queries similarly to the proof of BK-CMA,    sampling inputs that trigger the trapdoor branch. Finally, as long as  $\advA$'s  forgery  $(m^*,(R^*,z^*))$ does not trigger the trapdoor branch, $\advB$  wins the CDL game if $\advA$ forges. 
%This is a rough idea that will need significant work. 
%In particular, we will need tricks to make the above work with iO, and there are substantial additional issues in the proof to resolve. For one,  $f$ in CDL has domain  $\GG$, not $\GG \times \bits^*$ as above. We will also need to  ensure  outputs of $f$ on random inputs look  indistinguishable from those of $\ELF$ on random inputs. Finally, in the reduction to CDL, we will  need to sample inputs to $\hashH$ of the form $m \| g^z/X^c$ that trigger the trapdoor branch.

While this  establishes instantiability, iO is not a notion  cryptographic hashing can achieve. If it proves too difficult to  get an instantiation   based on properties achievable by cryptographic hashing, we will at least  \emph{eliminate the use of iO}, lending evidence that a ``worldly'' hash function suffices.   %Finally, an element $R$ received by the forger as a result of a signing query can be reused in its forgery.  We hope these issues can be resolved by using forms of puncturable PRFs, \emph{etc.} 
%To build up our ideas, we will start with relaxations of unforgeability incomparable to BK-CMA  such as a bounded number of signing queries and the forger declaring its  signing queries up front.  To the best of our knowledge,  no  instantiability result is known for Schnorr signatures in these  settings either.

%\heading{Extending to advanced functionalities.} As in the case of tightness, it will be interesting to look at instantiability of advanced schemes based on Schnorr like the $\mathsf{Bulletproofs}$ SNARK. We will see if we can extend the above ideas to such settings. 
%Finally, we hope to connect the approach with our new paradigm for instantiating ROM encryption transforms previously described.  There, we want to obfuscate a program which has a trapdoor branch that executes on a single point (the point hidden by the MB-AIPO) and outputs the output point of the MB-AIPO. Here we want the trapdoor branch to execute on \emph{polynomially-many} input points given by the the responses to the adversary's signing queries. 

\iffalse
We plan to build $H$ using a combination of  ELFs with other primitives that are  non-interactive and, if not falsifiable themselves, implied by falsifiable assumptions. Notably, borrowing from our results on tight security of Schnorr signatures, we plan to use \emph{circular Diffie-Hellman} (CDL) in $\GG$. In particular, let $f$ be such that $\CDL[\GG,f]$ holds. We also need a distribution $\calD$ on $\GG$ such that: given a $K$ key, it is possible to both sample an element from $\calD$ along with its discrete-logarithm and test whether an element is in the support of $\calD$. Furthermore, without $K$, a sample from $\calD$ looks like a random element, and it's hard to generate any element in the support of $\calD$ Then, our rough idea (which does not  work) is to obfuscate the program:
 \begin{center}
 \begin{tabular}{c} 
 \begin{minipage}{1in}
\begin{tabbing}
123\=123\=123\=123\=123\=\kill
\textbf{Program} $\mathsf{P}[K, K_\mathsf{elf}](R \| m)$: \\
\> If $R \cdot X^c \in \mathsf{Supp}(\calD)$ then return $\mathsf{ELF}_{K_\mathsf{elf}}(R \| m)$ \\
\> Else return $f(R\|m)$
\end{tabbing} \end{minipage}
\end{tabular}
\end{center}

The idea for the proof would initially be to handle signing queries as in our BK-CMA proof, but sample group elements $g^z \in \calD$. The adversary can only generate signing queries $(R,m)$ with $R \notin \calD$ unless it reuses $R$ it received from a signing query.  
\fi



% We plan to extend the result using two key ideas: \textbf{Idea 1:} Combine the ELF with an $f$ such that $\CDL[\GG,f]$ holds. \textbf{Idea 2:}  Wrap the combination in an obfuscated program using indistinguishability obfuscation.

%We stress that this will give a \emph{theoretical} construction of a hash function that instantiates Schnorr signatures that would not be used in practice. The  existence of such a hash function makes the goal of instantiating Schnorr signatures with a cryptographic hash function  plausible.

%\subsubsection{Extensions to Other Fiat-Shamir-based Protocols}
%Crites \emph{et al.}~\cite{EPRINT:CriKomMal21} turn existence of reduction from UF-CMA to BK-CMA for Schnorr into a knowledge assumption in its own right and prove it in the AGM.  Instead,

%Our initial idea for how to do this is to wrap the 
%To do this, we  plan to  build an $H$ that instantiates Schnorr under UF-CMA by putting an ELF along with an  $f$ satisfying ($\GG$,$f$)-CDL inside an obfuscated program. In the proof, we want $H$ to ``hit'' $f$ during the forgery verification but the ELF during signing oracle simulation. This may be possible as the reduction has the flexibility to choose $R$  in the signing oracle simulation.  %Conceptually, this approach could be seen as strengthening the above-mentioned construction of CI hash function by using ELFs, similar to PI's prior work~\cite{AC:MurONeZah22}.
%To start, we will consider a relaxation  of the problem to  \emph{composite-order} groups, as in a prior framework developed by the PI~\cite{AC:GLOW12} for proving security of many signature schemes. 

 \subsection{Instantiating ROM Encryption Transforms via ELFs}
 

%A particularly vexing case of ROM instantiability concerns \emph{transforms}; in other words,  compilers that take one or more ``base schemes'' (that may or may not use ROs) and output a ``target scheme'' that uses ROs.  
%We say that the transform is secure if for any secure base schemes the output target scheme is secure under the appropriate security notions.  
%xSo we refer to a ROM transform as  \emph{instantiable} if there are standard-model hash functions that can  replace the ROs used by the target scheme such that the latter is secure when the base schemes are. %This means the transform cannot ``work'' in the standard model in general.
%Examples of RO model transforms include~\cite{}, and~\cite{}; have been shown uninstantiable in this sense~\cite{}. 

\iffalse
\heading{Approaches toward instantiability.}
Instantiating transforms in the ROM is difficult because they must resist contrived base schemes. As such, one approach is to identify practical \emph{sub-classes} of base schemes for which we can  instantiate a transform, meaning the transform works as long as the base schemes live in these sub-classes. Such sub-classes are often characterized by security properties, \emph{i.e.} security properties beyond those needed in the ROM. Indeed,  we may not expect that a scheme can support the same class of base schemes once ROs are dropped. We may also  have to evade impossibility results.  In this paper,  we provide such results for two  popular transforms.
\fi

% This raises the questions:
%\begin{quote}
%For transforms that are not known to be uninstantiable, can we instantiate them?
%For transforms that are known to be uninstantiable, can we propose alternative transforms that are instantiable?  Or can we prove instantiability restricting to ``practically relevant'' classes of base schemes?
%\end{quote}
%Here practically relevant would be defined in terms of properties that schemes the transform is likely to be applied to in practice would satisfy. 
Now, we will study instantiability of highly influential  ROM encryption transforms that upgrade weak to strong encryption, above all focusing on the Fujisaki–Okamoto ($\mathsf{FO}$) transform that is ubiquitous in NIST PQC schemes.
  %Combining new techniques for ELFs with other novel primitives, we will show  instantiability of two  ROM ``encryption transforms'' that, like Schnorr signatures, are ``crown jewels'' of the ROM. Specifically, we will consider the Optimal Asymmetric Encryption Padding ($\mathsf{OAEP}$) trapdoor permutation based transform~\cite{EC:BelRog94} and the Fujisaki-Okamoto ($\mathsf{FO}$) hybrid-encryption transform~\cite{JC:FujOka13}.  
 %At a high level, these transforms upgrade a ``weak'' public-key encryption scheme meeting a notion like one-wayness (OW-CPA) to a ``strong'' one meeting security under chosen-ciphertext attack (IND-CCA). IND-CCA is the gold-standard notion of secuirty for encryption where the adversary can make queries for decryptions of ciphertexts of its choice other than its challenge.
%\emph{uninstantiability} results to some extent. % for FO~\cite{TCC:BrzFarMit15} and a black-box uninstantiability result for OAEP~\cite{EC:KilPie09}.
%There was a flurry of work on instantiability of these schemes following the general uninstantiability results~\cite{}. but the thorny issue was never resolved\
%Accordingly, the main question we study is:
%\
%\begin{quote}
%    \textit{Do there exist standard-model hash functions that suffice to instantiate IND-CCA2 secure OAEP and FO?} %(In the case of OAEP, we must use non-blackbox techniques.  %In the case of FO, we must  consider additional properties  of the base PKE scheme compared to the RO model results.)
%\end{quote}
%For simplicity, we only discuss $\mathsf{FO}$, which is pervasive today in  post-quantum cryptography. 
\iffalse
$\mathsf{OAEP}$ takes a trapdoor permutation (TDP) $\calF$ and produces a public-key encryption scheme whose public key is an instance $f$ of the TDP.  It uses two ROs $\hashG,\hashH$ and the encryption algorithm has the form: 
\[
\calE^{\mathsf{OAEP}}_{f}(m;r) = f(s \| t) \quad \mbox{where} \quad s = \hashG(r) \xor  m \| 0^\zeta \quad \mbox{and} \quad t = \hashH(s) \xor r 
\enspace,
\]
where $\zeta \in \N$ is the redundancy parameter. Decryption inverts the TDP and the padding transform to recover recovers $s \| t$; it parses $s = m \| z$ and returns $m$ if $z = 0^\zeta$, otherwise it returns $\bot$.
%In the case of the full-domain hash signatures~\cite{CCS:BelRog93}, black-box uninstantiability was bypassed by Hohenberger \emph{et al.}~\cite{EC:HohSahWat14} using program obfuscation, which partly inspires our work.
\fi
$\mathsf{FO}$   takes a public-key encryption   scheme $\PKE$ (or key-encapsulation mechanism) and a symmetric-key encryption scheme $\mathsf{SE}$ to produce a hybrid public-key encryption scheme  $\mathsf{FO}[\PKE,\mathsf{SE}]$ secure against chosen-ciphertext attack. Encryption has the form:
\[
\calE^{\hyb}_{pk}(m;r) = 
\calE^{\asy}_{pk}(r; \hashH(r)) \| \calE^{\sy}_{K}(m) 
\quad \mbox{where} \quad
K = \hashG(r \| c_1), c_1 = \calE^{\asy}_{pk}(r; \hashH(r))
\enspace,
\]
%In the initial paper, $\hashH$ takes as additional input $c_2 = \calE^{\sy}_{K}(m)$, which allows the authors to prove the FO transform works in the RO model for public-key schemes that are one-way CPA and \textit{well-spread}, and symmetric-key encryption schemes that are IND-CPA. 
%Later Hofheinz, H\"{o}velmanns, and Kiltz~\cite{TCC:HofHovKil17} dropped the additional input to $\hashH$ and traded the well-spread condition on 
where $\calE^{\asy}_\pk$ is the encryption algorithm of $\PKE$ and $\calE^{\sy}_K$ is the encryption algorithm of $\mathsf{SE}$. Decryption of $c_1 \| c_2$ first decrypts $c_1$ to recover $r$ and returns $\bot$ if  $c_1 \neq \calE^{\asy}_{pk}(r; \hashH(r))$. Else it returns decryption of $c_2$ under $K = \hashG(r \| c_1)$. We will also analyze   variants that use   ``implicit rejection.''

We refer to a ROM transform as  \emph{instantiable} if there are real hash functions that can  replace the random oracles in the output scheme, such that the latter is secure whenever the base schemes are, under the appropriate security  notions. % We focus on instantiability of $\mathsf{FO}$ in this sense. %For $\mathsf{FO}$, a \emph{post-quantum} instantiation is clearly essential.  Hence, we ask:  %If the base PKE scheme is ``rigid'' as defined by~\cite{}, the re-encryption check upon decryption can be eliminated.

%Other modifications have been proposed in the \emph{quantum} RO model (see Section~\ref{sec-related}) particularly because of the NIST PQC competition~\cite{nist}.
%It is thus of growing importance to establish instantiability of FO.
%This method seems preferable since one-time IND-CCA symmetric encryption is cheap, and so we build on this transform. 

%More generally  they showed uninstantiability of all ``admissible'' encryption transforms. 
%FO has gained renewed interest recently in the context of  post-quantum cryptography (PQC), as some variant of it is used by most of the KEMs submitted to the NIST PQC competition~\cite{nist}.
%%Ideally one would  We endeavor to retain instantiability for any IND-CCA PKE scheme, so we seek to slightly modify the transform.
% For OAEP, Shoup~\cite{C:Shoup01} showed that it cannot be secure for \emph{every} one-way trapdoor permutation.  But this result does not apply to practical TDPs, let alone the commonly used RSA TDP.  Further, there are black-box impossibility results imply that one either has to use non-blackbox assumptions on the hash functions or on the TDP~\cite{}.   Accordingly, we target full instantiation of the OAEP using non-blackbox assumptions.
% Further results about the schemes are discussed below.

%The main question we study is:


%here are negative results  we will need to circumvent. % Kiltz and Pietrzak~\cite{EC:KilPie09} show $\mathsf{OAEP}$ does not have a black-box proof of IND-CCA security under any assumption implied by a random TDP.  
%Brzuska \emph{et al.}~\cite{TCC:BrzFarMit15} show that for $\mathsf{FO}$ there exists an IND-CPA secure base PKE scheme on which the scheme output by $\mathsf{FO}$ is insecure. %Looking ahead, we will make stronger or orthogonal assumptions on the base PKE scheme as compared at IND-CPA.

\subsubsection{Our Prior Results}

%The PI's prior results at ASIACRYPT 2022~\cite{EC:MurOneZah22} show instantiations of $\mathsf{FO}$ and $\mathsf{OAEP}$. The results have various limitations we later propose to overcome. 
%Here we overview the results of our prior work, which forms the basis for our proposed work.
%Namely, Tte focus of the proposed work will be on extending our treatment of $\mathsf{OAEP}$ and $\mathsf{FO}$ below to the \emph{post-quantum} setting, which is of rapidly growing importance. 

At ASIACRYPT 2022~\cite{AC:MurONeZah22}, we proposed a unified paradigm for instantiating a variant of  $\mathsf{FO}$, and the $\mathsf{OAEP}$~\cite{EC:BelRog94} transform, using the same hash function built from indistinguishability obfuscation, ELFs, and puncturable PRFs:  %(We are mainly concerned with instantiating one of the two random oracles used in each of these transforms; the other one is then easier to instantiate.)  % The unified paradigm is reminiscent of the technique proposed for instantiating Schnorr signatures above, although we do not yet know a formal connection. %Our hash function  employs indistinguishability obfuscation (iO)~\cite{C:BGIRSVY01,SIAM:BGIRSVH16} to obfuscate the composition of a punctured pseudorandom function (PPRF)~\cite{AC:BonWat13,EC:BoyGilIsh15,CCS:KPTZ13} with an extremely lossy function (ELF)~\cite{C:Zhandry16}. 
%extremely lossy function (ELF) using indistinguishability obfuscation~\cite{C:BGIRSVY01,SIAM:BGIRSVH16} (whose recent constructions~\cite{STOC:JainLS21,EC:WeeW21,STOC:GayP21} make feasibility results using iO compelling), (2) an  extension of an idea due to
%Sahai and Waters~\cite{STOC:SahWat14} and 
%Black-box uninstantiability of the full-domain hash signatures~\cite{CCS:BelRog93} was bypassed by Hohenberger \emph{et al.}~\cite{EC:HohSahWat14} using program obfuscation, which partly inspires our work.
%In the security proofs, we extend an idea  of~\cite{AC:BrzMit14a} and use multi-bit auxiliary-input point function obfuscation (MB-AIPO)~\cite{EC:CanDak08} in the obfuscated program as well. % MB-AIPO obfuscates  a program $P_{x,y}$ which on input $x'$ outputs $y$ iff $x=x'$; in our case $x$ will be the randomness for the challenge ciphertext. %ELFs are also a new addition to our extension of BM's paradigm. % as mentioned above.   %Note that our unified paradigm indicate OAEP and FO may use their hash functions in similar ways, something not apparent from  RO-model analyses.
 %To explain ELFs~\cite{C:Zhandry16}, we first recall the notion of a lossy function, a trapdoor-less version of lossy trapdoor functions~\cite{DBLP:journals/siamcomp/PeikertW11}. 
%A lossy function key can be generated in one of two modes, the injective or the lossy mode, where the first induces an injective function and the second induces a highly non-injective one.
%Furthermore, keys generated via these two modes are indistinguishable to any efficient adversary. Note that the lossy function image cannot be \emph{too} small, else there would be a trivial distinguisher.  
%ELFs achieve \emph{much more lossiness} by reversing the order of quantifiers. Namely, for an ELF, for every adversary there exists an (adversary-dependent) indistinguishable lossy key-generation mode. The induced function can even have an appropriate  \textit{polynomial}-size image. 
%Zhandry~\cite{C:Zhandry16} constructs ELFs based on exponential DDH, where the lossy mode depends on the run-time of the adversary. % ELFs isolate a significant but overlooked property of ROs that can surprisingly be realized under standard assumptions
%Zhandry shows ELFs can be used to instantiate ROs in a variety of contexts but does not address  chosen-ciphertext security.
 $\hashG := \iO  (\ELF \circ\PRF)$. %in a direct instantiation, where as usual $\iO$ is an indistinguishability obfuscator~\cite{C:BGIRSVY01,SIAM:BGIRSVH16}, $\ELF$ is an ELF in injective mode, and $\PRF$.
%Here the PRF  and ELF keys are hardcoded into the obfuscated program.
The main observation was: 
 \begin{quote}    \emph{ELFs allow answering decryption queries in a standard-model proof of chosen-ciphertext security for $\mathsf{FO}$ by simulating  ``observability'' used in the classic  ROM proof.}
 \end{quote}
Intuitively, in the proof, after we switch $\ELF$  to lossy mode, our reduction  to DL answers the adversary's decryption queries by enumerating the ELF's image. (This  comprises the space of   possible symmetric keys.) This is similar trying   each random oracle output so far, \emph{i.e.}~leveraging observability. 
%The problem is that the ELF output used to form  the challenge ciphertext will not look random to such a reduction, because it runs in long enough time to break the ELF.
%To resolve this, '
Because of technical issues, we also incorporated auxiliary-input multi-bit point function obfuscation (MB-AIPO)~\cite{EC:CanDak08} in the proof as well, extending an idea of~\cite{AC:BrzMit14a}.

%into the hash function as explained later in more detail. In the remainder, we focus on  $\mathsf{FO}$. %Intuitively MB-AIPO allows obfuscating any program $P_{x,y}$ that outputs $y$ on input $x$ and otherwise $\bot$. %In our proof for $\mathsf{FO}$, $x$ $y$ is the challenge symmetric key. %Unfortunately, the auxiliary input, \emph{i.e.}~the side information the adversary has on $c_1 \| r,K$ in our proof is  hard to analyze and in general MB-AIPO is subject to impossibility results~\cite{}. %Namely, incorporating MB-AIPO into $\hashG$ allows us to change $\hashG$ to output a truly random value on a hidden point used in formation of the challenge ciphertext.  

%MB-AIPO obfuscates  a program $P_{x,y}$ that on input $x'$ outputs $y$ iff $x=x'$; in our case, $x$ will be the randomness used in forming the challenge ciphertext.  %In essence, we change the hash function to  output a  \emph{truly random} point on the secret input used in forming the challenge ciphertext, and otherwise evaluates the ELF. % Notably, the main hash function we construct to instantiate FO and OAEP \emph{is the same in both cases}, which is what we mean by ``paradigm.''
  

%the latter also uses the (by now standard) punctured-programming approach of Sahai and Waters~\cite{STOC:SahWat14}. Thus, overall, our paradigm could be seen as replacing a component of an extension to the hash function construct of~\cite{AC:BrzMit14a} with an ELF. 

\iffalse

\heading{Results about OAEP.}
First, we show OAEP is fully instantiable for the sub-class of TDPs we call $\zeta$-FB-NM TDPs.  
But the TDP must also  satisfy some additional properties depending on what strength of auxiliary-input point-function obfuscation (AIPO)~\cite{C:Canetti97,TCC:BitPan12,TCC:BelSte16} exists \emph{non-constructively}. To instantiate $\hashG$, below we use iO~\cite{C:BGIRSVY01,SIAM:BGIRSVH16}, puncturable PRFs~\cite{STOC:SahWat14,CCS:KPTZ13,EC:BoyGilIsh15}, and ELFs~\cite{C:Zhandry16}; in the case assuming AIPO, $\hashH$ is arbitrary but fixed.
Recall that AIPO for $x$ is essentially an obfuscated program that outputs 1 if the input is $x$ (else $\bot$); here the adversary also gets some auxiliary information sampled jointly with $x$.
Moreover, then we need  the TDP to satisfy a ``non-blackbox'' version of second-input extractability (SIE) introduced in~\cite{CCS:BarPoiZan12}, which roughly requires that given an image point and part of the preimage, the full preimage can be recovered.  We call the black-box version BB-SIE and the non-blackbox one NBB-SIE2 (the `2' denoting interactivity). We also require ``one-wayness with hint,'' which is another way of saying BB-SIE should \emph{not} hold. As BB-SIE holds for low-exponent RSA unconditionally, this case requires the \emph{high-exponent} regime for RSA-OAEP. 
In the next case we assume  \emph{multi-bit AIPO} (MB-AIPO)~\cite{EC:CanDak08,AC:BrzMit14b} exists for a certain distribution. 
Note \emph{multi-bit AIPO}, where the obfuscated program outputs $y$ if the input is $x$ (else $\bot$) and the adversary has auxiliary information $z$ sampled jointly with $x,y$, is impossible for all distributions~\cite{AC:BrzMit14b}, though the special case we need may be realized below.
We additionally need that the TDP is \emph{black-box} SIE as well as partial one-way~\cite{C:FOPS01}.  Remarkably, we show we actually only need POW in the \emph{low-exponent} regime for RSA-OAEP; all the other properties needed either hold unconditionally or follow, including FB-NM. %MB-AIPO allows us to avoid this stronger assumption on RSA (which is non-blackbox and interactive). 
Finally, we provide an MB-AIPO construction for our needs based on RSA, requiring  a conjecture about its composability. %As NBB-SIE is implied by SIE one could view the move from AIPO to MB-AIPO as trading not-SIE for POW, which seem incomparable. 

Second, we turn to FO. Inspired by the modular analysis of FO by Hofheinz \emph{et al.}~\cite{TCC:HofHovKil17}, we first focus on instantiability of the part that transforms PKE from one-wayness under plaintext-checking attack (OW-PCA\footnote{This notion says it is hard to break one-wayness relative to an oracle that tells for an input $(m,c)$ if $c$ decrypts to $m$.}) to IND-CCA, namely transform 3.2.2 in~\cite{TCC:HofHovKil17}. 
We tweak this transformation (however, in certain cases the tweak adds no overhead) and make a novel  ``fixed-message non-malleability'' (FM-NM) assumption on the PKE in addition to OW-PCA.  Furthermore, we need MB-AIPO for a special distribution. The instantiation of $\hashG$ is the same as in OAEP. Combining this with the symmetric-key encryption being  authenticated encryption, we are able to show our transform is fully instantiable.  
Our modular analysis isolates the properties needed from PKE scheme, which may already be satisfied by some schemes.
Alternatively, one can use the encrypt-with-hash (EwH)~\cite{C:BelBolONe07} part of the $\mathsf{FO}$ transform for generically upgrading OW-CPA PKE. Unfortunately, EwH is uninstantiable~\cite{TCC:BrzFarMit15,FOCS:GoyKopWat17} so we make the additional assumption that the PKE scheme is \emph{lossy}~\cite{EC:BelHofYil09}, similar notions being introduced in \emph{e.g.}~\cite{TCC:KolNao08,C:PeiVaiWat08}) and use pairwise-independent hashing for the instantiation as per Hemenway and Ostrovsky~\cite[Corollary 2]{AC:HemOst13}.  Note that by~\cite{TCC:BrzFarMit15,FOCS:GoyKopWat17} any assumption weaker than IND-CPA will not suffice for EwH, whereas OW-PCA and FM-NM are weaker than IND-CPA (but not applied to an EwH scheme here).
%Namely, we show that EwH is then a lossy TDF in the terminology of~\cite{STOC:PeiWat08}, with enough lossiness to argue OW-CPA.  
Thanks to ``re-encryption on decryption,'' it can be seen we then get OW-PCA and FM-NM for free. 
Additionally, assuming lossiness, we only need an MB-AIPO for \emph{statistically} unpredictable distributions. We show a new such construction from ELFs.
. % For technical reasons, both the special input and the programmed point must be hidden even to the obfuscated program.  For this, we use multi-bit point function obfuscation~\cite{}.

%\heading{Hash function design and rationale.} \adam{Mohammad please fill this in.}
\fi

%\heading{Results on $\mathsf{OAEP}$.}
%For simplicity, we talk about results that pertain to messages chosen independently of the public key; otherwise we need to use sub-exponential hardness of the assumptions.
%We show that low-exponent RSA-OAEP is instantiatiable under the same assumption on the base scheme  used by Fujisaki \emph{et al.}~\cite{C:FOPS01} to prove chosen-ciphertext security in the ROM, partial-domain one-wayness. We also exploit algebraic properties that are unconditionally satisfied by low-exponent RSA.
%We show that OAEP is fully instantiable under suitable assumptions.  %First, we assume the TDP satisfies our notion of $\zeta$-FB-NM. \adam{Explain where this is used} %which  from standard assumptions for low-exponent RSA.%\footnote{ Shoup~\cite{JC:Shoup02} shows an NM assumption is necessary on the , but we make an NM assumption on the underlying TDP here.  
% As explained above, the additional assumptions needed on the TDP depends on the strength of point-function obfuscation assuming to exist. Then 
%Here we take into account public-key-dependent messages, and hence the sub-exponential assumptions we use.
%Here we instantiate $\hashG$ in OAEP as $\iO(\ELF(\PRF_K(\cdot)))$ as described above where $\iO$ is an indistinguishability obfuscator~\cite{C:BGIRSVY01,SIAM:BGIRSVH16}, $\ELF$ is an injective-mode ELF, and $\PRF$ is a puncturable pseudorandom function~\cite{AC:BonWat13,EC:BoyGilIsh15,CCS:KPTZ13}. 
%The PRF key and ELF function are hardcoded into the obfuscated program, in the punctured programming style of~\cite{STOC:SahWat14}. 
%To instantiate $\hashH$ we use a one-wayness extractor~\cite{EC:HsiLuRey07} with polynomial-length output. Note that ELFs are non-blackbox  so we circumvent~\cite{EC:KilPie09}.
%In the proof (and not in the construction), multi-bit point function obfuscation with auxiliary input (MB-AIPO) is used. %An MB-AIPO adversary is given the obfuscation and auxiliary information  depending on the real point function input and output points.  %In addition to this obfuscation, the adversary is given auxiliary information depending on the real input and output points. Their goal is to guess whether the output point of the obfuscated function is real (consistent with the auxiliary information) or random (independent of the auxiliary information).

%which can also be constructed from ELFs~\cite{C:Zhandry16} or iO~\cite{AC:BelSteTes14}.


% The Fujisaki-Okamoto (FO) transform~\cite{} takes a ``weak'' public-key encryption scheme and lifts it to chosen-ciphertext security by using it in a hybrid construction with a  symmetric-key encryption scheme. 
Following Hofheinz \emph{et al.}~\cite{TCC:HofHovKil17}, we  view $\mathsf{FO}$ as consisting of two sub-transforms.
We will first focus on instantiability of the sub-transform of $\mathsf{FO}$ that uses $\hashG$  to upgrade PKE satisfying one-wayness under plaintext-checking attack (OW-PCA) to chosen-ciphertext security (IND-CCA), cf.~transform 3.2.2 of Hofheinz \emph{et al.}~\cite{TCC:HofHovKil17}. In particular, this sub-transform is \emph{not} subject to an uninstantiability result of~\cite{TCC:BrzFarMit15} --- that result applies to the other sub-transform that uses Encrypt-with-Hash~\cite{C:BelBolOne07}, which we discuss later.  In OW-PCA, the adversary may submit queries $(c,m)$ to learn if $c$ decrypts to $m$. In our ASIACRYPT 2022  work, we instantiated a variant of this sub-transform where encryption is:
\[
\calE^{\hyb}_{pk}(m;r) = 
\calE^{\asy}_{pk}(r;z) \| \calE^{\sy}_{K}(m \| r) 
\quad \mbox{where} \quad
K = \hashG(r\| \calE^{\asy}_{pk}(r;z))
\]
where $z$ denotes random coins for $\calE^{\asy}_{pk}(\cdot)$.
Decryption first recovers $r$ from the asymmetric  part of the  ciphertext, computes the symmetric key with the hash function, and then decrypts the symmetric part of the ciphertext to get $m\|r'$. Finally, $m$ is returned iff $r'=r$, otherwise $\bot$. Here, we needed to assume that   $\mathsf{SE}$ is \emph{authenticated encryption}. 
%Unfortunately, this result has major shortcomings: 
%\begin{newenum}
    
\heading{Main question.} Our prior result introduces useful ideas but has significant shortcomings: the instantiation applies  to a \emph{variant} of $\mathsf{FO}$ not used in practice,  it is not \emph{post-quantum},  and it relies on heavyweight (iO) or suspect (MB-AIPO) assumptions.  These combine to form a  structured hash function bearing little resemblance to cryptographic hashing used in practice. Thus, we ask:

%\item The  instantiation  only pertains to a variant of $\mathsf{FO}$ not used in practice.
%\item The instantiation is not \emph{quantum-secure}, in particular current ELFs  rely on elliptic-curves.
%\item  The instantiation uses the  heavy-hammers of iO and auxiliary-input MB-AIPO.
%\item  The notions employed to build %$\hashG$ cannot plausibly be achieved by cryptographic hashing.
%\end{newenum}
%In sum, our prior result gives little practical insight into $\mathsf{FO}$. To fix this, we ask:
\begin{ques}
    Can we instantiate the Fujisaki-Okamoto transform using notions that can  be: (1) plausibly satisfied by cryptographic hashing and (2) realized under post-quantum assumptions?
\end{ques}

%\noindent Here, we leave open what kind of assumptions we will use for now.
% In our  instantiation we replace  $\hashG$ with $\iO(\ELF(\PRF_K(\cdot)))$ as discussed above.
%Note that this modificated transform is basically as efficient as the original; if the base symmetric-key encryption is randomness-recovering then $r$ can simply be used as coins for $\calE^{\sy}_{K}(\cdot)$.
%Moreover, if the symmetric-key encryption is already \emph{randomized} and \emph{randomness-recovering}, then $r$ can safely be used as its coins as there is no additional overhead. 
\iffalse
Overall, we highlight the changes/novelties as compared to~\cite{TCC:HofHovKil17}:
\begin{newitemize}
  %This will allow us to avoid the stronger OW-PCVA assumption (see~\cite{TCC:HofHovKil17} for the definition) on the public-key scheme, which the adversary also has a ciphertext validity oracle.
\item The transform in~\cite{TCC:HofHovKil17} is ``KEM-to-KEM'' whereas ours takes advantage of the interplay between the public-key and symmetric-key encryption schemes, as in the original FO transform.
\item We include $r$, the ``message'' encrypted by the public-key scheme, in the message encrypted by the symmetric-key scheme, and check that the values match upon dehttps://en.wikipedia.org/wiki/Garbled_circuitcryption.
    \item We do not use ``implicit rejection''   to intentionally suppress decryption error messages about the starting PKE scheme.  %Indeed, a ciphertext validity check cannot be performed 
    %The concern here is whether the IND-CCA adversary is able to use its decryption  oracle to simulate a ``ciphertext validity'' oracle.  In our case, however, give first cip it is hard for the adversary to come up with c_2 to make the decryption query c_1,c_2 to the adversary. In the random oracle this fact was captured by the fact that G is random oracle and adversary need to query r on oracle G to be able to for corresponding c_2 for c_1 however this is not possible due to PKE being OW-PCA. In the standard model this fact is captured by ELF assumption on G. The adversary cannot come up with corresponding c_2 for c_1 given c_1 and G(r) unless it invert c_1 which is impossible due to PKE being OW-PCA. 
\end{newitemize}
\fi

%In particular, we do not perform ``implicit rejection'' (returning random when a ciphertext is invalid) and we encrypt the encapsulated key as part of the message for the symmetric-key encryption scheme.
%We also explicitly model the symmetric-key encryption scheme (\emph{i.e.}, we treat the whole hybrid encryption rather than just KEM),
%Note that in practice we use pseudorandom generator to  transforming a short, uniform string called the seed $r$ into a longer, pseudorandom output string which we usually use as randomness in the encryption schemes. Thus, the changes are conservative.   

%In practice, the encapsulated key would result in only 128 bits overhead in the message length for the symmetric-key encryption scheme.  Given our instantiability result, our modifications may be preferable in practice. 
%We show this modified part of the FO transform is instantiable under suitable assumptions.  % where the input point has the form $r^* \| c_1^*$, the output point is $K^*$, and the auxiliary input has the form $(t,d,c^*,\pk',m)$ where $c^* = c^*_1 \| c^*_2$ is an encryption of $m$.  
%Furthermore, the oracle provided to the adversary is either a public-key ciphertext validity checker or, as a separate assumption, a symmetric-key ciphertext validity checker. 
%We instantiate $\hashG$ exactly as we do in $\mathsf{OAEP}$ above, illustrating that $\mathsf{OAEP}$ and $\mathsf{FO}$ rely on related properties of the hash functions, which is not obvious \emph{a priori}.
%Finally, if the base  PKE scheme is \emph{lossy}~\cite{EC:BelHofYil09}, we can instantiate $\hashH$ in the part of $\mathsf{FO}$ that uses Encrypt-with-Hash~\cite{C:BelBolONe07}  by~\cite[Corollary 2]{AC:HemOst13}. %, we can instantiate $\hashH$ using a combinatorial hash function. However, some schemes of interest such as Kyber are not lossy~\cite{}.
% This does not  contradict~\cite{TCC:BrzFarMit15} because  lossiness is stronger than IND-CPA.  
 
 \subsection{Instantiating Fujisaki-Okamoto: Practice-Oriented and Post-Quantum} 


% While  the PI's prior work was a good first step that observed the power of ELFs here, our proposed work aims to give improved instantiations of $\mathsf{FO}$ that are much more practically relevant. 
 
%We will solve all  major shortcomings of the prior instantiation,, which are: (1) The  instantiation  only pertains to a variant of $\mathsf{FO}$; (2) The instantiation is  \emph{quantum-insecure}  because  ELFs  rely on elliptic-curves; (3) The instantiation uses the  heavy-hammers of iO and auxiliary-input MB-AIPO (which is impossible in general~\cite{}). (4) The  notions cannot be plausibly achieved by cryptographic hashing. In sum, our prior instantiation gives little practical insight. We therefore ask:

%Again, we mainly focus on the sub-transform of $\mathsf{FO}$ that upgrades OW-PCA secure PKE to iND-CCA secure PKE. We touch on instantiability of the other part of the transform later.
%There even more problems as well, such as complexity leveraging, which makes the security proof loose. %Thus, while of interest as a theoretical feasibility result, the prior instantiation remains quite unsatisfactory from a practical point of view. 

 \iffalse
 In particular, we delineate some key drawbacks to the prior instantiation result as follows:
 \begin{enumerate}
 \item The prior instantiation result applies to a variant of $\mathsf{FO}$ transform that is not used in practice.
     \item The prior instantiation result relies on  indistinguishability obfuscation (iO), a heavyweight assumption that has as-yet unclear post-quantum viability.
     \item The prior instantiation result uses  auxiliary-input multi-bit point function obfuscation, a non-standard assumption subject to negative results~\cite{}.
     \item The prior instantiation uses a ``complexity leveraging'' argument in the case of public-key-dependent messages, which leads to a  loose security bound.
     \item The only known construction of ELFs is quantum-insecure. 
To obtain a much better instantiability result, we will  to introduce several novel cryptographic notions. 


%     \item The lossiness assumption needed to instantiate $\hashH$ is not satisfied by \emph{e.g.}~$\mathsf{Kyber}$.
 \end{enumerate}
\fi

%\begin{ques}
%    Can we give a post-quantum instantiation of the   Fujisaki-Okamoto transform using notions that could  be satisfied by cryptographic hashing and built under post-quantum assumptions?
%\end{ques}
%We will introduce  a handful of  new tools before proceeding with our instantiation.

To answer the above question, we   introduce several new notions of  independent interest.


\subsubsection{New Cryptographic Toolbox} \label{sec-tools}

\heading{Tunable lossy functions.}  ELFs seem hard to build under  post-quantum assumptions, so we will develop new \emph{relaxations} that can be. In particular, we propose \emph{tunable lossy functions} (TLFs), which  interpolate between  standard LFs in the sense of~\cite{JACM:PeiWat11} and ELFs. In a TLF, the lossy key can be generated so the image has sub-exponential or even quasi-polynomial size. We can build such a post-quantum TLF from  the sub-exponential  \emph{learning-with-errors} (LWE) assumption based on the LF from LWE of~\cite{JACM:PeiWat11}. Such a TLF is   sufficient because we can ``offload'' missing hardness onto other primitives, which can themselves be realized from sub-exponential LWE.  %As mentioned  in Section~\ref{sec-kr}, we will investigate ``targeted'' notions of TLFs~\cite{}  as well.  %Indeed, an ELF is = ``complexity absorbing'' in that, while it relies on  exponential hardness, other primitives in the higher-level construct can rely on  polynomial-hardness. (In the proof, the adversary against the latter  will  enumerate the ELF's lossy image.) But if these other primitives  use LWE, we can    ``balance'' the hardness of LWE to sub-exponential overall.

\iffalse
\heading{Proposed instantiation.}
We propose to instantiate  $\hashG$ in   $\mathsf{FO}$ with the hash function $\mathsf{TLF}(\mathsf{UCE}(\cdot))$, where $\mathsf{UCE}$ is universal computational extractor~\cite{C:BelHoaKee13} and $
\mathsf{TLF}$ is a tunable lossy function. In other words, we will completely eliminate the use of iO and MB-AIPO as compared to our prior work. Moreover, both TLFs and UCEs can be constructed in theory under post-quantum assumptions. The high-level intuition for the proof is that the TLF is used to emulate ``observability'' in the ROM and the UCE, which can be seen as a hardcore function for any computationally unpredictable source, is used to extract a randomly-looking symmetric key.
\fi


\heading{Security under functional plaintext-checking attack.}
We consider a novel assumption on a public-key encryption scheme  we call \emph{one-wayness under functional plaintext-checking attack} (OW-FPCA), defined wrt.~a function  $F$. In OW-FPCA wrt.~$F$, the adversary is given $(\pk, c^*)$ where $c^*$ encrypts a random message $m^*$ under public key $\pk$. and tries to output $m^*$ given an oracle that on input $(c,y)$ returns 1 iff $F(m) = y$, where $m$ is the decryption of $c$ under $\sk$.
%Intuitively, OW-FPCA interpolates between OW-PCA and OW-CCA.

In our result, roughly, $F = \hashG(K,\cdot)$. (We gloss over the fact that a ciphertext is also an input to $\hashG$.) 
Ultimately, we want to  ``offload'' the  additional requirement on $\PKE$ to  hash function $\hashG$ in the case $\PKE$ is \emph{deterministic}, which it will be due to the Encrypt-with-Hash sub-transform:

\begin{conj} \label{conj-ext}
    If $\PKE$ is \emph{deterministic} (\emph{i.e.}, is a trapdoor function) and one-way, and $\hashG$ is \emph{extractable}~\cite{TCC:CanDak09}, then $\PKE$  is OW-FPCA wrt.~$\hashG$.
\end{conj} 
%In other words, if $\hashG$ is extractable then we can completely avoid any new assumptions on the PKE.

\noindent Above, a hash function if extractable~\cite{TCC:CanDak09} if the only way to produce a hash value is to ``know'' a preimage of it.  We discuss this assumption more below. %for every adversary that outputs a hash value, there is a non-blackbox extractor getting the adversary's code and coins that outputs some preimage.
%Our prior work at PKC 2020~\cite{PKC:CaoONeZah20} defines a ``many-times'' notion of extractability we may use.
%Thus, at the end of the day, we still want to make the usual OW-PCA assumption on the base PKE scheme. %he idea for the proof is that the OW-PCA adversary runs the extractor on input $y$ whenever the OW-FPCA adversary makes a query $(c,y)$, to get the message $m$. It then queries $(m,c)$ to its own oracle, forwarding the response. 

\heading{UCE with adaptive leakage.} A \emph{universal computational extractor}~\cite{C:BelHoaKee13} $\mathsf{UCE}$ is another powerful notion for  instantiating random oracles. It asks that hash values look random given side information on the inputs.   More formally, the adversary is split into ``source'' $S$ and ``distinguisher'' $D$. The game samples a  bit $b$ and a hash key $K$. Then, $S$ runs with access to an oracle that on input $x$ returns $\mathsf{UCE}(K,x)$ if $b = 1$ and random otherwise. Eventually, $S$ outputs ``leakage'' $L$, which is passed to $D$, who is also given $K$ but no oracle and tries to guess $b$. For non-triviality,  $S$'s queries must be \emph{unpredictable} given $L$. 
We only need \emph{one-query} UCE, where $S$  makes one query.  It will be important for us that one-query UCE has \emph{post-quantum} realizations~\cite{EC:CheMao25}.

We introduce (one-query) UCE with \emph{adaptive leakage}, where $L$ can  depend on $K$. Note that if $L$ depends arbitrarily on  $K$,  there is a trivial attack. However, we consider a game where the adversary outputs 
split leakage $L = (L_1,L_2)$ such that $L$ determines $x$ but $x$ remains unpredictable given either part. Furthermore, we split  $D$ into $D_1$ and $D_2$.  First, $D_1$ is given hash key $K$ and $L_1$ to output leakage $L_3$. Now, $D_2$ is given $K$ and $L' = (L_2,L_3)$ and tries to guess $b$. For non-triviality,  $x$ must be unpredictable given $L'$. This notion  will be quite natural in our  application as $L_1 = \sk$ for the base PKE scheme and $L_2 = (\pk,c)$ where $c$ encrypts  $x$ under $\pk$.  To justify it, we  will show:
\begin{conj} \label{conj-fo2}
The post-quantum one-query UCE of~\cite{EC:CheMao25} is secure with adaptive leakage.
\end{conj}

We will also show that a random oracle satisfies our new notion; in this sense, it can act as a design criterion for cryptographic hashing.. %Also, we note that Farshim and Mittelbach~\cite{FM16} proposed a much more complex adaptive extension of UCE.

%As a sanity check, we will also show a random oracle is a one-query UCE wrt.~adaptive leakage. 
%In other words, we aim to show that UCE for adaptive leakage continues to be a plausible heuristic for cryptographic hashing as well as realizable in theory from post-quantum assumptions. 
%, and thus one-query UCE for adaptive leakage continues to be a plausible assumoption for cryptographic hashing.
%In fact, under natural conditions on the PKE, the  TDF obtained via  Encrypt-with-Hash on a OW-CPA PKE schceme is an adaptive TDF in the ROM.

%also use auxiliary-input point function obfuscation (AIPO), which is much weaker than MB-AIPO and plausibly satisfied by cryptographic hashing. 

\subsubsection{Proposed Result} \label{sec-inst-fo}


%We sketch how the above pieces could fit together to give the desired instantiability result for $\mathsf{FO}$.

Our preliminary instantiation of $\hashG$ in the  the OW-PCA-to-IND-CCA sub-transform  of $\mathsf{FO}$  
is: \begin{newmath}\hashG := \mathsf{TLF} \circ \mathsf{UCE} \;, \end{newmath}%
that is, the key for $\hashG$ consists of a key for $\mathsf{TLF}$, a tunable lossy function and a key for $\mathsf{UCE}$,  a one-query  UCE for adaptive leakage.% We stress that here we instantiate the OW-PCA-to-IND-CCA sub-transform of $\mathsf{FO}$; we instantiate  the other sub-transform below. We aim to prove:
\begin{conj} \label{conj-fo}
    Let $\hashG$ be as above. Then the  OW-PCA-to-IND-CCA sub-transform of  $\mathsf{FO}[\PKE,\mathsf{SE}]$ is IND-CCA secure assuming   $\PKE$ is OW-FPCA wrt.~$\hashG$ and $\mathsf{SE}$ is IND-CPA.
\end{conj}
Our  preliminary instantiation  has   simple intuition: 
 \begin{quote}
\emph{$\mathsf{TLF}$ allows simulating ``observability'' in the ROM to answer decryption queries,  while $\mathsf{UCE}$ allows extracting a uniform symmetric key for the challenge ciphertext.}
\end{quote}

\noindent Compared to our prior result, using OW-FPCA wrt.~$\hashG$ instead of OW-PCA on the base PKE here addresses numerous technical issues and we do not resort to a non-standard variant of $\mathsf{FO}$. Moreover, while UCE is a ``two-stage'' assumption, the \emph{one-query} restriction (arising from the single challenge ciphertext  in IND-CCA) makes it easier (in our view) to reason about than even full-fledged UCE without adaptive leakage.

% In the proposed work, we will also work out the exact  post-quantum assumptions under which we can build our uanderlying primitives.

 \subsubsection{Completing the Fujisaki-Okamoto Instantiation} \label{sec-fuj-cont}
 
\heading{Instantiating Encrypt-with-Hash.} Finally, we  need to instantiate the other sub-transform of $\mathsf{FO}$ that  derandomizes the base PKE scheme $\PKE$  using hash function $\hashH$, known as  Encrypt-with-Hash~\cite{C:BelBolOne07}. Then we can  apply Conjecture~\ref{conj-fo2}  and  Conjecture~\ref{conj-fo} assuming $\hashG := \mathsf{TLF} \circ \mathsf{UCE}$ is extractable to conclude instantiability of $\mathsf{FO}$. Thus, we ask:
 \begin{ques}
     What assumptions on $\PKE$ and $\hashH$ imply  Encrypt-with-Hash~\cite{C:BelBolOne07} is  one-way?
 \end{ques}
Instantiating Encrypt-with-Hash is challenging because~\cite{C:BrzFarMit14} show it require a property of $\PKE$ not implied by IND-CPA. Our ASIACRYPT 2022 work~\cite{AC:MurONeZah22} observed that \cite{AC:HemOst13} yields an instantiation of Encrypt-with-Hash when $\PKE$ is \emph{lossy}~\cite{EC:BelHofYil09} and $\hashH$ is pairwise-independent. However, NIST-selected $\mathsf{Kyber}$~\cite{kyber} is not lossy. We will explore alternative  properties.

 \heading{Guiding hash function design.} Ultimately, we  seek assumptions on $\hashG$ and $\hashH$ that both admit post-quantum realizations and can also separately be \emph{plausibly achieved by cryptographic hashing}, bringing to light what hash function properties are needed for emerging  standards. In preliminary work, we have focused on $\hashG$. Key questions here include:
\begin{newitemize} 
     \item  \emph{Can we build a \emph{monolithic} hash function $\hashG$ that suffices to instantiate $\mathsf{FO}$? }
     Using a composition of distinct hash functions for the instantiation is hard to interpret in practice.
     \item \emph{Can we use a relaxation of TLFs achievable by cryptographic hashing?} Or, \emph{is there another mechanism that can simulate ROM observability in the standard model?} Recall TLFs cannot be constructed in the ROM~\cite{TCC:FisRoh23}.
    % A similar issue was discussed in Section~\ref{sec-kr}.
     \item \emph{Can we use a relaxation of  extractability  achievable by cryptographic hashing?} In practice, $\hashG$ will not be extractable since its output space is the set of all 
     $128$-bit keys. One  possibility could be a standard-model analogue of ``preimage-awareness''~\cite{EC:DodRisShr09}.  
\end{newitemize} 

By comprehensively addressing these questions, we will ensure maximal practical impact. Too often instantiability work remains purely theoretical; here it will be tied directly to practice, which is crucial for standards and deployment. While beyond the scope of this project, future work could argue the design criteria we identify for cryptographic hash functions from assumptions on their underlying compression functions or permutations, further grounding our results in practice. %Ultimately, we will  aim to use notions that one can be study for cryptographic hash functions via suitable assumptions on  lower-level their compression functions or permutations, as in~\cite{C:BelHoaKee14}.


%To see this, let's zoom in on why we needed a variant of  $\mathsf{FO}$ in prior work.  It is a step in the proof where we build an OW-PCA adversary against the base PKE scheme that runs the assumed IND-CCA adversary. At this point, the ELF is in lossy mode, so to answer decryption query $c_1 \| c_2$ the OW-PCA adversary enumerates the image of the ELF and checks which one correctly decrypts $c_2$. For each $K$, the OW-PCA adversary can check if  $K$  decrypts $c_2$ to some valid message. But before returning the message, the OW-PCA adversary needs to check that  $c_1$ encrypts an $r$ such that $\hashG(r \| c_1) = K$. This is enabled precisely by $c_2$  encrypting $r$ in addition to $m$. By switching to OW-FPCA, we avoid modifying $\mathsf{FO}$ and solve numerous other issues.    % Thus, to check the forgoing condition the IND-PCA adversary can check that $\hashG(r \| c_1) = K$ and that the plaintext-checking oracle on input $(r,c_1)$ returns 1. 

\iffalse
\subsection{Thrust Summary}

This thrust will show how  lossiness-based  hashing tools yield the first instantiation results for Schnorr signatures (used in TLS and Bitcoin)  and the Fujisaki-Okamoto transform (which underlies all NIST PQC encryption schemes). In particular, our tunable LF + adaptive UCE approach to instantiating  Fujiskai-Okamoto  is the first practical, post-quantum solution. Combined with our previous thrust, this completes  our  study of real-wold viability of ROM schemes: tight proofs under natural assumptions, and instantiability results under plausible hash-function properties.
\fi


 %then we can in particular obtain a one-way TDF by applying the  ``Encrypt-with-Hash'' step of $\mathsf{FO}$ and instantiating $\hashH$ with a combinatorial hash function as per~\cite[Corollary 2]{AC:HemOst13}. %However, as we later discuss, some such schemes of interest in practice like Kyber are not lossy. %However,  many randomized PKE schemes of interest such as Kyber are not lossy.  %(Note that this part of the transform \emph{is} subject to the uninstantiability result of~\cite{TCC:BrzFarMit15}, but lossiness is stronger than IND-CPA.) Indeed, we  obtain an instantiation of our modified $\mathsf{FO}$ transform when it is applied to any  one-way trapdoor function.
%\heading{Our goal: Post-quantum inseetantiability.}

\iffalse
\subsubsection{Saving Bandwidth in Code-Based Cryptography using  $\mathsf{OAEP}$}
% Indeed, $\mathsf{OAEP}$, while more bandwidth-efficient than $\mathsf{FO}$, is not often considered in post-quantum cryptography (PQC), whereas $\mathsf{FO}$ is used pervasively. Why is this? A primary  reason  appears to be that 

%The $\mathsf{FO}$ transform readily applies to the post-quantum setting but $\mathsf{OAEP}$ does not.  Therefore, we will first extend $\mathsf{OAEP}$ to this PQC \emph{in the ROM} and even given an application to watermarking large language models, which is clearly of independent interest.

Before treating post-quantum instantiability, we  will treat  the use of $\mathsf{OAEP}$ in post-quantum cryptography (PQC) in the first place. Despite the fact that $\mathsf{OAEP}$  is  more bandwidth-efficient than $\mathsf{FO}$ because it encrypts part of the starting TDP,  $\mathsf{OAEP}$ is not used much in PQC whereas $\mathsf{FO}$ is pervasive. Actually  


The main reason seems to be that that in PQC we do not have ( TDPs but rather trapdoor \emph{functions} (TDFs). 
Such TDFs come about in two ways.
The first is via the Encrypt-with-Hash~\cite{C:BelBolONe07} step  of $\mathsf{FO}$ applied to some randomized PKE scheme. However, the base randomized PKE scheme is often assumed IND-CPA, so $\mathsf{OAEP}$ is not needed. 
The second way is from \emph{code-based cryptography}; in particular, the Neiderriter TDF~\cite{}. Additionally,  the McElice randomized PKE scheme~\cite{} is randomness-recovering (\emph{i.e.}, the decryptor recovers the coins of the encryptor) and hence can  be viewed as a TDF. This domain \emph{can}{ benefit from $\mathsf{OAEP}$ because these TDFs are only assumed one-way, and thus cannot be used to encrypt a message directly without first applying $\mathsf{OAEP}$. 
Note we just need the resulting PKE scheme to be IND-CPA; then we can obtain an IND-CCA secure hybrid encryption encryption scheme via use the second step of $\mathsf{FO}$. 
Thus, we raise the question: % Indeed, Classic McElice~\cite{}, a finalist of the NIST PQC competition, applies a variant $\mathsf{FO}$ to the Niederriter TDF. (The variant includes another hash output in the ciphertext, which is needed for security in the QROM.)

\begin{ques}
    What properties of a trapdoor function (TDF) are necessary so that applying  $\mathsf{OAEP}$  to it is IND-CPA secure in the (quantum) ROM?
    \end{ques}

Targhi and Unruh~\cite{} previously  analyze $\mathsf{OAEP}$ applied to TDFs in the \emph{quantum} ROM, but they modify $\mathsf{OAEP}$ in a way that makes it less  efficient.  NTRU~\cite{} later adopted this modified version.  Later work analyzing $\mathsf{OAEP}$ without any modification in the quantum ROM only considers trapdoor \emph{permutations}~\cite{}. In the proposed work, we will  fill the gap and analyze plain $\mathsf{OAEP}$ applied to TDFs in the quantum ROM. In particular, we aim to use partial-domain one-wayness~\cite{} on the TDF, and we will analyze whether this property is satisfied by Neiderriter and McElice. Additionally, note that regardless of the  underlying padding transform used, the resulting PKE scheme will not be not IND-CCA for \emph{every} partial-domain one-way TDF. The reason is the TDF may be malleable in the sense that given $f,y$ one can efficiently find $y' \notin \mathsf{Image}(f)$ such that $f^{-1}(y') = f^{-1}(y)$.  Analogously to $\mathsf{FO}$, we will analyze two ways of dealing with this: Performing a re-encryption check on decryption, or assuming ``rigidity'' of the TDF~\cite{}, which has been shown for Neiderriter~\cite{}. Finally, we will concretely  calculate the bandwidth savings of using $\mathsf{OAEP}$ rather than $\mathsf{FO}$ for specific schemes at different security levels and explore applications  where this savings is important.
\fi

\iffalse
% We will also extend $\mathsf{OAEP}$ to encrypt long messages without using hybrid encryption. 
%The  idea is to define $\mathsf{OAEP}^*$ to take a TDF $\calF$ with input length $k$ and produce a PKE scheme whose public key is an instance $f$ of the TDF. It uses two ROs $\hashG,\hashH$, and the encryption algorithm has the form: 
%\[
%\calE^{\mathsf{OAEP}^*}_{f}(m;r) = f(x_1..x_k) \| x_{k+1}..x_{|x|-k} \quad \mbox{where} \quad s = \hashG(r) \xor m \| 0^\zeta,~~t = \hashH(s) \xor r,~~x = s \|t
\enspace,
%\]
% where $\zeta \in \N$ is the redundancy length and $s_i .. s_j$ denotes bits $i$ through $j$ of a string $s$. Note that as opposed to $\mathsf{FO}$, we encrypt part of the message under the pubic-key part of the scheme. One benefit of using $\mathsf{OAEP}^*$ to encrypt long messages versus using  $\mathsf{OAEP}$ in hybrid encryption is that we should only need a one-wayness, not partial-doimain one-wayness, assumption on the TDF to encrypt long messages. However, cryptographic hashing is slower than authenticated encryption (AE) used in $\mathsf{FO}$. To fix this, we will  introduce a novel variant of $\mathsf{OAEP}^*$ of independent interest, which is competitive with $\mathsf{FO}$ in terms of computational efficiency. Namely, we will set  $s = \calE_K(m \| 0^\zeta)$ where $K = \hashG(r)$ and $\calE$ is the encryption algorithm for an AE scheme.% In particular, we want to find a security notion for AE under which this is secure.
 \iffalse
Although it seems to be overlooked in the  literature, such a gain in bandwidth is also possible with $\mathsf{FO}$ \emph{but only if the base PKE scheme is IND-CPA}; \emph{e.g.}, in the case of NIST PQC competition winner Kyber~\cite{}. That is, define a corresponding transform $\mathsf{FO}^*$ like this:
 \[
\calE^{\hyb}_{pk}(m_1 \| m_2;r) = 
\calE^{\asy}_{pk}(m_1 \| r; \hashH(r)) \| \calE^{\sy}_{K}(m_2) 
\quad \mbox{where} \quad
K = \hashG(r \| c_1), c_1 = \calE^{\asy}_{pk}(r; \hashH(r))
\enspace,
\]
where we parse a message $m$ as to parts $m_1 \| m_2$ of appropriate length.  If the base PKE scheme (which can be either randomized or deterministic) is only \emph{one-way} then $\mathsf{FO}^*$ is clearly insecure. 

\fi
%Note that if we start with a randomized PKE scheme and apply Encrypt-with-Hash, we only need $\mathsf{OAEP}$ when the starting scheme is one-way instead of IND-CPA, although $\mathsf{OAEP}$ may still be useful here as a hedge. The main setting we have such one-way schemes is code-based cryptography, specifically, the Neiderriter TDF and (randomized) McElice PKE. The latter is randomness-recovery so can be viewed as a TDF without applying Encrypt-with-Hash. Thus, the primary application of $\mathsf{OAEP}^*$ will be to code-based schemes.

%We could have alternatively defined $\mathsf{OAEP}^*$ to use hybrid encryption.   Therefore, we will also consider a variant of $\mathsf{OAEP}^*$ that uses authenticated encryption in a novel way:
%\[
%\calE^{\mathsf{OAEP}^*}_{f}(m;r) = f(x_1..x_k) \| x_{k+1}..x_{|x|-k} \quad \mbox{where} \quad s = \hashG(r) \xor m \| 0^\zeta,~~t = \hashH(s) \xor r,~~x = s \|t
%\enspace,
%\]
% For example, given $pk$ the adversary could trivially distinguish between encryptions of $m'$ and $m''$ such that for every  $r$,  $\calE^{\asy}_{pk}(m' \| r; \hashH(r))$ starts with a zero bit but  $\calE^{\asy}_{pk}(m''' \| r; \hashH(r))$.


\iffalse
Overall, we have the following tasks:
\begin{itemize}
    \item Extend $\mathsf{OAEP}$ and $\mathsf{OAEP}^+$ to trapdoor functions and characterize security of it in this case.
    \item Extend the scheme to use authenticated encryption for $\hashH$ instead of cryptographic hashing.
    \item Calculate the concrete savings of using these schemes versus $\mathsf{FO}$ for important PQC schemes.
    \item Identify applications where encryption of short messages is important, as well as applications of hybrid encryption where bandwidth is at a premium.
\end{itemize}
 \fi\fi

%\heading{Proposed work.}
%A major problem with the above is  we really want \emph{post-quantum instantiations}, as these encryption transforms are used in a post-quantum setting. In other words, we want to instantiate these schemes and show security against quantum adversaries in the standard model. However, \emph{existing ELFs are not post-quantum, and constructing post-quantum ELFs is a   difficult problem we do not anticipate solving.} Instead, we will   look at relaxations and variants of ELFs that  suffice in our applications yet are realizable using post-quantum assumptions.

\iffalse
\subsubsection{Making $\mathsf{OAEP}$ Work for Code-Based Cryptography}


\iffalse
\heading{Encrypting short messages.}
We will first consider encryption of \emph{short} messages under PKE. Such a setting has a number of applications. For example, Oblivious DoH (ODoH) and related privacy‑preserving HTTP systems encrypt short request tokens, query identifiers, and path hashes. Similarly, IoT provisioning and secure bootstrapping encrypts short onboarding payloads (\emph{e.g.},  device ID, configuration blob, timestamp). Suppose we want to use post-quantum IND-CCA secure PKE. Conventionally,  we would apply $\mathsf{FO}$ to a base post-quantum PKE scheme, and  the resulting hybrid encryption scheme would have ciphertexts $(c_1,c_2)$ totaling about $k + |m|$ bits.

Specifically, recall that in the resulting hybrid encryption scheme, $c_1$ is used to encrypt randomness used to derive a symmetric-key. What if 

\iffalse
\iffalse
\begin{itemize}
    \item Zether, a confidential smart‑contract framework on Ethereum, where encrypted bids and account tags often require embedding metadata alongside numeric values—these fit within FrodoPKE’s direct message size;
%• Oblivious DoH (ODoH) and related privacy‑preserving HTTP systems, where encrypted request tokens, query identifiers, and path hashes are typically under 128 bytes and currently use KEM+DEM—FrodoPKE allows replacing that with pure PKE, reducing complexity;
%• IoT provisioning and secure bootstrapping, where onboarding payloads (e.g. device ID, configuration blob, timestamp, MAC) often fit in ≤128 bytes. FrodoPKE enables stateless, post‑quantum, low-entropy configuration encryption without hybrid layers.
%\end{itemize}
This setting has 
improve bandwidth efficiency for short messages across the board and bandwidth efficiency for long messages for code-based cryptography.
Before improving instantiability of the transforms, we will develop a method for reducing bandwidth in PQC using $\mathsf{OAEP}$, which is of independent interest. The main idea is to apply OAEP to a TDF that is either obtained directly from code-based cryptography or from applying the Encrypt-with-Hash deradomization step of $\mathsf{FO}$
The main idea is to use $\mathsf{OAEP}$
We propose to investigate the direct encryption of structured 128-byte messages using post-quantum trapdoor functions (TDFs)—specifically FrodoPKE, which supports IND‑CPA‑secure payloads up to ~128 bytes [FrodoSpec]. This enables CCA‑secure, derandomized encryption (e.g., using a Fujisaki–Okamoto style transform [Fujisaki00]) of self‑contained messages, avoiding hybrid encryption entirely.
This is particularly compelling for real-world systems such as:
• Zether [Zether20], a confidential smart‑contract framework on Ethereum, where encrypted bids and account tags often require embedding metadata alongside numeric values—these fit within FrodoPKE’s direct message size;
• Oblivious DoH (ODoH) and related privacy‑preserving HTTP systems, where encrypted request tokens, query identifiers, and path hashes are typically under 128 bytes and currently use KEM+DEM—FrodoPKE allows replacing that with pure PKE, reducing complexity;
• IoT provisioning and secure bootstrapping, where onboarding payloads (e.g. device ID, configuration blob, timestamp, MAC) often fit in ≤128 bytes. FrodoPKE enables stateless, post‑quantum, low-entropy configuration encryption without hybrid layers.
Our work will include implementation of these use cases, performance benchmarking, and formal modeling of derandomized FrodoPKE encryption for real-world application contexts like IoT, blockchain, and privacy-preserving protocols.
with an IND-CPA secure PKE scheme such as Kyber and want to upgrade it to IND-CCA. When bandwidth is at a premium, our idea is to modify $\mathsf{FO}$ to improve ciphertext rate, defined as the length of the ciphertext divided by the length of the message. Namely, we now encrypt part of the message under the PKE scheme, \emph{i.e.}~we  parse a message as $m = m_1 \| m)2$ and consider 
\[
\calE^{\hyb}_{pk}(m_1 \| m_2;r) = 
\calE^{\asy}_{pk}(r'; \hashH(r')) \| \calE^{\sy}_{K}(m_2)~\mbox{where}~r' = r \| m_1, K = \hashG(r' \| c_1), c_1 = \calE^{\asy}_{pk}(r'; \hashH(r'))
\enspace,
\]
\fi
\fi

$\mathsf{FO}$ is used pervasively in post-quantum cryptography. We argue that $\mathsf{OAEP}$ is of interest in this setting  as well. Particularly, we will show how to use $\mathsf{OAEP}$ to obtain bandwidth improvements over using  $\mathsf{FO}$ for code-based cryptography.  This will significantly broaden the applicability of $\mathsf{OAEP}$, as well as motivate the study of its post-quantum instantiability. Our analyses here will be in the quantum ROM, \emph{i.e.}, where the adversary's RO queries can be in superposition.

 $\mathsf{OAEP}$ was classically applied to one-way trapdoor \emph{permutations} (TDPs) namely RSA. For post-quantum cryptography, our idea is to analyze $\mathsf{OAEP}$ applied to trapdoor \emph{functions}, towards the aim of getting. The main TDFs of interest are from \emph{code-based cryptography}, particularly the Neiderriter TDF~\cite{} used in Classic McElice~\cite{}. Also note that  the  McElice randomized PKE scheme~\cite{} is  randomness-recovering (\emph{i.e.}, the decryptor recovers the  coins used to encrypt) and  hence can be viewed as a TDF. 
 Our main observation is 
 \begin{quote}
\emph{Upgrading code-based TDFs to IND-CCA secure PKE via $\mathsf{OAEP}$, rather than $\mathsf{FO}$, as  done \emph{e.g.}~in classic McElice, would yield bandwidth gains.}
\end{quote}
Namely, the gains are because (part of) the message would be encrypted ``inside'' the TDF.
Such an optimization is not possible when using $\mathsf{FO}$ because the TDF  is one-way and would not protect partial information about the message.

Unfortunately, prior work fails to establish the resulting PKE schemes would be secure.
Specifically, while Targhi and Unruh~\cite{}  proves $\mathsf{OAEP}$ secure for TDFs in the quantum ROM, they modify $\mathsf{OAEP}$ in a way that makes it less  efficient. NTRU adopts this modification~\cite{}. 
Later work proves the unmodified $\mathsf{OAEP}$ secure in the quantum ROM,  but they only consider TDPs. Thus, we ask:
\begin{ques}
    What properties of a TDF $\calF$ (\emph{e.g.}, Neiderriter) make $\mathsf{OAEP}$ applied to $\calF$ secure in the quantum ROM?
\end{ques} 
We will consider both IND-CCA and IND-CPA as the target security notion. When encrypting short messages, we can encrypt the entire message inside the TDF, so the appropriate target would be IND-CCA.  When encrypting long messages, the appropriate target  would be IND-CPA because we can upgrade the scheme to a bandwidth-optimized IND-CCA secure hybrid encryption. We plan to study how exactly to do this. Roughly, encryption will  have the form:
\[
\calE^{\hyb}_{pk}(m_1\|m_2;r) = 
f(s \| t) \| \calE^{\sy}_{K}(m_2) 
\quad \mbox{where} \quad
  s = \hashG(r) \xor  m_1 \| 0^\zeta, \quad t = \hashH(s) \xor r, \quad K = \hashG(s \| t)
\enspace.
\]
Note that this step is not a direct application of the second part of the $\mathsf{FO}$ transform, so we will give a separate analysis. 

%Furthermore, for long messages we should  not require partial one-wayness~\cite{C:FOPS01} of the TDF rather than standard one-wayness.

%Additionally, Many post-quantum randomized PKE schemes such as Kyber are converted to TDFs as part of $\mathsf{FO}$ by using Encrypt-with-Hash~\cite{C:BelBolONe07}.We can again save bandwidth by  encrypting part of the message ``inside'' this TDF, which significantly broadens the scope of our proposed work.  In particular, to achieve IND-CCA security here, for short messages we  need to apply $\mathsf{OAEP}$, though not for long messages because the PKE scheme is assumed IND-CPA. We will look at applications where one encrypts short messages and work out the concrete bandwidth savings for various schemes for different message-lengths and security levels.

%One issue is that there is an attack under IND-CCA  if given $f,y$ one can efficiently find $y' \notin \mathsf{Image}(f)$ such that $f^{-1}(y') = f^{-1}(y)$. (This can't happen when using a TDP.) We will deal with this in similar ways to  $\mathsf{FO}$. We will   calculate the  bandwidth savings of using $\mathsf{OAEP}$ instead of $\mathsf{FO}$ for specific schemes like Classic McElice.


%To obtain a post-quantum instantiation of FO, we will aim to construct \emph{post-quantum}  ELFs. Recall that ELFs are currently only known from exponential-DDH, which is quantum insecure, and exponential assumption are \emph{necessary} to construction ELFs. In particular, we will look at constructions from post-quantum iO plus forms of the learning parity with noise (LPN) assumption that are plausibly post-quantum, exponentially-secure.  In particular, we will consider LPN with constant noise and linear number of samples.  Recent work constructs (trapdoor, but non-extreme) lossy functions from related assumptions~\cite{}, lending plausibility to the approach. Another assumption we will consider is the recently introduced MQ assumption~\cite{}. Even a heuristic construction would be of interest at this point, given the stark lack of progress on this question.


%A primary concern about idealized models is when they can be \emph{instantiated}, meaning when we can safely replace the of computation that is abstracted away by the model by public, efficiently-computable code.  The alarming fact is that there are many \emph{un}instantiability results, \emph{e.g.}~\cite{}, showing schemes secure in popular idealized models may have \emph{no such instantiation at all}, let alone a canonical instantiation like $\mathsf{SHA256}$ in the case of a RO. It is thus of great interest for both practice and theory \emph{which} idealized-model schemes are instantiable and \emph{how} to instantiate them.

%in algebraic settings. 
%We discuss potential applications of our notions we will explore, and they will serve as a backdrop for our new assumptions in the following thrust. 

%Building on prior success in the RO model, we undertake a broad approach to   instantiability:

\iffalse
\begin{ques} \label{ques-std}
   Is there  a unified framework for instantiating many idealized models across  a range of applications of interest?
\end{ques}
\fi

\iffalse
\subsection{Instantiating Idealized Models the UCE Way}
 Here we are interested in taking a broader approach  and developing a unified framework for instantiating a \emph{range of  idealized models}. Namely, we will develop a framework to that capture what it means to ``behave like $X$'' where $X$ is an idealized object like a random oracle or a generic group. % We start with  the background for our main approach.

\heading{Possible approaches.}
There are many notions geared toward instantiating the ROM  such as extremely lossy functions~\cite{C:Zhandry16}, correlation-intractable hash functions~\cite{JACM:CanGolHal04}, \emph{etc.} It may be possible to generalize such notions  to other idealized models, \emph{e.g.}~one might  define extremely lossy groups or correlation-intractable groups. We take a different approach based on \emph{universal computational extractors} (UCEs)~\cite{C:BelHoaKee13}, which prior work already shows is promising in this regard. Broadly, we ask:
%\paragraph*{Moving  to other idealized models.} There is currently not much work on notions designed instantiate the GGM and AGM across a range of schemes. 
%In fact, we aim for a broader agenda:
\begin{ques}
 How far can we push UCE  to instantiate idealized models beyond the ROM?
\end{ques}

We stress that our goal here is \emph{definitional} in nature. We will later sketch a few applications, such as instantiating the $\mathsf{ECIES}$ encryption scheme~\cite{}.

\heading{Instantiating random oracles with UCEs.}
To motivate  UCEs, it is instructive to see why a simple attempt at instantiating a RO in general fails (which it must, due to uninstantiability results). Namely, let us  try to define what it means for a hash function to look ``like a RO.''  Fix a keyed hash function $\Hash \colon K \times D \to R$. Define  the following game: Sample a challenge bit $b \getsr \bin$, a hash key $\hk\getsr K$, and a RO $\rho \colon D \to R$.  Then, adversary $\advA$  runs with access to an oracle $\HashO$ which, when queried on $x \in D$, returns $\Hash(\hk,x)$ if $b =1$ and $\rho(x)$ otherwise. $\advA$'s goal is to guess $b$.  But this is easy: $\advA$ on input $\hk$ queries $\HashO(x)$ for arbitrary $x \in D$ and returns 1 iff $\Hash(\hk,x)$. 

A solution this issue was proposed by Bellare \emph{et al.}~\cite{C:BelHoaKee13,C:BelHoaKee14},  with  their  notion of \emph{universal computational extractors} (UCEs). 
 The idea  is to split the adversary $\advA$ into two stages called ``source'' $\advs$ and ``distinguisher'' $\advd$.
 As above, let $\Hash \colon K \times D \to R$ be a keyed hash function. UCE is defined via the following game: Sample a challenge bit $b \getsr \bin$, a hash key $\hk \getsr K$, and a RO $\rho \colon D \to R$. Then, $\advs$ runs with access to an oracle $\HashO$ which, when queried on $x \in D$, returns $\Hash(\hk, x)$ if $b = 1$, and $\rho(x)$ otherwise. Eventually, $\advs$ outputs leakage $L$, which is passed to $\advd$, who is also given the hash key $\hk$ but no oracle access. Distinguisher~$\advd$ then tries to guess $b$. 
 
%While there are many of other notions geared towards instantiating ROs~\cite{}, we particularly  focus on building off of  UCEs.  Indeed, in  advanced cryptographic schemes, what often makes proving security in the standard model difficult is that correlated and leaky information about the secrets is available to the adversary. This concern is exactly what UCE addresses.
To avoid ``trivial'' attacks, one must put restrictions on  $\advs$.  Bellare \emph{et al.}~define  \emph{unpredictability} (UP). meaning the following game should be hard to win:  $\advs$ runs with access to a RO and produces leakage~$L$. Then $\advp$ runs on input $L$ and wins if it guesses any of $\advs$'s queries. 
 Thus,  intuitively, UCE instantiates a RO  when honest inputs are unpredictable but  may be correlated or leaky. 
 
% \heading{What's the point?} The point is we can assume \emph{e.g.}~$\mathsf{SHA256}$ is a UCE whereas assuming it's a RO makes no sense; there are supporting theoretical constructions of UCEs in special cases~\cite{}. Indeed, when we are using $\mathsf{SHA256}$ we want to rigorously understand what we are assuming about it, so we may study such properties. Independently of  applications, we believe the definition UCE and the various extensions of it we develop have intrinsic philosophical value. 
 % This means for UP sources there is no  generic attack.  %In fact, in certain cases UCEs have been built using iO~\cite{}, showing they are instantiable.

%An interesting question, which will become more complex in other idealized models we consider below, is how dow we know that unpredictability is enough to rule out all ``trvial'' ottacks and what does trivial mean exactly?  For UCEs, it means that the notion is achieved \emph{by a RO itself}.  


%Although we primarily view this as an important definitional question in its own right, we give  a few applications as well.% For example, we will treat novel notions of subversion-resistant non-interactive zero-knowledge proofs as well a and instantiations as the popular ECIES encryption scheme.

%By later incorporating \emph{knowledge extraction} into the framework, we can facilitate moving results in the  AGM~\cite{} to the standard model, \emph{e.g.}, about SNARKs~\cite{} and multisignatures~\cite{}. %If we are successful in Section~\ref{xxx}, we will have that these frameworks are \emph{instantiable}. 
%Finally, in applications requiring ``proofs about proofs,'' \emph{i.e.}, where a party must present a proof that a given proof verifies, one could use UCE or its algebraic extensions to avoid proving statements about circuits containing RO or generic group gates, which we cannot do in a provable way.


\iffalse
\paragraph*{Prior work.}
 Soni and Tessaro (ST)~\cite{EC:SonTes17} previously defined an ideal primitive and implementation, with a UCE-inspired notion of \emph{public-seed pseudorandomness} (psPR).  Namely, an ``ideal primitive'' is a pair  $\mathbf{I} = (\Sigma,\calT)$ where $\Sigma$ is a set of functions having the same domain and range and $\calT$ is a (not-necessarily-uniform) probability distribution on $\Sigma$.  The corresponding ideal model provides all parties with oracle access to a function $f$ chosen according to  $\calT$. A function family $\FF \colon \keys \times \Dom \to \Rng$ is an \emph{implementation} of $\mathbf{I}$ if $\FF(K,\cdot) \in \Sigma$ for all $K \in \keys$.  For $\mathbf{I}, \FF$, and adversary $\advA = (\advs,\advd)$, ST define the prSP game that chooses $b \gets \bits$ and $K\getsr \keys$, runs $\advs$ with oracle access to either  $\FF(K,\cdot)$ (if $b=1)$ or $f$ sampled according to $\calT$ (if $b =0$), to get leakage $L$, which is given to $\advd$ along with $K$; $\advd$ then has to guess $b$. A particular case of psPR that ST consider are  
\emph{public-seed pseudorandom permutations} (psPRPs), which extends UCE to permutations.

While our  definitions could  be stated in the psPR framework,  ST do not consider algebraic idealized models or idealized cryptographic primitives.  Moreover, a treatment of PSP in such cases does not  follow from their work, because there is no formula for specifying the  ideal primitive and correspondingly what will be the key for $\FF$ (\emph{i.e.}, play the role of the hash key in UCE).  There is similarly no formula for specifying how, given some $\mathbf{I}, \FF$, the source should be restricted so that the resulting notion is  both (1) achievable and (2) useful. We anticipate these  as being key challenges. 
\fi 
%While we could still use ST's framework for our definitions, we opt not to do so for readability.  

%In particular, they show the restrictions that must be imposed on the adversary must change.  More generally, they define the notion of an \emph{ideal primitive}. However, their notion does not cover groups, and it does not seem possible to give a unified treatment of all ideal primitives under UCE-style notions, as there is no general recipe for how to restrict the source against an arbitrary ideal primitive. 

%\begin{task}
%In each setting below, formalize UCE-style notions and identify (1) source classes for which they are achievable and (2) applications of the corresponding notion.
%\end{task}
%While we focus mainly on the definitional aspect, we envision that our notions will serve to ``explain'' some  \emph{ad hoc} assumptions and yield new ones useful in cutting-edge cryptography.

\iffalse
\paragraph*{Gold-standard instantiations.}
For an idealized-model scheme, there is typically a \emph{canonical} instantiation, such as using $\mathsf{SHA256}$ in place of a RO. To justify the canonical instantiation, we want to show it is \emph{gold-standard} in the following sense. We want to identify a property  $\mathsf{P}$ or multiple properties $\mathsf{P}_1,\ldots,\mathsf{P}_n$  that are interesting in their down right and are (1) plausibly satisfied by object used for the canonical instantiation, \emph{e.g.}, $\mathsf{SHA256}$ or a suitable elliptic curve group,  (2)   sufficient to prove security of $\mathsf{S}$, and (3) achievable in principle based on plausible cryptographic assumptions. Note the construction in (3) may be impractical and is not meant to be used in practice. Rather its  \emph{existence} gives plausibility to security of the canonical instantiation used in practice. %An example, the Fiat-Shamir~\cite{} transform for some $\Sigma$-protocols can be instantiated using correlation intractable (CI) hash functions~\cite{}, $\mathsf{SHA256}$ is plausibly CI, and CI hash functions can be built in principle from iO~\cite{}. 

Not all  instantiations are gold-standard.  For example, ELFs~\cite{C:Zhandry16} cannot be achieved from ROs~\cite{TCC:GunFis23}, so presumably $\mathsf{SHA256}$ is not an ELF. Also, RO instantiations using iO~\cite{}   typically build custom hash functions such that when they are plugged into $\mathsf{S}$ yield the desired security.  
\fi


 % The in RO model, Schnorr has a loose reduction to DL~\cite{} in the standard model, a tight reduction to DL in the AGM~\cite{}, and a tight reduction to a non-standard but standard-model assumption~\cite{}. On the other hand, Chen \emph{et al.}~\cite{} show that in the Shoup's GGM~\cite{EC:Shoup97} one can use an extremely simple hash function securely, albeit for small message spaces.  (We stress that the idea is \emph{not} to build the hash function using the GGM.) 

 %Another key aspect of our approach will concern new notion of \emph{programmability}~\cite{} of either the group or hash functions.   We will also  look for ways around programmability that are not precluded by negative results~\cite{EC:FisFle13}.


%Besides UCE,  we will port other notions from hash function to groups, such as programmability~\cite{}, correlation-intractability~\cite{}, extreme lossiness~\cite{}, and more, giving us a diverse toolkit to work with.
%This will be especially useful for studying the interplay between properties of groups and properties of hash functions needed in the protocols.
%We note that some properties of groups we use may \emph{not} satisfied by appropriate elliptic curve groups.  Rather, we ``manufacture'' groups with some of the properties from indistinguishability obfuscation and other assumptions.  Of course, we do not suggest to use these manufactured groups in practice.  However, their \emph{existence} makes it more likely appropriate elliptic curve groups would still ``work'' in the protocols because the protocols are instantiable in principle. 
%For applications, the idea will be to leverage such assumptions on the group so as to mitigate requirements on other components of the protocols in question, such as hash functions. 
%To see this, note that an idealized model can be defined via  a \emph{labeling oracle} $\LabelO$ that maps elements of a set $S$ under consideration to random handles, as well as additional oracles $\calO_1,\ldots, \calO_n$  that may call $\LabelO$. In an instantiation, we require  that a (keyed) public, efficiently-computable function  replaces  $\LabelO$ wrt.~``split'' adversaries as in UCE.  Thus, a natural question is:  %As the adversary can call $\LabelO$ so could seemingly simulate the additional oracles itself, but actually the game may make \emph{unpredictable} queries that the adversary \emph{cannot} simulate.  Our main question is:

% We recall the UCE notion and then our preliminary work. 

%To understand the proposed work  it will be necessary to understand my preliminary work in TCC 2022~\cite{}. We therefore give an overview below, starting with PGG's predecessor UCE.

%Here we describe our recently published preliminary work on the main problem at hand.  We especially point out the key insights from the work that we will take forward.


\subsubsection{Pseudogeneric Groups} \label{sec-prime}

\iffalse
The case of prime-order groups was previously examined by my prior work. Our notion, \emph{pseudo-generic groups} (PGGs), is one way of porting UCE to the GGM.  However, here we view PGG as part of a different framework.
Namely, inspired by the wide-ranging  applicability and interest in algebraic PRFs~\cite{}, we will study algebraic hash functions and ask:
\begin{ques}
    Can we build algebraic UCEs?
\end{ques}
Zhandry shows ELFs, which can be realized from exponential-DDH,  yield UCEs for one-query computationally unpredictable sources~\cite{}. While it achieves a very weak version of UCE, this gives evidence such algebraic hash functions exist.
To make the question more precise, fix a group parameter-generation algorithm~$\group$ which, on input the security parameter~$\secparam$, outputs group parameters~$\params$ consisting of a group operation~$\cdot$, an arbitrary group generator~$g$, and a prime~$p \in [2^{\secpar-1}, 2^\secpar)$. Implicit in these parameters is a set~$\GG$ such that~$(\GG, \cdot)$ forms a cyclic group of order~$p$ with generator~$g \in \GG$. We write the sampling of group parameters as~$(\params := (\cdot, g, p)) \getr \group(\secparam)$, with the understanding that~$\params$ implicitly defines the underlying set~$\GG$. Now let $\Hash \colon K(\params) \times D \to \GG$. We want $\Hash$ that are UCE. 
%Our approach here to building algebraic UCEs is more general as compared to the PI's prior work; see below.

We will in particular investigate if known algebraic PRFs (or modifications thereof) are UCEs. This is motivated by the fact that the  Dodis-Yampolskiy (DY) PRF~\cite{} is precisely the construction of UCEs from groups proposed in~\cite{},although it was not viewed that way.  However, it was only proved UCE for low-degree sources.  In the proposed work, we will aim to treat the more general case, considering other constructions. Even going beyond groups, we will also examine \emph{post-quantum} UCEs from lattices, \emph{e.g.}, based on SWIFFT~\cite{} or LWE-based PRFs~\cite{}.
%As a first step, we will examine whether the Goldreich-Goldwasser-Micali PRF is a \emph{correlated-input hash} (CIH)~\cite{} if the underlying PRG is a RO. CIH is a relaxation of UCE introduced by the PI where the leakage consists only of hash values.  This is handy as it implies CIH can hold for \emph{computationally} unpredictable sources, which UCE cannot~\cite{}.
\fi

%\begin{ques}
%    How can we formulate UCE in the case of prime-order groups?
%\end{ques}
% We start with prime-order groups. 
We start with the case of prime-order groups and the GGM~\cite{EC:Shoup97}, as treated in the PI's work at TCC 2022. 
Here  all algorithms have oracle access to a random generic-group encoding  $\sigma \colon \Z_p \to \GG$ as well as an oracle $\Op$ that on inputs $h_1,h_2 \in \GG$ returns $\sigma(\sigma^{-1}(h_1) + \sigma^{-1}(h_2) \bmod p)$.  When run, an algorithm get $\sigma(1)$ and the encodings of other application-dependent elements.
For our instantiations, it will important to consider not a fixed group $\GG$ but  
a group parameter-generator~$\group$ which, on input~$\secparam$, outputs~$\params$ consisting of group  $(\GG,\cdot)$,  a generator~$g$, and prime~$p \in [2^{\secpar-1}, 2^\secpar)$. Namely,  $(\GG,\cdot)$ forms a group of order~$p$ with generator~$g$. We write~$(\params := (\GG, g, p)) \getsr \group(\secparam)$. %Now, we ask:
% with the understanding that~$\params$ implicitly defines the underlying set~$\GG$. We let other parameters depend on $\GG$. % Let $\Hash \colon K(\params) \times D \to \GG$. 
%One possibility is to ask that $\Hash$ be UCE. However, we will take a  different angle, offering a new perspective on the PGG notion introduced in~\cite{TCC:BFHO22}.

%\begin{ques}
%In a UCE-style definition for groups, what should play the analogue of the hash key?
%\end{ques}


%We  consider two different answers, giving rise to two  corresponding notions. 

%\paragraph*{Pseudo-generic groups, Take 1.} 
The idea is having a random group generator play the role of the hash key in UCE. For $\params := (\GG, g_0, p)$, consider the exponentiation hash function $H_{\mathsf{exp}} \colon \GG \times \Z_p \to \GG$ where  $H_{\mathsf{exp}}(g,a)  := g^a$ for $a \in \Z_p$.   It is easy to see that $H_{\mathsf{exp}}$ is \emph{not} UCE. We are interested in  restrictions on the source such that the notion is achieved by $H_{\mathsf{exp}}$.  We call $\advs$'s oracle $\Exp$ here, as it either computes exponentiation (if $b=1$) or a random mapping (if $b=0$).

%The PI's prior work proved that generic groups \emph{even with preprocessing} (aka.~AI-GGM) are PGG for masking sources. It also show PGG for masking sources  generalizes the Uber assumption~\cite{EC:BonBoyGoh05} and allows achieving advanced security notions via ElGamal-style constructs. 
%Sample a secret bit $b \sample \bin$, $(\params := (\circ, g, p)) \sample \group(\secparam)$, a random generator $h \sample \GG$, and a generic-group encoding $\encode \colon \ZZ_p \to \GG$. Then, $\advs$ given $\params$ interacts with an exponentiation oracle $\Exp$ which, on input $x \in \ZZ_p$, returns the real group element $h^x$ if $b = 1$, and a generic element $\encode(x)$ otherwise. The source can pass some leakage~$L$ to $\advd$, who is also given the random generator $h$ but loses access to the oracle and has to guess $b$.

\iffalse
\paragraph{Pseudo-generic groups, Take 2.}
Here we let the entire group parameters (except its order) play the role of the hash key  in UCE.  Let us call it PGG2, now referring to the prior notion as PGG1.  %A usual, we  split the adversary $\advA$ into two stages, a source $\advs$ and distinguisher $\advd$. 
The PGG2 game samples a secret bit $b \getsr \bin$, group parameters~$(\params := (\GG, g, p)) \getsr \group(\secparam)$ and generic-group encoding $\encode \colon \ZZ_p \to \GG$.  Then, $\advs$ given $p$ interacts with a group operation oracle $\Op$ that on input $h_1,h_2 \in \GG$ returns $h_1 \cdot h_2$ if $b=1$ and $\encode(\encode^{-1}(h_1) + \encode^{-1}(h_2))$ otherwise, as well as an exponentiation oracle $\Exp$ which, on input $x \in \ZZ_p$, returns $g^x$ if $b =1$ and $\encode(x)$ if $b=0$.  The source can pass some leakage $L$ to $\advd$, who is also given $\params$ but loses the oracles, and tries to guess $b$.  %We will later see that it is easier to construct groups satisfying PGG2, motivating its consideration.
\fi 

\heading{Restrictions on the source.}
%\begin{ques}
%What are necessary or sufficient conditions on the source for PGG1 (r
Unlike UCE, UP is not sufficient for PGG to be achievable.
An insight of the PI's prior work~\cite{TCC:BFHO22} was to require a notion of \emph{algebraic unpredictability} (AU).
In the AU game, the source $\advs$ runs with access to the ideal oracle while querying $x_1, \ldots, x_q$, and produces leakage~$L$. The unbounded predictor $\advp$ runs on input $L$ and must guess a linear combination of the queries, \emph{i.e.} it outputs $(\alpha_0, \alpha_1, \ldots, \alpha_q)$ not all zero and wins if $\sum_{i=1}^q \alpha_i x_i = \alpha_0$. 
Beyond AU, we  required  parallelizability and limited post-processing conditions.
 %  To see why, consider $\advs$ that on input $\params$ chooses $x \getsr \ZZ_p$, submits $h \gets \Exp(x)$, computes $h' \gets h^{1/x}$, and returns $L = h'$; then, $\advd$ on input $(g,L)$ returns $(g=L)$. 
%In particular~\cite{TCC:BFHO22} shows that a generic group is PGG1 for masking sources.  In the proposed work, we will consider different  source classes, for example  \emph{reset sources}~\cite{C:BelHoaKee13}, as well as limited adaptivity among queries. Here, a challenging task will be:
%\begin{task}
%Characterize \emph{necessary} restrictions on $\advs$ for PGG1/PGG2 to be achievable in the GGM.
%\end{task}
%an query $\Exp$ adaptively.  %For example, we will explore novel source classes such relaxations of block-sources. Block sources played a major role in the related study of  deterministic encryption~\cite{C:BolFehONe08}, of which UCE is a (trapdoorless) generalization.  
%For PGG2,  considering that $\advs$ only performs its group operations via $\Op$, we may be able to restrict these queries to drop AU, making PGG1 and PGG2 incomparable in general. % This will hopefully open up very different applications. %The question is how useful the resulting notion will be in applications. To appreciate the usefulness of PGG2, consider that in most applications of the Uber assumption, the specific polynomials used in the assumption are fixed by the scheme and not chosen by the adversary as a function of the group parameters. PGG2 will suffice in such applications. 


%An important goal will be to determine when PGG1 vs.~PGG2 is useful in applications and for what class of sources.  We conjecture that, like PGG1, PGG2 for masking sources implies the Uber assumption and holds in the AI-GGM. 
%\paragraph*{Applications.}
%The PI's prior work shows PGG1 for masking sources generalizes the Uber assumption~\cite{EC:BonBoyGoh05,PAIRING:Boyen08} so gives all implications thereof, addresses the need to move beyond this assumption, and it also implies advanced security properties for ElGamal like key-dependent message (KDM)-security, related-key attack (RKA)-security, \emph{etc.} We observe that PGG2 for masking sources also  generalizes the Uber assumption as long as the polynomials do not depend on the group parameters. In many applications this is the case, as the polynomials are fixed by the scheme in question. In the proposed work, we will comprehensively study when PGG1 vs.~PGG2 is applicable. One benefit of PGG2 is that the oracles can use $g_0$ rather than $g$, which can be fixed given $\GG$. %Note that one  benefit of PGG2 is that it captures \emph{fixed-generator} assumptions~\cite{C:BarMaZha19}. % We will aim to determine when PGG1 vs.~PGG2 are applicable.

%Based on our preliminary investigation below, we believe necessary conditions for PGG1 could be (1) limited post-processing (\emph{e.g.}, $\advs$ cannot exponentiate its query answers) and (2) there does not exist a multilinear polynomial $f \in \Z_p[X_1,\ldots,X_q]$ that vanishes on $\advs$'s queries with good probability. 

%To prevent trivial attacks, one needs more than unpredictability of the $\advs$'s queries to $\Exp$, as seen by the following: the
% source $\advs$ samples $x_1, x_2 \sample \ZZ_p$, queries $h_i \gets \Exp(x_i)$, and computes $x_3 \gets x_1 + x_2$ and $h_3' \gets h_1 \circ h_2$. It then queries $h_3 \gets \Exp(x_3)$ and passes the bit $(h_3 = h_3')$ to $\advd$, who simply returns it. The advantage of such  $\advA = (\advs, \advd)$ in the PGG game is almost~$1$, even though $\advs$ is unpredictable since $x_1$ and $x_2$ are random.
%In particular, a core concep
%or UCEs, the answer is unpredictable sources. 
%In our context, unpredictability alone is not sufficient:
%The restrictions on $\advs$'s queries to $\Exp$ will be the same as for PGG2 below, previously introduced in the PI's TCC 2022 paper.  In particular, a core conceptual contribution is a notion \emph{algebraic unpredictability} (AU).
%In the AU game, the source $\advs$ runs with access to the ideal exponentiation oracle while querying $x_1, \ldots, x_q$, and produces leakage~$L$. Predictor $\advp$ runs on input $L$ and must guess a linear combination of the queries, i.e., outputs $(\alpha_0, \alpha_1, \ldots, \alpha_q)$ not all zero and wins if $\sum_{i=1}^q \alpha_i x_i = \alpha_0$.  Beyond AU, as in~\cite{} we impose parallelizability and limited post-processing restrictions that we do not discuss in detail here. 

\iffalse
In the case of PGG1 (as opposed to PGG2 below) further restrictions are needed: Consider $\advs$ that samples $x_1,x_2 \sample \Z_p$ and queries $g_i \gets \Exp(x_i)$ for $i \in \bits$. It then queries $g_3 \gets \Op(g_1,g_2)$ and passes $L = (g_1,g_2,g_3)$ to $\advd$, who returns 1 iff $g_1 \cdot g_2 = g_3$.  At first blush, this could be addressed by requiring (basic) unpredictability of the queries to $\Op$.  However, this is tantamount to requiring unpredictability of the \emph{responses} to $\advs$'s queries to $\Exp$. Looking ahead, this would severely limit potential applications. Instead, we require unpredictability of responses to $\advs$'s $\Op$ queries.   
\fi 

%As usual, the group operation gives rise to an exponentiation algorithm~$\mathsf{exp}(h, x)$ whose output is denoted as~$h^x$. We will often omit explicitly writing the operation~$\circ$.
%To port the UCE definition to cyclic groups, our first idea is to let a random group generator $g$ play the role of the hash key, and to use exponentiation with base $g$
% (i.e., the map $x \mapsto g^{x}$)
%in place of the hash.


\iffalse
\paragraph*{Pseudo-generic Groups: Take 2.}
In PGG1, $\advs$'s queries to $\Exp$ cannot depend on $\params$, which limits  applications.   Next we define another way of adapting UCE to the GGM, which we call PGG2. (Let us now call the original PGG notion PGG1 for clarity.)  PGG2 can be seen as applying UCE not to the GGM itself but the exponentiation function keyed by the generator. Compared to PGG1, we drop the $\Op$ oracle.  The PGG2 security game for a group parameters-generator~$\params$  is as follows: Sample a secret bit $b \sample \bin$, $(\params := (\cdot, g, p)) \sample \group(\secparam)$, a random generator $h \sample \GG$, and a generic-group encoding $\encode \colon \ZZ_p \to \GG$. Then, $\advs$ given $\params$ interacts with an exponentiation oracle $\Exp$ which, on input $x \in \ZZ_p$, returns the real group element $h^x$ if $b = 1$, and a generic element $\encode(x)$ otherwise. The source can pass some leakage~$L$ to $\advd$, who is also given the random generator $h$ but loses access to the oracle and has to guess $b$.

We keep the same restrictions on $\advs$'s queries to $\Exp$ as in PGG1, but we drop the additional requirement of unpredictability of responses from $\Op$.  As such, we believe that PGG1 and PGG2 are \emph{incomparable}.  Indeed, PGG1 is weaker than PGG2 in that the source's queries to $\Exp$ do not depend on the group parameters. However, PGG2 is weaker than PGG1 in that it does not require the group operation to be idealized.

\fi
%We want that every such $\advA = (\advs, \advd)$ has a small advantage in this game.
%Thus, the PGG notion captures the intuition that if an adversary does not know the random generator $g$ of $\GG$, exponentiation with base $g$ looks like it returns random elements from $\GG$.  

%To prevent trivial attacks,  one needs more than unpredictability as seen by the following: the
 %source $\advs$ samples $x_1, x_2 \sample \ZZ_p$, queries $h_i \gets \Exp(x_i)$, and computes $x_3 \gets x_1 + x_2$ and $h_3' \gets h_1 \circ h_2$. It then queries $h_3 \gets \Exp(x_3)$ and passes the bit $(h_3 = h_3')$ to $\advd$, who simply returns it. The advantage of such  $\advA = (\advs, \advd)$ in the PGG game is almost~$1$, even though $\advs$ is unpredictable since $x_1$ and $x_2$ are random.
%In particular, a core concep
%or UCEs, the answer is unpredictable sources. 
%In our context, unpredictability alone is not sufficient:
%Accordingly, a core conceptual contribution of our work was \emph{algebraic unpredictability} (AU).
%In the AU game, the source $\advs$ runs with access to the ideal exponentiation oracle while querying $x_1, \ldots, x_q$, and produces leakage~$L$. Predictor $\advp$ runs on input $L$ and must guess a linear combination of the queries, i.e., outputs $(\alpha_0, \alpha_1, \ldots, \alpha_q)$ not all zero and wins if $\sum_{i=1}^q \alpha_i x_i = \alpha_0$.  Beyond AU, we impose parallelizability and limited post-processing restrictions that we do not detail here.  Subject to these restrictions, we prove a generic group is PGG, meaning there are no other trivial attacks.
 

\subsubsection{Pseudogeneric Group Actions}


\iffalse
Beyond prime-order groups, we will consider a number of algebraic settings.  We only discuss two of them below, but we may also consider hidden-order and composite-order groups.
%As idealized models are common in a variety of algebraic settings, towards realizing them we investigate the PGG notion beyond basic group structure. 

\paragraph*{Bilinear Groups.}
We will develop notions of  \emph{bilinear} PGG1 (bPGG1) and \emph{bilinear} PGG2 (bPGG2). For simplicity, let us consider bPPG1 for bilinear-group parameters $\params := (\GG, \GG_T, g_0, p,e)$ with non-degenerate symmetric bilinear map  $e \colon \GG \times \GG \to \GG_T$. In the preliminary definition, we will give the source access to \emph{two} oracles, which are either exponentiation in $\GG,\GG_T$ relative to unknown random generators $g,g_T$ respectively $(b=1)$ or independent random encoding functions $(b=0)$.
\iffalse
%  We sketch our preliminary definition of BPGG1.
For clarity, let us stick to symmetric pairings and just fix a bilinear group described by a tuple $(e,\GG,\GG_T,p)$ where $e$ is  an efficiently computable, non-degenerate bilinear map $e \colon \GG \times \GG \to \GG_T$ for groups $\GG,\GG_T$ of prime order $p$. Here let a pair $(g,g_T)$ of random  generators for $\GG, \GG_T$ respectively play the role of the hash key in the UCE definition.  That is, we provide the adversary oracle access to \emph{two} exponentiation oracles with bases $g,g_T$ respectively.  In the bilinear PGG game, we first sample a secret bit $b \sample \bin$,  random generators $g \sample \GG$ and $g_T \sample \GG_T$, and generic-group encodings $\encode \colon \ZZ_p \to \GG$ and  $\encode_T \colon \ZZ_p \to \GG_T$. Then, a source $\advs$ interacts with exponentiation oracles $\Exp,\Exp_T$. The first, on input $x \in \ZZ_p$, returns $g^x$ if $b = 1$, and  $\encode(x)$ otherwise.  Similarly, the second, on input $x \in \ZZ_p$, returns the real group element $g_T^x$ if $b = 1$, and a generic element $\encode_T(x)$ otherwise. The source can pass some leakage~$L$ to a distinguisher $\advd$, who is also given $g,g_T$ but loses access to the oracles and has to guess $b$.  
 \fi
Building on AU, we will consider \emph{quadratic  unpredictability} (QU).
For QU, the source $\advs$ runs with access to the ideal exponentiation oracles $\encode,\encode_T$ while querying $x_1, \ldots, x_q$ to $\encode$ and $x_{T,1}, \ldots, x_{T,q}$, and produces leakage~$L$. The unbounded predictor $\advp$ runs on input $L$ and outputs 
\begin{newmath}
(\alpha_0,\alpha_1,\ldots,\alpha_q, \beta_{1,1}, \ldots, \beta_{q,q}, \alpha_{1,T}, \ldots, \alpha_{q,T})
\end{newmath}%
not all zero and wins if 
\begin{newmath}
\sum_{i=1}^q \alpha_i x_i + \sum_{j,k=1}^{q} \beta_{j,k} x_j x_k + \sum_{\ell=1}^q \alpha_{\ell,T} x_{\ell,T}  = \alpha_0 \;.
\end{newmath}%
After identifying various source classes to be considered, we will endeavor to establish analogues to the results above for prime-order groups. In particular, here we will use our notions to study  advanced security properties of cutting-edge encryption schemes in bilinear groups, like registered inner-product encryption~\cite{DBLP:conf/asiacrypt/FrancatiFMMRV23} and registered attribute-based encryption~\cite{EC:HLWW23}. %We envision most of the applications with come with the extensions below.
\fi
  %  \end{task}
 %   Comparatively, the PI's prior work only showed such a result for ElGamal.

%We envision that key applications will come when the models are extended below. 
%\begin{ques}
%      What are necessary and sufficient conditions on the source for bPGG1 (resp.~bPGG2) to be achievable in 5the  bilinear AI-GGM?   
%end{ques}

%Beyond QU,  we will study analogous questions to the case of non-bilinear groups above.
%To start with, we  aim to show analogous results to the non-bilinear case, namely bPGG1 and bPGG2 for masking sources imply  the bilinear Uber-assumption and that they hold in the bilinear AI-GGM. 


%\paragraph*{Capturing $\mathsf{KOALA}$.} $\mathsf{KOALA}$~\cite{} is a novel assumption useful in cutting-edge cryptography~\cite{}. We aim to show $\mathsf{KOALA}$ is captured by multi-source, knowledge-based  BPGG1 or BPGG2. 
%It is easier to explain the ``weak'' $\mathsf{KOALA}$ assumption also given in~\cite{}.  Consider an algorithm $\mathsf{Samp}$ that outputs $\vecu_1,\ldots,\vecu_t \in \Z_p^m$ where the $\vecu$'s are correlated individually pseudorandom.  Then $$
%To understand $\mathsf{KOALA}$, consider the matrix Diffie-Hellman (MDDH) assumption~\cite{C:EHKRV13} on a  (bilinear) group $\GG$ of order $p$, which for some $\ell > k \in \N$ and matrix $M \in \Z_p^{\ell \times k}$ states that $(g,g^{Mx})$ is indistinguishable from $(g,g^u)$ where $x \getsr \Z_p^k$ and $u \getsr \Z_p^\ell$. $\mathsf{KOALA}$ interprets MDDH as a knowledge  assumption, stating that if a distinguisher $\advd$ breaks MDDH, $\advd$ must ``know'' a vector in the nullspace of $M$.  Agrawal, Shota, and Wichs~\cite{}  consider distributions on $M$, hence they give auxiliary input to $\advd$ and the extractor the view of $\advA$.

%By giving a broader framework for $\mathsf{KOALA}$, we hope to find weaker assumptions that suffice for the above applications.  We note that optimal broadcast encryption has recently been constructed for evasive LWE~\cite{}, but it seems difficult to give a  framework for that assumption.
 
 
\iffalse
Moreover, we believe 
there are vastly different restrictions that also work (both in the prime order and bilinear cases), and we plan to identify them and find applications in the proposed research.  For example, instead of using parallel sources, one might use \emph{block sources}; cf.~my work on deterministic encryption (DE)~\cite{C:BolFehONe08}. DE is closely related to UCE as the latter generalizes correlated-input hashing, which in turn is just DE without decryption.  Thus, we believe there will be applications to PGG as well. 
\fi


\iffalse
\paragraph*{Establishing plausibility.} In our preliminary work, we prove that \emph{a generic group is PGG}. While this is a promising result, we view it as imperative to more establish soundness of PGG. Indeed,  in making assumptions that test the limits of standard-model provability, one has to be extremely careful about impossibility/uninstantiability results. To address this concern, I propose several  approaches.

\begin{newitemize}
\item First, I will define the \emph{Augmented GGM} (AGGM) akin to the Augmented RO model (AROM)~\cite{}.  The AROM is a way of a showing RO-model construction resists techniques for proving uninstantiability; the AGGM will be such a way for GGM-based constructions. In particular, we wish to prove that a \emph{generic group is PGG in the AGGM}. This will significantly build on the proof in the preliminary work.
\item
Second, we will seek constructions of  PGG groups based on  accepted assumptions (say a DL-hard group,  program obfuscation,  homomorphic encryption, and NIZKs).  While such constructions would be not practical, they would  also refute uninstantiability of PGG under these assumptions.  Such an approach was taken  by Agrikola and Hofheinz~\cite{}
for Uber-hard groups and by Agrikola \emph{et al.}~\cite{} for the AGM. 
\end{newitemize}
%We next describe some ways we that we will generalize the PGG notion.
\fi

%\paragraph*{Adding adaptivity.} A major obstacle to applying PGG in certain applications is that the achieved security notion is \emph{non-adaptive}.  The core reason is that in PGG the source must make \emph{all} queries before seeing information about generator $g$.  The analogous issue was address in the UCE setting my Farshim and Mittelbach~\cite{FSE:FarMit??}. We endeavor to lift their techniques to the group-theoretic setting...
%We anticipate bPGG1 and bPGG2 will be useful in revisiting security of schemes that are proven in the bilinear GGM, for example the structure-preserving signature of~\cite{JC:FucHanSla19}.
%We believe the most compelling applications of bPGG1 and bPGG2 will be made possible by extensions to these notions below.

%\paragraph*{Composite order.}  Our idea for extending bPGG1 and bPGG2 to composite-order is motivated by the generalized subgroup decision assumption~\cite{TCC:BelWatYil11}.  In the composite-order setting, the subgroup decision assumption~\cite{TCC:BonGohNis05,TCC:BelWatYil11}   implies ``$q$-type''  assumptions~\cite{EC:ChaMei14}. We would like to study whether this phenomenon generalizes to cases where exponents are correlated or leaky.


%\paragraph*{Hidden order.}
%In the hidden-order setting, we no longer need \emph{limited post-processing} for the PGG source (because the adversary exploiting post-processing is  inefficient), which may open up novel applications.


%\begin{task}
%    Generalize the results of~\cite{EC:ChaMei14} to the PGG context (leaky, correlated exponents). 
%\end{task}

%\paragraph*{Hidden-order groups.} 
%Hidden-order groups present an intriguing setting for PGG because the adversary that wins the game using post-query exponentiation is no longer efficient.  This means that in  knowledge-based versions of PGG in this setting, we may simply ask that for a successful adversary must ``know'' a multilinear polynomial vanishing at its queries.% This style of definition is akin to $\mathsf{KOALA}$.  %Hidden-order groups are important in  constructions of verifiable delay functions and succinct arguments. %n particular, we will study applications of our notions to a component of these known as \emph{proofs-of-exponentiation}~\cite{}. 

\heading{Defining PGGA.} Efficient (commutative) group actions such as CSIDH~\cite{AC:CLMPR18} are of growing importance for post-quantum cryptography.  Zhandry~\cite{ITCS:Zhandry24}, building on Duman \emph{et al.}~\cite{AC:DHKKLR22}, introduced an idealized model called the \emph{generic group action model} (GGAM). % In fact, Zhandry~\cite{} argues that the GGAM is preferable to the AGM in this setting.  
Remarkably, a UCE-style definition is a perfect fit for instantiating the GGAM  because in the GGAM the adversary effectively just has a group exponentiation oracle; there is no group operation oracle. So, we can simply drop the $\Op$ oracle in PGG and obtain a notion we call \emph{pseudo-generic group actions} (PGGA). A major benefit is that this notion will have much simpler restrictions on the source than PGGs for the notion to be achievable, and thus broader applicability. %%We want to show a generic group action satisfies PGGA. Then, to enable broad protocol design as well as to prove negative results:
%Furthermore, we will consider the notion wrt.~\emph{quantum} adversaries:
%\begin{task}
%   Formulate and prove an Uber assumption for group actions that admits a broad class of decision oracles. First prove it in the GGAM and then under PGGA, against \emph{quantum} adversaries.
%\end{task}
%Decision oracles are important in applications~\cite{AC:DHKKLR22}.   Then we aim to find quantum money/lightning schemes under other  assumptions than~\cite{ITCS:Zhandry24} using our framework.


\heading{A construction of PGGA from obfuscation.} Here we will go further than simply defining PGGA and giving applications, and give a supporting theoretical construction using indistinguishability obfuscation (iO)~\cite{DBLP:journals/jacm/BarakGIRSVY12,FOCs:GGHRSW13,JACM:JaiLinSai24} and puncturable pseudorandom functions~\cite{STOC:SahWat14}. % We describe our  construction using a puncturable pseudorandom \emph{injection}~\cite{TCC:BonKimWu17}. However, our construction can be adapted to suitable puncturable PRFs. 
    Let $\calO$ be an iO and  $F \colon \bits^k \times \bits^\lambda \to \bits^{q(\lambda)}$ be a \emph{puncturable pseudorandom injection} (PRI)~\cite{TCC:BonKimWu17} (\emph{i.e.}, which is invertible given the key) for some polynomial $q$. Our group action construction uses two programs with hardcoded  key $K \in \bits^k$ and $\lambda$-bit prime $p$.
 \vspace{-1mm}
 \begin{center}
 \begin{tabular}{c|c} 
 \begin{minipage}{1in}
\begin{tabbing}
123\=123\=123\=123\=123\=\kill
\textbf{Program} $\mathsf{P}_0[K, p](L_0,L_1)$: \\
\> If $F^{-1}_K(L_b) \in \Z_p$ for $b\in \bits$ then \\
\>\> $L \gets F_K(F^{-1}_K(L_0) + F^{-1}_K(L_1) \bmod p)$ \\
\>\> Return $L$ \\
\> Else return $\bot$
\end{tabbing} \end{minipage}
 &
\begin{minipage}{1in}
\begin{tabbing}
123\=123\=123\=123\=123\=\kill
\textbf{Program} $\mathsf{P}_1[K, p](L)$: \\
\> If $F^{-1}_K(L) \in \Z_p$ return $1$ \\
\> Else return $0$\\\\
\end{tabbing} \end{minipage}
\end{tabular}
\end{center}

\noindent The associated group scheme  $\Gamma_0[\calO,F] = (\Gamma_0.\Setup,\Gamma_0.\Val,\Gamma_0.\Add)$ is defined as follows:
\begin{newitemize}
    \item On input $\secparam$,  $\Gamma_0.\Setup$ chooses $K \getsr \bits^k$ and $\lambda$-bit prime $p$, and sets  $\widehat{\mathsf{P}}_b \getsr \calO(\mathsf{P}_b[K,b])$ for $b \in \bits$. Its output  is  $((\widehat{\mathsf{P}}_0, \widehat{\mathsf{P}}_1),F(1),p)$.
    \item On inputs $(\widehat{\mathsf{P}}_0,  \widehat{\mathsf{P}}_1),h$,  $\Gamma_0.\Val$ outputs $\widehat{\mathsf{P}}_1(h)$.
        \item On inputs $(\widehat{\mathsf{P}}_0,  \widehat{\mathsf{P}}_1),h_0,h_1$,  $\Gamma_0.\Add$ outputs $\widehat{\mathsf{P}}_0(h_0,h_1)$.
\end{newitemize}

%As $F_K$ is an \emph{injection},  $\Gamma_0[\calO,F]$ enjoys unique representation.
%For this task, the PI will collaborate with a new faculty member at UMass Amherst in quantum computation,  Filip Rozpedek. 

% with the hope of identifying novel assumptions and applications, \emph{e.g.} related to quantum money~\cite{ITCS:Zhandry24}. 
%Additionally, we will consider \emph{preprocessing attacks}. Note that in the case of DL, the optimal preprocessing attack in groups~\cite{ITC:RotSeg22} simply performs a random walk that should also work with group actions as well. For an upper bound, since group actions are just groups with less structure, existing techniques~\cite{EC:CorKog18} should again apply. However, it is unclear if existing techniques apply for more complex assumptions like PGGA2.


\iffalse 

 \subsection{Extensions  to the Algebraic Notions and Applications} \label{sec-know}

%  We next discuss extensions to our basic  framework to capture. Broadly, we want to capture  cutting-edge assumptions under a common framework and  address instantiability of cutting-edge schemes. %, which we will  prove are satisfied by generic groups.

% We will also consider all these approaches for a wide range of CE-analogues. 

%We stress that in UCE, the source $\advs$ either talks to the real hash function (if $b=1$) or a random oracle oracle (if $b=0)$.  However, in PGG, the source does \emph{not} talk to  the real group or a random group.  Rather, one may think of it as talking either \emph{exponentiation} or a \emph{random labeling function} $\sigma$.   Recall that if talking to  a random group the adversary gets an oracle that on labels $x,y$ returns $\sigma(\sigma^{-1}(x) \circ \sigma^{-1}(y))$. Falling more inline with the UCE case, it is tempting to define a  version of PGG  that replaces $\sigma^{-1} (\cdot)$ here with the discrete log function in the real group when the PGG challenge bit $b=1$, but then the game would be inefficient.  In other words, \emph{a key insight  is that in porting UCE to other idealized models, we must consider in which ``aspecct'' of the model we wish to instantiate.}
%Nevertheless, jumping ahead, we will explore PGG variants in this direction. 
%As before, for this notion not to be void we must put restrictions on the queries that $\advs$ is allowed to make. First, observe that $\advs$ must be unpredictable, because without any such requirement $(\advs, \advd)$ can mount the attack for UCEs sketched above.
%Indeed, similarly to the UCE definition, unpredictability prevents trivial attacks such as where $\advs$ queries $1$ and $2$ to get $g$ and $h$, respectively, then computes $h' \gets g \circ g$, and returns $1$ if and only if $h' = h$. (Notice that $1$ and $2$ are predictable.
%We argue now that due to the presence of a group structure on $\GG$ that was missing in the UCE setting, further and quite delicate conditions are needed. 

%For applications, it will be useful to further extend prime-order PGG1/PGG2 along a few directions. The extensions mostly apply to our later notions in other algebraic settings as well.

%For further applications, it will be useful to extend the above notions in a few ways. 
%We extend the notions a few ways needed in applications.=
 
\paragraph*{Multiple sources.}
We will consider multiple sources  $\advs_1, \ldots, \advs_n$ who each  output a leakage and field element pair $(L_i,x_i)$, the $x_i$'s then interacting via some polynomials and passed to $\advd$ with all the leakages.  (There is no longer an $\Exp$ oracle.) For example, in the case of PGG1 for two sources, fix bivariate polynomials $p_1,p_2,\ldots,p_m$ and consider sources $\advs_1,\advs_2$ where $\advs_1$ on input $\params$ outputs $(L_1,x_2)$ and $\advs_2$ on input $\params$ outputs $(L_2,r_2)$.  Then  $\advd$ receives $(L_1,L_2,g)$ for random $g \getsr \GG$ as well as either $g^{p_i(x_1,x_2)}$  for all $i \in [m]$, or $m$  random group elements.
%z\begin{ques}
%      What are necessary and sufficient conditions on the sources for multi-source bPGG1 to be achievable in the  bilinear AI-GGM?   
%\end{ques}
Surprisingly, this generalization  captures \emph{fixed-generator} assumptions (cf.~\cite{C:BarMaZha19}).  For example, taking $p_i(x_1,x_2) = x_1 x_2 + x_1^i$ and $L=x_2$ we get Assumption 3 of~\cite{C:BarMaZha19}. (AU follows from the game-based Schwartz-Zippel Lemma in~\cite{TCC:BFHO22}.) %The benefit is it now follows from the techniques of the PI's prior work~\cite{TCC:BFHO22} that the assumption holds in the GGM. More generally, we ask: 
%We will further study whether such fixed-key notions are possible for UCE/psPRP and, if so, develop applications to using a fixed-key blockcipher.  
%\begin{ques}
% Which existing  assumptions are captured by multi-source PGG notions?
%\end{ques}
We can also capture \emph{Matrix Diffie-Hellman} (MDDH)~\cite{JC:EHKRV17}.  %wherein a  non-uniform matrix of exponents is randomized in the exponent. % $\calD$-MDDH asks that for  \smash{$\mathbf{A} \in \Z_p^{\ell_1 \times \ell_2}$} with $\ell_1 > \ell_2$ sampled from a distribution $\calD$ and $\vecx \getsr \Z_p^{\ell_2}$, $(g,g^{\mathbf{A}}, g^{\mathbf{Ax}})$ is indistinguishable from $(g,g^{\mathbf{A}},g^{\vecu})$ where \smash{$\vecu \getsr \Z_p^{\ell_2}$}.  By capturing such assumptions, we will be able to analyze them in a unified way and find weaker assumptions that suffice. 

\iffalse
\begin{ques}
    Can we define a \emph{search} version of our PGG-style notions?
\end{ques}


There is \emph{search} version of MDDH called  \emph{kernel MDDH}~\cite{AC:MorRafVil16}, which asks that is it hard to find (in the exponent) a vector in the kernel of $\mathbf{A}^\top$. This  suggests a search version of our PGG-style notions that replaces $\advd$ with a ``finder'' $\advf$.  Assume for simplicity that the leakage $L$ is just the vector of responses $L = (Y_1,\ldots,Y_q)$ that $\advs$ receives to its $\Exp$ queries, which corresponds to an algebraic analogue of \emph{correlated-input hashing}~\cite{TCC:GoyOneRao11}, a relaxation of UCE. Then it should be hard for $\advf$ given $(g,L)$ to  output $a_1,\ldots,a_q \in \Z_p$ such that $\prod_{i=1}^n Y_i^{a_i} = 1_{\GG}$.
\fi

% \begin{ques}
%     Can we extend our PGG-style notions to take \emph{knowledge extraction} into account?
% \end{ques}


\paragraph*{Adaptivity.} UCE/PGG are \emph{non-adaptive}; \emph{e.g.}, for PGG1  the source's queries must be independent of the generator $g$. This will hinder us in capturing ``flexible'' assumptions where the adversary's target to output can depend on the problem instance. For  UCE, Farshim and Mittelbach~\cite{FSE:FarMit16} developed an adaptive extension called \emph{interactive computational extractors} (ICEs).  We will refine ICEs and extend them to  algebraic settings. 
\iffalse
In tandem, we will study:
\begin{ques}
    Can we develop a \emph{search} (vs. decisional) version of PGG1/PGG2 and extensions?
\end{ques}
\fi
In particular, we want to capture %$\emph{e.g.},  the Kernel %Matrix DH  Assumption~\cite{AC:MorRafVil16}, $q$-SDH~\cite{JC:BonBoy08} t
 $q$-SDH~\cite{JC:BonBoy08} and its extension to ARSDH~\cite{EC:LipParSii24},   only  justified in the AGM. Note that these are \emph{search} assumptions, so we may use search relaxations of our framework in such cases.


\paragraph*{Knowledge.} To enable AGM instantiations, we will add a \emph{knowledge} (aka.~extractability) component to our framework. 
%(With some caveats~\cite{AC:ZhaZhoKat22,EPRINT:JaeMoh24}, a proof in the AGM implies a proof in the GGM, but PGG1/PGG2 do not instantiate all applications thereof.)
%We propose to synthesize these cases into a  notion that answers our question above. 
%Roughly, a knowledge assumption asserts that an adversary producing some group elements or distinguishing two distributions of them must ``know'' some discrete logs. % For example, the \emph{knowledge of exponent assumption} (KEA)~\cite{C:Damgaard91,C:BelPal04} asserts that given $g,g^a$, an adversary that produces $(C,Y)$ such that $C^{a} = Y$ must ``know`` the discrete-log of $C$. This is formulated via existence of a suitable non-blackbox extractor $\calE$.
%Indeed, at first glance, it may seem we can turn \emph{any} assumption in the AGM into a standard-model one by requiring  a corresponding extractor $\calE$ outputs an explanation of the adversary's output. For example, \emph{knowledge-CDH} (kCDH) would say:  For any $\advA$ that on inputs $g,g^{x},g^y\in \GG$ outputs $g^{xy}$, there is a non-blackbox extractor $\calE$ that given the view of $\advA$ outputs $a,b,c \in \Z_p$ such that $a + bx + cy = xy$.  The problem is that kCDH reduces to DL in the \emph{standard model}, following the same reduction as in the AGM~\cite{C:FucKilLos18}. 
%However, such a conversion from an assumption in the AGM (or its decisional version~\cite{TCC:RotSeg20}) to one in the standard model \emph{is} useful when there are both hard and easy instances of a problem.  For example, 
 %Our starting point will be $\mathsf{KOALA}$~\cite{PKC:BeuWee19,TCC:AgrWicYam20},   which requires that a successful  $\calD$-MDDH distinguisher as above ``knows'' a vector in the kernel of $\mathbf{A}^\top$. 
% and $\mathsf{KOALA}$ says the \emph{only} way it can be easy is if the distinguisher knows such a vector. 
%Interestingly, in $\mathsf{KOALA}$ the extractor does \emph{not} get the view of the distinguisher. We will consider both possibilities in that regard.
Developing a  framework for  knowledge assumptions was posed as an open problem by Ghadafi and Groth~\cite{AC:GhaGro17}.  We start with special cases to build up intuition.
Let us start with the case  that the  the leakage $L = (Y_1, \ldots, Y_q)$ passed to $\advd$ in the PGG1 game is simply the vector of responses given to $\advs$ on its $\Exp$-queries. This is an algebraic analogue of correlated-input hashing~\cite{TCC:GoyOneRao11}. As the knowledge component,   we require that for a successful adversary $\advA = (\advs,\advd)$,  $\advd$ ``knows''  $a_1,\ldots,a_q \in \Z_p$ such that $\prod_{i=1}^q Y_i^{a_i} = 1_{\GG}$. Interestingly, we do not need any restrictions on $\advs$ here. Next, consider a more general case that  $L = (Z_1,\ldots,Z_{q'})$ for some $q' \in \N$ is a vector of $\emph{arbitrary}$ group elements. Here we also require that  $\advs$ ``knows'' an explanation of each $Z_i$, \emph{i.e.}, some $a_{i,1}, \ldots, a_{i,q} \in \Z_p$ such that  $Z_i =  \prod_{j =1}^{q} Y_j^{a_{i,j}}$.

Such restricted $L$ should suffice in some applications.
However, we would like to treat the general case that $L$  contains arbitrary \emph{bits}, which is much more challenging to reason about.
Nevertheless, we may still be able to make assumptions about the adversary's knowledge in this case. Our main idea is to require that for a successful adversary $\advA = (\advs,\advd)$, $\advs$ ``knows'' an $f \in \Z_p[X_1,\ldots,X_q]$ that vanishes on its $\Exp$-queries. For example, consider  $\advs$ that queries $x \getsr \Z_p$ to get $X$ and $x^2$ to get $Y$, then passes $L = 1$ if $X^x = Y$ and $L=0$ otherwise, and $\advd$ outputs  $L$. Then $f(X_1,X_2) := X_1^2 - X_2$.

\begin{ques}
    Can we synthesize the above cases into a   framework for  knowledge assumptions?   
    \end{ques}


 % Note we \emph{can} allow computational AU sources in the GGM~\cite{TCC:BFHO22}.
%For implications, we will start with KEA1 and KEA3~\cite{C:Damgaard91,C:BelPal04}, and the  $q$-power KEA~\cite{AC:Groth10a}. We also
%In some applications, we will need to  consider ``benign'' but \emph{computational} AU sources (\emph{i.e.},  auxiliary input). ``Benign'' means immune to obfuscation-based attacks~\cite{C:BrzFarMit14,STOC:BCPR14,AC:BoyPas15}.
Results in the AGM, unlike the GGM, are \emph{reductions}. We will study whether our  framework reduces to a ``nice'' assumption in the (decisional) AGM~\cite{C:FucKilLos18,TCC:RotSeg20} and thus instantiates applications that  fit a specific game structure, where we assume knowledge  extractors rather than the adversary outputting explanations. We conjecture the impossibility of~\cite{C:Zhandry22b} will not apply and  further investigate feasibility in Section~\ref{sec-know}. 
Assumptions we  aim to generalize  include $\mathsf{KOALA}$~\cite{PKC:BeuWee19,TCC:AgrWicYam20}, a knowledge version of MDDH  used in  breakthroughs~\cite{EC:AgrYam20,TCC:AgrWicYam20, C:AgrYadYam22}. 
   Note we have no  framework for evasive LWE~\cite{EC:Wee22}, which is used in similar applications.  Such a framework facilitates cryptanalysis as well as constructions from weaker assumptions.
   
%We will then show extractability makes UCE-style notions applicable to signature schemes for the first time. We will focus on recent  \emph{two-round multisignature schemes with key aggregation} (cf.~\cite[Tables 2 and 3]{EC:PanWag24}), useful in blockchains. The analyses of the most popular schemes such as $\mathsf{HMBS}$~\cite{AC:BelDai21} and $\mathsf{MuSig2}$~\cite{C:NicRufSeu21} require the AGM for tight security (in the RO model). We want to instantiate the AGM here. The trend in the community has been to design \emph{alternative} schemes with tight such proofs, \emph{e.g.}, $\mathsf{Toothpick}$~\cite{EC:PanWag24}.  However, we do not expect practitioners to move away from the most popular schemes. Therefore, we need better analyses of the latter. 
   % Then we will study instantiability of schemes set in the GGM or AGM, such as  some ``real world'' zkSKARKs~\cite{cryptoeprint:2024/721}  and efficient witness encryption schemes~\cite{PKC:CamFioKho24,cryptoeprint:2024/264}.  Note that there is some controversy about the AGM~\cite{C:Zhandry22a}.  The challenge will be to fit such applications into the PGG template.
 
Recently, Bauer \emph{et al.}~\cite{EPRINT:BFHK23}  proposed a framework for knowledge assumptions they call the ``Uber-knowledge'' (UK) framework. We believe the UK framework   captures somewhat different assumptions than ours. 
We will conduct a detailed comparison as part of the proposed work. 
\fi

%(like $\mathsf{KOALA}$) is decisional. %and in their framework the ``verification equation'' is restricted to be a one-dimensional polynomial. %We will conduct a detailed comparison as part of the proposed work. %Rin ours the leakage $L$ can contain arbitrary bits, which is important for capturing assumptions where the adversary's output includes field elements.
\iffalse
\paragraph*{Applications to randomness subversion.} We will leverage our framework to study  \emph{randomness subversion}, where randomness is  controlled by Big Brother, for \emph{succinct non-interactive arguments of knowledge}  (SNARKs). %We first explain the direction and then the relevance of our framework. %Indeed, we will see that a broad framework for knowledge assumptions will be needed to analyze security in this setting.
Several SNARKs  assume a \emph{structured reference string} (SRS) that is honestly generated.   A subverted SRS was studied  by~\cite{PKC:Fuchsbauer18,JC:ALSZ21}, following~\cite{AC:BelFucSca16} in the non-succinct case. In particular, they define \emph{subversion-zero-knowledge} (sub-ZK) and \emph{subversion-soundness} (sub-SND), where the adversary  chooses the SRS. But sub-SND can't be achieved because the adversary knows the trapdoor, and the only current-generation SNARK analyzed under sub-ZK  is $\mathsf{Groth16}$~\cite{EC:Groth16} (by~\cite{PKC:Fuchsbauer18}). Moreover, to show sub-ZK, they  require the prover to perform expensive ``consistency checks'' on the SRS. In practice,  such checks are not performed, and there is a ``proof-of-concept'' attack~\cite{CCS:Fuchsbauer19}. We aim to broaden and understand such attacks:   %We aim to both broaden scope to other popular zkSNARKs and better understand the attack-space:
\begin{task}
    Characterize the space of sub-ZK attacks possible on  $\mathsf{Groth16}$~\cite{EC:Groth16}, $\mathsf{Plonk}$~\cite{EPRINT:GabWilCio19}, and $\mathsf{Marlin}$~\cite{EC:CHMMVW20}  without  consistency checks (or very minimal checks).
\end{task}


We will start with KZG commitments~\cite{AC:ChuKatZho13}.
We will also  consider a relaxation where the adversary  has  \emph{partial}  control over the SRS, \emph{i.e.}, where the SRS is  \emph{unpredictable} (UP). This could happen in a ceremony such as~\cite{AC:KMSV21} when no one is honest. %SNARKs resilient  to such attacks could also allow lighter-weight such ceremonies, cf.~\cite{C:GraVadZuc06} in the case of coin-tossing. 
Here, we hope to show positive results for sub-ZK and sub-SND.
Finally,  suppose the \emph{prover's} randomness is bad but the witness is UP. In particular, we will consider ``transparent''  SNARKs (STARKs) such as~\cite{C:GLSTW23} here. Such settings are related to foundational work on ZK, \emph{e.g.}~\cite{FOCS:CanPasShe07,EPRINT:DahLin20,TCC:BitCho20}, but   new for SNARKs.

Several aspects of our framework will be useful in avoiding idealized models here. For example, sub-ZK  involves knowledge assumptions because the simulator needs to extract a trapdoor from an adversarially-chosen SRS. Moreover, when the SRS or witness is UP, exponents may be correlated and leaky; UP and algebraic UP can be bridged by \emph{linear-dependence destroyers}~\cite{TCC:BFHO22}.  For STARKs, we will consider a novel combination of UCE plus \emph{extractability}~\cite{TCC:CanDak09}.
\fi
  %For SND, there is  \emph{knowledge-soundness} (KSND) and the stronger and more practical \emph{simulation-extractability} (SIM-EXT). Recent  work  shows KSND of  $\mathsf{Groth16}$~\cite{EC:Groth16} in the UK framework~\cite{EPRINT:BFHK23} and  KSND of $\mathsf{Plonk}$~\cite{EPRINT:GabWilCio19} and $\mathsf{Marlin}$~\cite{EC:CHMMVW20} in the RO model~\cite{cryptoeprint:2024/994}. However, for SIM-EXT, currently such results require  the AGM (with oblivious sampling). So:
\iffalse
\begin{task}
Show  SIM-EXT of $\mathsf{Groth16}$,  $\mathsf{Plonk}$, and $\mathsf{Marlin}$  without relying on the AGM.
\end{task}
Actually, $\mathsf{Groth16}$ can only achieve ``weak'' SIM-EXT.
We would like to show the above \emph{with efficiency optimizations (such as ''relinearlization'') used in practice}.
\fi


%We will further study instantiability of the RO model for zkSNARKs and related protocols like Bulletproofs~\cite{SP:BBBPWM19}, where it is used for the Fiat-Shamir (FS) transform~\cite{C:FiaSha86}.  Not much is known about instantiating 

%A ``dream result'' would be to instantiate Fiat-Shamir~\cite{C:FiaSha86} for $\mathsf{Plonk}$ and $\mathsf{Marlin}$, even in the standard case of KSND.  This  \emph{e.g.}~by building on the above and techniques developed to instantiate Schnorr signatures in Section~\ref{sec-fs}


 % Second, we will use framework to instantiate popular threshold and multisignature schemes analyzed in the AGM; see Section~\ref{sec-thresh}. %(In particular, without capturing knowledge, PGG does not seem to apply to zkSNARKs or signatures.)
 %Third, we will aim to instantiate recent witness encryption schemes analyzed in the bilinear GGM~\cite{PKC:CamFioKho24,cryptoeprint:2024/264}.%Finally, we will study recent threshold signatures and multisignatures that for efficiency (either by virtue of a tight reduction or the construction itself) are set in the AGM.  % Note that the AGM has come under controversy beyod merely being instantiable~\cite{AC:ZhaZhoKat22,C:Zhandry22b}, making transitioning the above results to the standard model all the more important. % special-case witness encryption~\cite{PKC:CamFioKho24,cryptoeprint:2024/264}. 
% \adam{Need applications} % We will conduct a detailed comparison as part of the proposed work. %We will perform a detailed comparison as part of the proposed work.  % For example, they do not capture $\mathsf{KOALA}$.  To be fair, their goal was not to capture decisional assumptions, but we believe they do not even  

%We will finally synthesize these cases into a generation notion of \emph{knowledge} bPGG1/bPGG2 with multiple sources.% We want to (1) justify it in the bilinear GGM and (2) use it to analyze SNARKs, a primary application of the AGM:  %Note that this approach largely obviates the need to specify explicit restrictions on the source. %Then we will use it to eliminate use of the AGM in recent SNARKs:

\iffalse
\paragraph*{Application to zkSNARKs.}
For the most efficient versions of recent zkSNARKs that are actually deployed, the analyses often rely on the AGM/GGM. We ask:  %Indeed, recent research dropping these idealized models incurs an efficiency losses~\cite{}.

\begin{ques}
    Can we avoid using the AGM/GGM in zkSNARKs without an efficiency loss?
\end{ques}

Recent research suggests their use could be replaced with non-falsifiable/knowledge assumptions~\cite{} or even falsifiable assumptions~\cite{}.  Indeed, the latter does not contradict the seminal Gentry-Wichs result~\cite{} as these SNARKs use the RO model~\cite{C:FiaSha87}.  
Ut is a good topic to study which “non-falsifiable assumptions” one can use to get optimal efficiency

\fi

%Here, some examples are my TCC paper (not cited yet by you - we proposed AGM with oblivious sampling, AGMOS, there), BFHK23.

%Second, mention Gentry-Wichs. And that GW is not applicable in the ROM (it only states about the plain model). Most practical SNARKs (except Groth16 and its precursors) use ROM. In that case, it is possible you can just use falsifiable assumptions. Here, LPS24 is a good citation. (LPS24 has efficiency loss but I have a new paper, hopefully eprinted the next week, where we show Plonk is secure in the ROM under falsifiable assumptions.)

%\paragraph*{Applications to zkSNARKs.} 
%There are some related very interesting questions. Like can you do sim-ext under falsifiable assumptions, without losing efficiency. I am working on it (just started this week, currently alone.) Recently Faonio, Fiore, Russo (eprint 2024, accepted to ACM CCS) showed Plonk is sim-ext in the AGM with oblivious sampling

%A challenge will be when $L$ contains both group elements and bitstrings, which is already tricky in the AGM itself~\cite{C:Zhandry22b}.
%\begin{task}
%    Reduce kbPGG1/kbPPG2 to entropic $q$-DL in the bilinear AGM, or show this is not possible. 
%\end{task}

%\begin{task} \label{task-agm}
%    Prove knowledge-soundness of \emph{e.g.}~Groth16~\cite{EC:Groth16}, Sonic~\cite{CCS:MBKM19}, Plonk~\cite{EPRINT:GabWilCio19}, Marlin~\cite{EC:CHMMVW20}, or variants based on kbPGG1/kbPGG2 for multiple sources.
%\end{task}

%Proving knowledge-soundness of SNARKs without resorting to the AGM is a well-established line of work.  Recently,  Bauer \emph{at al.}~\cite{EPRINT:BFHK23} and Lipmaa \emph{et al.}~\cite{EC:LipParSii24} provide new assumptions/frameworks for doing so. Comparatively, we will aim to make our framework easy to use and applicable to diverse schemes.   Note that as the application to SNARKs does not require \emph{decisional} assumptions,  we may develop search versions of our assumptions here instead. % The upshot of our results will be that kbPGG1/kbPGG2 are likely not impossible (whereas the AGM is uninstantiable~\cite{C:Zhandry22b}) as we will show in Section~\ref{sec-feas}.
%By generalizing it, we hope to find weaker assumptions than $\mathsf{KOALA}$ that suffice in its applications to optimal broadcast encryption~\cite{} and multi-input attribute-based encryption~\cite{}.
%One curious issue is that in $\mathsf{KOALA}$ the extractor does \emph{not} get the view of $\advd$, which would be natural and consistent with the decisional AGM~\cite{TCC:RotSeg20}.

% Our rough idea is to ask that $\emph{if}$ there is an adversary winning the PGG game \emph{then} there is a non=blackbox extractor for it that outputs $a_1,\ldots,a_q$ as above.  
 %Note that one cannot simply make the entire AGM a standard model assumption, in particular because nested extractors have a blow-up in runtime. However, we will investigate whether key applications of the AGM can be moved to our standard-model framework: 
%\begin{ques}
 %   Are SNARKs (\emph{e.g.}, Groth16~\cite{C:FucKilLos18}, Sonic~\cite{}, Plonk~\cite{}, Marlin~\cite{}) or multisignatures (\emph{e.g.}, ) provable under knowledge PGG?
%\end{ques}

%\begin{ques}
%    Are multisignatures currently proven secure in the AGM (\emph{e.g.}, ) provable in our standard-model framework?
%\end{ques}

%Some of these will be in the bilinear setting.
%Actually, considering (b)PGG1 and (b)PGG2 in the \emph{decisional} AGM is not necessary for SNARKs.  We will therefore develop a \emph{search} version of our framework. As a first step, it we will treat algebraic analogues of the PI's notion of \emph{correlated-input hashing}~\cite{TCC:GoyOneRao11}, the reason being that  the leakage $L$ will only con sist of group elements. (Handling adversaries who receive both group elements and bitstrings in the AGM is trick~\cite{C:Zhandry22b}.)  %wheras 


%As another application, we will try to capture the $\mathsf{KOALA}$ assumption of~\cite{PKC:BeuWee19,TCC:AgrWicYam20}.  $\mathsf{KOALA}$ extends MDDH to ask that a successful distinguisher must ``know'' a vector in the nullspace of $\mathbsf{A}$.  $\mathsf{KOALA}$ has applications to  obfuscation~\cite{PKC:BeuWee19},  broadcast encryption~\cite{EC:AgrYam20,TCC:AgrWicYam20}, and attribute based encryption~\cite{}.  By capturing it. we aim to study of whether weaker or alternative assumptions suffice. There are subtleties to resolve, for example, in $\mathsf{KOALA}$, the extractor does \emph{not} get the adversary's view, which differs from usual knowledge assumptions. 

% A second way is to extract knowledge from the source instead of the distinguisher.  Namely, we ask that for every successful $A = (\advs,\advd)$  there is a non-blackbox extractor $\mathcal{E}$ that outputs a  multilinear polynomial $f \in \Z_p[X_1,\ldots,X_q]$  vanishing on $\advs$'s queries.  For example, consider $\advs$ that on input $\params$ chooses $x \getsr \Z_p$ and queries $x$ and $x^2$ to get $h_1,h_2$, then passes $L=1$ if $h_1^x = h_2$ and $L=0$ otherwise; $\advd$ just outputs $L$. Then $\mathcal{E}$ outputs $f[X_1,X_2] := X_1^2 - X_2$.
 

% There is \emph{search} version of MDDH called  \emph{kernel MDDH}~\cite{AC:MorRafVil16}, which asks that is it hard to find (in the exponent) a vector in the kernel of $\mathbf{A}^\top$. This  suggests a search version of our PGG-style notions that replaces $\advd$ with a ``finder'' $\advf$.  Assume for simplicity that the leakage $L$ is just the vector of responses $L = (Y_1,\ldots,Y_q)$ that $\advs$ receives to its $\Exp$ queries, which corresponds to an algebraic analogue of \emph{correlated-input hashing}~\cite{TCC:GoyOneRao11}, a relaxation of UCE. Then it should be hard for $\advf$ given $(g,L)$ to  output $a_1,\ldots,a_q \in \Z_p$ such that $\prod_{i=1}^n Y_i^{a_i} = 1_{\GG}$.

\iffalse
 \subsection{UCE-Style Security Notions for Encryption}

\heading{Motivation.} Here we will  consider UCE-style definitions to instantiate   ideal authenticated encryption (AE)~\cite{C:BarFar18} and ideal pubic-key encryption (PKE)~\cite{C:ZhaZha20}. Note that AE and PKE are ``higher level primitives'' as compared to hash functions and groups. Therefore, the  motivation for these idealized models is different from the previous ones, namely for ideal AE and PKE the motivation is to construct schemes meeting \emph{all} ``game-based'' security notions. Intuitively, this means one can security deploy schemes that are ``like'' the idealized objects in these models in any scenario whose security is captured by a game, which is a very useful property.  % Indeed, when  deploying a scheme, it is difficult to predict all applications in which it will be used, and therefore all security properties needed, ahead of time. 

Prior work interprets ``like'' via  indifferentiability~\cite{TCC:MauRenHol04}, which they achieve via constructions using lower-level idealized models like the ROM.  This approach, however, has a deficiency in that while the \emph{ideal object} in the ideal AE or ideal PKE models satisfies all game-based  security notions, the indifferentiable  constructions may not. The  issue is that indifferentiability only holds for ``single-stage'' games~\cite{EC:RisShaShr11} (where the adversary is stateful), and many important security notions are multi-stage. As UCE is multi-stage, by defining UCE-style security notions  for AE and PKE, we can prove indifferentiable schemes  meet many \emph{multi-stage} security notions as well, which is a strict improvement over prior results. This agenda parallels the work of~\cite{C:BelHoaKee14} on UCE.

%\begin{ques}
%What  classes of \emph{multi-stage} security notions can we achieve for AE and PKE?  
%\end{ques}
%\begin{ques}
%What notions of multi-stage indifferentiabilty can be achieved in the standard model?
%\end{ques} 

%To achieve  multi-stage indifferentiability, one can go  for reset-indifferentiability~\cite{EC:RisShaShr11,AC:Zhandry21}, which  requires a stateless simulator. 
%However, it seems hard to avoid the state here. Following~\cite{C:BelHoaKee14},  we propose to work in the standard model and use analogues of UCE instead.
%Broadly, we will consider both idealized and standard-model notions of multi-stage security.

%In particular prior work on ideal PKE only considered indifferentiable constructions in the RO model; we will consider constructions in both the RO model and GGM.  %, and showed it implies multi-stage indifferentiability. %,  and unsplittability~\cite{EC:Mittelbach14}, which applies to any multi-stage game that can be viewed as ``keyed hash function game'' where only the last stage gets the key.
%However, the simulators constructed for the schemes of~\cite{C:BarFar18,C:ZhaZha20} \emph{do} maintain state. Will study whether \emph{weak} reset-indifferentiabilty~\cite{AC:Zhandry21} can be achieved; this notion suffices and can be easier to construct. In particular,  indistinguishability against query-unbounded distinguishers implies weak indistinguishability\cite{AC:Zhandry21}.
%\begin{ques}
%Are there indistinguishable (aka.~pseudorandom) AE or PKE schemes against query-unbounded adversaries, or more generally weak reset-indifferentiable AE or PKE schemes?
%\end{ques}
%We  stress that such constructions need not be the same as the ones from~\cite{C:BarFar18,C:ZhaZha20}
 %Instead, we propose to use extensions of UCE to establish multi-stage security.W  
\heading{UCE for authenticated encryption.} We will first define a notion of UCE-AE. While the encryption key $K$ seems like an obvious choice to play the role of the hash key in such a definition, we are skeptical there are applications  of the resulting notion and in preliminary work we  instead choose the \emph{encryption randomness} $r$ to pay the role of the hash key, for reasons explained later. 

We then want to build AE achieving our notion based on a conventionally-secure AE scheme and a UCE hash function, by adapting the construction of  key-dependent message secure symmetric-key encryption in the ROM by Black \emph{et al.}~\cite{C:BlaRogShr02}. Namely, we plan to replace the former building block in their construction with conventionally-secure AE and the random oracle with a UCE hash function.  In fact, Bellare \emph{et al.} previously used conventionally-secure symmetric-key encryption and a UCE hash function in the same construction, showing the resulting scheme is KDM-secure in the standard model. In fact, their proof uses the encryption randomness as a hash key, a main reason  chose to do in our definition of UCE-AE. 
To prove our construction meets UCE-AE in the standard model, we will generalize the proof technique   of~\cite{C:BelHoaKee13}. 
%by~\cite{EC:SonTes17} to show $\mathsf{BRS}[\mathsf{SE},H]$ is CP-indifferentiable.
%In particular, it will suffice by a result of~\cite{EC:SonTes17} that translating indifferentiability  to UCE-security for constructions meeting a relaxde notion of  indifferentiability in the ROM.

\iffalse
In an authenticated encryption (AE) scheme $\mathsf{AE} = (\calK,\calE,\calD)$, $\calE$ takes a key $K$, message $m$ and randomness $r$. (To simplify, we don't use a  nonce,  associated data, \emph{etc.} here). AE-security began with~\cite{AC:BelNam00} and evolved over time.
\begin{ques}
    Which input to AE encryption should play the role of the hash key in UCE?
\end{ques}
The secret key seems like the obvious choice, but we are skeptical there are applications  and choose $r$ instead. Then we conjecture that the definition of UCE-AE (which will extend UCE to \emph{injections})  generalizes KDM-security.  The proof should be similar to that of~\cite{C:BelHoaKee13} that $\mathsf{BRS}[\mathsf{SE},H]$~\cite{C:BlaRogShr02} is KDM-secure if $H$ is UCE and $\mathsf{SE}$ is IND-CPA.  Indeed, their proof uses  $r$ as a hash key. To achieve UCE-AE, we plan to generalize the result of~\cite{C:BelHoaKee13} on  $\mathsf{BRS}[\mathsf{SE},H]$~\cite{C:BlaRogShr02} from KDM-security to UCE-AE.

% With the above choice, ideal AE with key-space $\calK$, message-space $\calM$, and expansion-space $\calX$ provides  oracle access to  $\rho(\cdot,\cdot)$ and $\rho^{-1}(\cdot,\cdot)$, with $\tau \getsr \calX$ and $\rho$ drawn at random from the set of injections   $\mathsf{Inj}[\calK \times \calM \to \calC]$ such that (1) $|\rho(K,m)| = \tau$ for all $K\in \calK,m\in\calM$  and (2) $\rho(K,m) \neq \rho(K,m')$ for $m \neq m'$. Note that the ideal primitive differs from~\cite{C:BarFar18}. 
%Then, the UCE-AE game for $\mathsf{AE}$ against adversary $\advA = (\advs,\advd)$ chooses $\tau \getsr \calX$ and a random $b\getsr \bits$. Then $\advs$ gets oracle access to   $\calE(\cdot,\cdot,\tau)$  and its inverse if $b=1$, or to $\rho(\cdot, \cdot)$ and its inverse sampled as in the ideal primitives above if $b=0$, but \emph{not} $\tau$. It passes leakage $L$ to $\advd$, who gets $\tau$ but no oracles and has to guess $b$.
%Here the ideal primitive and definition of UCE-AE based on this choice. Intuitively, this extends UCE to keyed families of \emph{injections}. 

We will study restrictions needed on the source for UCE-AE to be achievable, as well what notions it implies.
%If the definition turns out not to be useful, we will choose $K$ instead of $\tau$ to play the role of the hash key in UCE.
To achieve it, we plan to generalize the result of~\cite{C:BelHoaKee13} on  $\mathsf{BRS}[\mathsf{SE},H]$~\cite{C:BlaRogShr02} from KDM-security to UCE-AE.
To this end, ST~\cite{EC:SonTes17} says that a CP-indifferentiable (a  substantial weakening of indifferentiability)  construction of ideal primitive $\mathbf{I}$ from  ideal primitive $\mathbf{J}$  is psPR   wrt.~reset sources if $\mathbf{J}$ is. So, casting UCE-AE in the psPR framework, we want to show:
\begin{task} \label{task-ro}
    Show  $\mathsf{BRS}[\mathsf{AE}, H]$ is CP-indifferentiable in the RO model if $\mathsf{AE}$ is AE-secure.
\end{task}
%Indeed, we conjecture that full indifferentiability does not hold for BRS, which is quite different from the schemes proposed  in~\cite{C:BarFar18}. 
However, we  need to use \emph{computational}  indifferentiability~\cite{C:ZhaZha20,AC:ZhaZha23} because of $\mathsf{SE}$. So,  we will see whether ST's result extends to that setting; if not, we will look for a direct proof of UCE-AE.
\fi
%Finally, we  will  investigate other    relaxations of multi-stage indifferentiability.
%thus showing that $\mathsf{BRS}[\mathsf{SE}, H]$ is UCE-AE wrt.~reset sources when $\mathsf{SE}$ is pseudorandom and  $H$ is UCE wrt.~reset sources. %We  will also investigate other   relaxations of multi-stage indifferentiability, \emph{e.g.}, where the adversary's queries in each stage are unpredictable.    %These may also be achievable in the standard model. Indeed, it is an interesting question where the ``boundary'' lies (in the form of restrictions to the definition) for indifferentibility  to be achievable in the standard model. 

\heading{UCE for public-key encryption.} We will also define a notion of UCE-PKE, again choosing the encryption randomness $r$ to play the role of the hash key in UCE.  Our approach to achieve it will be based on the widely-used ECIES~\cite{dhies} scheme.  In particular, we aim to show the \emph{first instantiability result about ECIES}.
\fi

\iffalse
Now we turn to:
%Now, to define a notion of UCE-PKE we ask:
\begin{ques}
    Which input to PKE encryption  should play the role of the hash key in UCE?
\end{ques}

Again, we conjecture that choosing the \emph{randomness} $r$  will be useful.
We will study what restrictions on the source are needed for UCE-PKE to be achievable, and what other notions it implies. %One interesting question is if it implies \emph{simultaneous} KDM+RKA-security, which would improve on the results of~\cite{TCC:BFHO22} for ElGamal (which show KDM-security and RKA-security under PGG separately). 
To achieve it, we propose a new approach  based on the  widely-used ECIES scheme~\cite{dhies}:
\begin{task} \label{task-2}
    Show that ECIES is CP-indifferentiable  in the GGM+RO+ideal AE models. 
\end{task}
%Of course, full indifferentiability would be even better, to strengthen the single-stage guarantees.
%Zhandry and Zhang~\cite{C:ZhaZha20} achieve indifferentiability using a less practical approach.
Actually, part of an ECIES ciphertext depends only on $pk$ and $r$, but this seems benign. So, we may need to modify  the definition of ideal PKE.
Then, we want to extend ST's result to constructions based on several primitives, giving us the 
\emph{first instantiability result about ECIES}. %We will need to investigate constructions of the building blocks. \emph{E.g.}, as a first cut,
Finally, we will study our UCE-style notions for the base primitives here wrt.~reset sources. % We will investigate other source classes as well. %In particular, considering PGG1 wrt.~reset sources  will be novel to the proposed work. 


\fi


% We will then extend the result to forms of indifferentiability such as context-restricted indifferentiability~\cite{SCN:MauJos18} while remaining in the standard  model. %Finally, we want to strengthen the above  result to achieve a multi-stage notion of indifferentiability where queries in each stage are \emph{unpredicatable}, staying the standard model. 
%\paragraph*{UCE-PKE.} \adam{Fill in}
% We  also need to study which multi-stage security notions UCE-AE implies.  Thus, the above task will  prove a stronger result than~\cite{C:BelHoaKee13}.
 
%\paragraph*{UCE-PKE.}
%Here, we study the ideal PKE model of Zhandry and Zhang~\cite{C:ZhaZha20}, which seems the most relevant in the context of UCE as it is public key.

 % However, a setting where we can drop this is \emph{hidden-order} groups. 
% \paragraph*{Adaptivity.} UCE and related notions are \emph{non-adaptive} in that the source's queries do not depend on the key.  Farshim and Mittelbach~\cite{FSE:FarMit16} addressed this for UCEs, proposing \emph{interactive computational extractors} (ICEs).  However, we find ICEs hard to work with,  and they lack the asymmetry between source and distinguisher. We ask whether a more straightforward notion suffices.  We briefly sketch a notion we will consider. Intuitively,  we $\advs$ can query any ``allowed'' function $f$ that on input $\hk$ outputs a pair $(\hk', x)$, the response to its query being  $H(f(\hk))$. Of course, the question is what ``allowed'' means. We will study this  and extend the treatment to PGG. % We anticipate these notions will be useful in handling adaptive signing queries in security proofs for signature schemes. 

%An initial possibility will be to consider PGG where the leakage $L$ output by $\advs$ consists only of group elements, which have to be explained in the sense of the AGM.  Then, we can also ask that $\advd$ provides an explanation of its decision in sense of~\cite{}, in terms of the elements in the leakage.
\iffalse
%To get some intuition, note that we can recast an assumption $\mathsf{X}$ in the AGM as as a standard-model knowledge assumption by assuming not that $\adv$  outputs a representation along with the solution but rather there exists a \emph{(non-blackbox) extractor}  $\mathcal{E}$ that outputs it.  Call such an assumption knowledge-$\mathsf{X}$.  For example, roughly speaking, knowledge-CDH says that if  $\advA$ given $(g,g^x,g^y)$ outputs $g^{xy}$, then an extractor getting $\advd$'s ``view'' (its coins and inputs) outputs $a,b,c \in \Z_p$ such that $ax+by+c = xy$.
%It follows by an identical argument as in the AGM~\cite{} that knowledge-CDH reduces to DL, so knowledge-CDH, or even knowledge-CDH + CDH, is not very helpful.
\fi
%Note that we can work in the standard model by asking that $\advs,\advd$ do not output explanations themselves but rather come equipped with suitable non-blackbox extractors. %It is instructive to see why such recasting does not work in general. For example,  consider knowledge-CDH, asking that if CDH-solver $\advA$ given $(g,g^x,g^y)$ outputs $g^{xy}$, then an extractor getting $\advd$'s ``view'' outputs $a,b,c \in \Z_p$ such that $ax+by+c = xy$.
%But it follows by an identical argument as in the AGM~\cite{} that such knowledge-CDH simply reduces to DL.
%PGG-style definitions  differ because  they have both easy and hard instances. 
 % Of course, we can further require $\advs$ and $\advd$ are algebraic here.  We will investigate these possibilities in the proposed work.

\iffalse
More specifically, in the AGM, if the adversary has $g_1,\ldots,g_n$ and outputs $h$, she must give  $a_1,\ldots,a_n \in \Z_p$ such that $\prod_{i=1}^n g_i^{a_i} = h$. The general framework  allows additional input and output as \emph{bits} with suitable restrictions, formalized by Zhandry in a circuit-based model~\cite{} in the case of non-bilinear groups and~\cite{} in an oracle-based model.
%This turns out to be to formalize, and the AGM has come under scrutiny in several recent works~\cite{}.  In the proposed work, we will build on the  formalization of the AGM due to  Zhandry~\cite{}. %We seek to model PGG in (an adaptation of) the AGM and show (1) that it is implied by a standard assumption in the AGM and (2) it is qualitatively weaker than the full AGM and in particular can be turned into a novel standard-model knowledge assumption.
Also note that while Zhandry shows that the AGM is  weak for multi-stage assumptions in general, there is no reason this applies to our specific assumption.
Now, we will carry out our agenda in several steps.  First:
\begin{task}
    Formulate  a \emph{search} (vs.~decisional version of BPGG. 
\end{task} 
For this, $\advs$ will be replaced with a \emph{DL-finder}. The motivation is that the standard definition of algebraic may then be applied to it.  We will also consider an algebraic \emph{source} $\advs$, where both its queries and output must be explained.  For this, we will 
 develop an oracle-aided extension of the AGM to accommodate the $\Exp$ oracle provided to the source in BPGG. %$Following~\cite{C:FucKilLos18} in the case of ROs, we will include this  oracle's responses to $\advs$ in  the basis with respect to which distinguisher $\advd$ must output an explanation; that is, the oracle's responses themselves do not need to be explained by the adversary.
 Finally, for decisional BPGG, we will use a decisional  version of the AGM (DAGM) due to Rotem and Segev~\cite{TCC:RotSeg20}  We dub the final assumption (also with multiple sources) \emph{knowledge} BPGG (kBPGG). 
% Let's say leakage~$L$ produced by $\advs$ contains all the oracle responses $h_1,\ldots,h_q$,  additional group elements $g_1,\dots,g_z$, and  some bitstring.  (In the AGM, group elements and bitstrings must be distinguished.)  Then, adapting Definition 3.1 of~\cite{TCC:RotSeg20}, for a PGG adversary $(\advs, \advd)$ with advantage $\varepsilon$ we require that there exists a bit $b'$ such that on a run of the PGG game with the challenge bit $b$ fixed to $b'$, $\advd$ outputs an explanation $\vecw$ such that, first of all,
%\begin{newmath}
%\prod_{j=1}^{z} g_j^{\vecw[q+j]} \circ g^{\vecw[q+j+1]} = 1_{\GG}
%\end{newmath}%
%with probability at least $\varepsilon/t^2$ over the run ($t$ is the adversary's runtime), and that the views of $\advd$ when ``projected'' onto $\vecw$ are distinct for $b\in\bits$.  
\fi
%An interesting question is why we only give the $\Exp$ oracle to $\advA$. One might also give a group operation oracle  $\textbf{Op}(x,y) := g^{\Dlog_{G,g}(x) + \Dlog_{G,g}(y)}$ for $x,y \in \Z_p$.  The reason is because the group operation oracle can be computed by the adversary on its own as $x+y$.
%This conition excludes the attack above, since $\advp$ can output $(0, 1, 1, -1)$ to win the game. One might try to modify the attack and let the source leak $(x_3, h_3)$ to~$\advd$, who, given $g$, can compute $g^{x_3}$ and compare it to $h_3$. But this also contradicts algebraic unpredictability, with a predictor returning $(x_3, 1, 1)$.
%Due to the existence of obfuscation-based attacks similar to those for UCEs~\cite{C:BrzFarMit14},  we allow the algebraic predictors to run in unbounded time. 
%\emph{fixed-generator} assumptions under PGG.  \emph{E.g.}, taking $p_i(x,r) = rx + x^i$ and $L=r$ we get Assumption 3 of~\cite{}. This is surprising because, as explained in~\cite{}, the PGG approach seemed tied to a random group generator.  It should also allow us to capture assumptions where exponents are  randomized by values unknown to the source, such as matrix DDH. 

\iffalse
\paragraph*{Parallel structure.}
It turns out that algebraic unpredictability by itself is not sufficient to rule out all generic attacks. Indeed, consider a source $\advs$ that samples $x \sample \ZZ_p$, queries $h_1 \gets \Exp(x)$ and $h_2 \gets \Exp(x^2)$, then computes \mbox{$h_2' \gets h_1^x$} and passes the bit $(h_2 = h_2')$ as leakage to $\advd$, who decides accordingly. Again, the advantage of $(\advs, \advd)$ in the PGG game is almost $1$, and now $\advs$ is even algebraically unpredictable. The issue here is that $\advs$'s queries satisfy a linear relation with coefficients that are themselves \emph{unpredictable but known to $\advs$} (in this case, $x \cdot x - 1 \cdot x^2 = 0$), an attack vector not ruled out by algebraic unpredictability. 

To address this problem, we consider \emph{parallel sources}. Loosely speaking, this means that $\advs$'s $\Exp$ queries are made in parallel by single-query sources $\advs_i(\state)$ which, other than receiving a common initial state, do not pass state among each other.
Indeed, the attack above was possible because $\advs$ could learn more than one oracle reply. %Note that, in this example, although the queries~$x$ and $x^2$ are allowed, the equality check $h_2 = h_2'$ requires knowledge of $h_2$ and~$h_2'$ (related to different queries), and hence violates the definition of a parallel source.


% Algebraic unpredictability is then required from the joint leakage of $\advs_i$.

%Note that, in the example above, the queries $x_1$ and $x_2$ are allowed, but the check $h_2 = h_2'$ requires knowing $h_2$ and $h_2'$ (related to different queries), and hence violates the definition of a parallel source.\par

\paragraph*{Restricted post-processing.}
Even considering only parallel sources does not rule out all trivial attacks. Indeed, one can modify the source~$\advs$ from above and make it parallel by setting $\state \gets x$, having $\advs_1(\state)$ compute $h_1$ and $h_2'$, and letting $\advs_2(\state)$ compute $h_2$. Leakage $(h_2', h_2)$ is passed to $\advd$, which returns the bit $(h_2 = h_2')$. This attack works because each $\advs_i(\state)$ still allows arbitrary post-processing of its oracle response (here, computing the exponentiation of $h_1$). Accordingly, we further restrict the class of sources and consider  parallel sources where each $\advs_i(\state)$ is allowed only \emph{structured post-processing} of its oracle replies.  
We show that in Shoup's GGM~\cite{EC:Shoup97} a generic group is a PGG for such sources, meaning PGG  captures a bonafide property of generic groups just as UCE captures a bonafide property or a RO.

\paragraph*{What is  gained?} 
PGG  ports UCE to GGM for algebraically unpredictable, parallel, and structured-postprocessing sources (which we call \emph{masking sources}). For simplicity, we often leave the source class implicit, just writing `PGG'. While masking sources seem very specific, they suffice for some applications~\cite{}.  Moreover, we will devise alternative restrictions on the source in the proposed work.
% PGGs are a new way we introduce of formalizing what it means for a group to be ``generic.'' 
%We stress that
 %PGG wrt.~any source class is a standard-model assumption, not an idealized model.  %We will also show  that PGG is weaker than the (decisional) algebraic group model (AGM)~\cite{}. 
Finally, PGG  should be  compared to the Uber assumption~\cite{}.  In preliminary work, we show Uber is implied by PGG.  Indeed, generalizing  the Uber assumption, PGG   handles (1) non-uniform exponents. (2) exponents that are not correlated by public polynomials, and (3)  leakage on the exponents. Looking ahead to the \emph{bilinear} group setting, we contend handling these cases will be crucial for transitioning proofs of schemes set in the GGM or the ostensibly weaker \emph{algebraic group model} (AGM) to the standard model.  %RaIn some sense, this gives evidence we can have the best of both worlds in terms of theory and practice.   
%As PGG  is built upon in the proposed work, we will see even more exciting applications. %It would thus be surprising if the assumption did not hold in appropriate elliptic curves.
\fi


\iffalse
   
\fi


\iffalse
\subsection{More Standard-Model Notions for Groups} \label{sec-port}

\begin{ques}
    What other useful hash function properties used for instantiating ROs besides UCE can be ported to instantiations of generic groups?
\end{ques}

\paragraph*{Correlation Intractability.}  Fix a hash function $\Hash \colon K \times D \to R$.  For $q\in\N$, $q$-input correlation-intractability~\cite{} asks for certain $2q$-ary relations $\calR$, given $k \sample K$ it is hard to find $x_1,\ldots,x_q \in D$ such that $(x_1,\ldots,x_q,\Hash_k(x_1),\ldots,\Hash_k(x_q) \in \calR$.  Positive results are known relative to some relations useful for NIZKs~\cite{}. 
Actually, a notion of correlation intractability for groups was already proposed by Dent~\cite{} to show uninstantiability of the GGM. 
However, version speaks to correlation intractability of the generic encoding function $\encode \colon \GG \to \Z_p$.   We consider a mapping in the other direction, namely exponentiation.
Taking inspiration from PGG1, we will  consider a notion asking that for certain  $2q$-ary relations $\calR$, given $(\params := (\circ, g, p)) \sample \group(\secparam)$ it is hard to find $x_1,\ldots,x_q \in \Z_p$ such that $(x_1,\ldots,x_q,g^{x_1},\ldots,g^{x_q}) \in \calR$. % Analogously to PGG, we may consider having the generator $g$ (but not the rest of $\params$) withheld from the adversary.

\paragraph*{Progammability.}
A programmable hash function is a keyed hash function to a group introduced by Hofheinz and Kiltz~\cite{}.   They ask that given $x_1,\ldots,x_m \in D$ and $z_1,\ldots,z_n \in D$ over the choice of $k \getsr K$, given associated trapdoor information depending on $g, h\in\GG$, for all $i\in [m], j\in[n]$, it is possible to recover $a_i,b_i \Z_p$ such that $\Hash_k(x_i) = g^{a_i} h^{b_i}$ and $b_j \in \Z_p$ such that $\Hash_k(z_j) = g^{b_j}$.   %To port this notion to groups, one possibility is simply consider the hash function to be exponentiation. Fixing $x_1,\ldots,x_m \in \Z_p$ and $z_1,\ldots,z_n \in \Z_p$, over the choices $g \sample \GG$ and some trapdoor information one can recover discrete logs of $X_i = g^{x_i}$ and each $g^{z_j}$

\paragraph*{Extreme Lossiness.}
Lossy functions were introduced by Peikert and Waters~\cite{} where they additionally assumed a trapdoor, which we do not. Lossiness for groups can be seen as subgroup-decision assumptions~\cite{}. Extreme lossiness, however, appears novel in this context.  Note that is is unclear that elliptic curve groups can achieve extreme lossiness.  Nevertheless, it may be possible to construct such groups from cryptographic assumptions. Such a result could show instantiations of generic groups \emph{in principle}.

\paragraph*{Plausibility and Achievability.}  We ask two main questions about the notions:
\begin{ques}
     (1) Are the notions achievable at all? (2) Are the notions achieved by appropriate elliptic curve groups?  
\end{ques}
\fi



\fi
\fi